% Traduction : phrase causale, finalité, futur dans la subordonnée, emphase (révision) — Espagnol Tle A — Cours signé OnBuch+
% Programme officiel MINESEC. Compiler avec : tectonic E21-causal-finalidad-futuro-subordinada.tex
\newcommand{\DOCMATIERE}{Espagnol}
\newcommand{\DOCNIVEAU}{Tle A}
\newcommand{\DOCMODULE}{Módulo 4 : Activités économiques et monde du travail}
\newcommand{\DOCLECON}{E21}
\newcommand{\DOCTITRE}{Traduction : phrase causale, finalité, futur dans la subordonnée, emphase (révision)}
% OnBuch+ — préambule « cours signés » ENRICHI (compilé avec tectonic / XeLaTeX)
% Système visuel « Pop » d'OnBuch+ : fond crème (jamais blanc), bordures encre
% épaisses, ombres « dures » décalées, titres Archivo Black en capitales.
% Version MATHÉMATIQUES : graphes (pgfplots/tikz), boîtes pédagogiques
% (définition, propriété, méthode, attention, le savais-tu, exercice résolu,
% exercice, TP, bilan), page de garde et signature OnBuch+ ancrée en bas.
%
% Macros attendues dans main.tex AVANT \input{preamble} :
%   \DOCMATIERE  \DOCNIVEAU  \DOCMODULE  \DOCLECON (numéro)  \DOCTITRE
\documentclass[11pt]{article}

\usepackage{fontspec}
\usepackage[spanish,french]{babel} % français = langue principale ; l'espagnol via \esp{..} / espagnol
\usepackage[a4paper,margin=2cm,top=3.4cm,bottom=2.8cm,headheight=1.4cm,footskip=1.3cm]{geometry}
\usepackage{amsmath,amssymb,amsfonts}
\usepackage{mathtools} % \xmapsto, \coloneqq... souvent utilisés par les modèles
\usepackage{cancel} % simplifications barrées (bilans de forces)
% \abs{z} (module, valeur absolue) et \norm{u} (norme), utilisés par les modèles
\providecommand{\abs}{}\let\abs\relax\DeclarePairedDelimiter{\abs}{\lvert}{\rvert}
\providecommand{\norm}{}\let\norm\relax\DeclarePairedDelimiter{\norm}{\lVert}{\rVert}
\usepackage[table]{xcolor} % [table] : fournit aussi \rowcolors (alternance de lignes)
\usepackage{pagecolor}
\usepackage{tikz}
\usetikzlibrary{arrows.meta,positioning,calc,shapes.geometric,shapes.misc,shapes.arrows,decorations.pathmorphing,decorations.pathreplacing,decorations.markings,patterns,shadows,mindmap,trees,fit,backgrounds,matrix,intersections,angles,quotes}
\usepackage{pgfplots}
\usepackage[version=4]{mhchem} % équations nucléaires \ce{^{235}_{92}U}
\usepackage[european,siunitx]{circuitikz} % schémas électriques
\pgfplotsset{compat=1.18}
\usepgfplotslibrary{fillbetween,groupplots} % aires entre courbes (calcul intégral)
\usepackage{enumitem}
\usepackage{fancyhdr}
% [most] charge toutes les bibliothèques tcolorbox (listings, minted, raster,
% documentation...) et gonfle inutilement la mémoire TeX sur un document long
% (~300 boîtes) : on ne charge que ce qui est réellement utilisé.
\usepackage[skins,breakable]{tcolorbox}
\usepackage{graphicx}
\usepackage{colortbl}
\usepackage{booktabs}
\usepackage{tabularx}
\usepackage{diagbox} % case d'en-tête diagonale des tableaux à double entrée
% Colonnes extensibles centrée / gauche / droite (C, L, R), souvent utilisées par les modèles
\newcolumntype{C}{>{\centering\arraybackslash}X}
\newcolumntype{L}{>{\raggedright\arraybackslash}X}
\newcolumntype{R}{>{\raggedleft\arraybackslash}X}
\usepackage{array}
\usepackage{multicol}
\usepackage{multirow}
\usepackage{siunitx}
% parse-numbers=false : accepte aussi \SI{2,5 \times 10^{-3}}{...}
\sisetup{locale=FR,output-decimal-marker={,},group-minimum-digits=4,parse-numbers=false}
\DeclareSIUnit\ohm{\ensuremath{\symup{\Omega}}} % U+2126 absent de la police
\DeclareSIUnit\division{div} % réglages d'oscilloscope (ms/div, V/div)
\DeclareSIUnit\tr{tr} % tours (vitesse de rotation, ex. \SI{1500}{\tr\per\minute})
\DeclareSIUnit\molar{M} % concentration molaire (ex. \SI{3}{\molar}, \SI{1}{\micro\molar})
\DeclareSIUnit\an{an} % année (le modèle confond parfois avec \year, un compteur TeX) — ex. \SI{5}{\kilo\meter\per\an}
\DeclareSIUnit\FCFA{FCFA} % franc CFA (ex. \SI{500}{\FCFA\per\kilo\gram})
\DeclareSIUnit\degreeBrix{\textdegree Brix} % teneur en sucre d'un sirop/jus (réfractométrie)
\DeclareSIUnit\month{mois} % le modèle utilise parfois \month, un compteur TeX, comme unité de durée
\usepackage{qrcode}
% \degree : commande fréquemment utilisée par les modèles pour un angle ou une
% température en dehors de \SI{}{} (non fournie par les paquets chargés).
\providecommand{\degree}{^{\circ}}
% \qed (fin de démonstration), utilisé par les modèles sans amsthm
\providecommand{\qed}{\hfill$\square$}
% \barre : double barre des valeurs interdites dans un tableau de variations (tkz-tab)
\providecommand{\barre}{\ensuremath{\|}}
% environnement proof (amsthm non chargé) utilisé par les modèles
\ifdefined\proof\else\newenvironment{proof}[1][Démonstration]{\par\noindent\textit{#1.}\ }{\qed\par}\fi
\usepackage{lastpage}
\usepackage{hyperref}
\hypersetup{hidelinks}

% ============================================================
% Palette « Pop » OnBuch+ — mêmes hex que la classe P (plus_ui.dart)
% ============================================================
\definecolor{poporange}{HTML}{F59321}
\definecolor{popdark}{HTML}{E07A0C}
\definecolor{popcream}{HTML}{FFF6E8}
\definecolor{popink}{HTML}{1C1714}
\definecolor{poppurple}{HTML}{7A5AE0}
\definecolor{popgreen}{HTML}{2E9E6A}
\definecolor{poppink}{HTML}{E0457B}
\definecolor{popblue}{HTML}{2F7DE1}
\definecolor{popgold}{HTML}{C98A12}
\definecolor{popmuted}{HTML}{8A7F76}
\definecolor{popline}{HTML}{EADFD2}
\definecolor{popgray}{HTML}{8A7F76}
\colorlet{popgrey}{popgray}
% Alias tolérants (le modèle nomme parfois les couleurs différemment)
\colorlet{popwhite}{white}
\colorlet{popblack}{popink}
\colorlet{popred}{poppink}
\colorlet{popyellow}{popgold}
% Teintes claires (fonds de tableaux/encadrés) — le modèle nomme parfois
% les couleurs avec un suffixe L (ex. popblueL) sans les avoir définies.
\colorlet{poporangeL}{poporange!20}
\colorlet{popdarkL}{popdark!20}
\colorlet{poppurpleL}{poppurple!20}
\colorlet{popgreenL}{popgreen!20}
\colorlet{poppinkL}{poppink!20}
\colorlet{popblueL}{popblue!20}
\colorlet{popgoldL}{popgold!20}
\colorlet{popredL}{popred!20}
\colorlet{popgrayL}{popgray!20}
% Teintes claires pour fonds de tableaux / zones
\colorlet{popblueL}{popblue!10!white}
\colorlet{poporangeL}{poporange!14!white}
\colorlet{popgreenL}{popgreen!12!white}
\colorlet{poppurpleL}{poppurple!12!white}
\colorlet{poppinkL}{poppink!12!white}
\colorlet{popgoldL}{popgold!14!white}

\pagecolor{popcream}

% ============================================================
% Typographie
% ============================================================
\defaultfontfeatures{Ligatures=TeX}
\setmainfont{PlusJakartaSans-Medium.ttf}[Path=../../../fonts/,BoldFont=PlusJakartaSans-Bold.ttf]
% Symboles, API et emojis absents de la police principale : repli sur DejaVu Sans (caractères actifs)
\newfontface\symfont{DejaVuSans.ttf}[Path=../../../fonts/]
\makeatletter
\newcommand\symactive[1]{\catcode#1=13 \begingroup\lccode`\~=#1 \lowercase{\endgroup\def~}{{\symfont\char#1}}}
\newcommand\symblank[1]{\catcode#1=13 \begingroup\lccode`\~=#1 \lowercase{\endgroup\def~}{}}
\newcommand\symhyph[1]{\catcode#1=13 \begingroup\lccode`\~=#1 \lowercase{\endgroup\def~}{\mbox{-}}}
\newcount\symc
\newcommand\symrange[2]{\symc=#1 \@whilenum\symc<#2 \do{\expandafter\symactive\expandafter{\the\symc}\advance\symc by 1 }}
\newcommand\symblankrange[2]{\symc=#1 \@whilenum\symc<#2 \do{\expandafter\symblank\expandafter{\the\symc}\advance\symc by 1 }}
\makeatother
\symrange{"0250}{"02FF}\symactive{"2713}\symactive{"2714}\symactive{"2717}\symactive{"2718}\symactive{"2610}\symactive{"2611}\symactive{"2612}
\symhyph{"2011}\symhyph{"2E1A}\symblankrange{"1F300}{"1FAFF}\symblank{"FE0F}
% Maths en sans-serif (Fira Math) avec chiffres et lettres droites de Plus
% Jakarta, pour que les formules se fondent dans le texte.
\usepackage[math-style=ISO,bold-style=ISO]{unicode-math}
\setmathfont{FiraMath-Regular.otf}[Path=../../../fonts/]
\setmathfont{PlusJakartaSans-Medium.ttf}[Path=../../../fonts/,range={up/num,up/{Latin,latin},\mathup}]
\usepackage{icomma} % virgule décimale sans espace en mode math
% Caractères mathématiques Unicode absents des polices de texte (Plus Jakarta,
% Archivo Black) : rendus par leur équivalent mathématique, en texte comme en
% formule — sinon ils s'affichent en carré vide (ex. « Arithmétique dans ℤ »).
\usepackage{newunicodechar}
\newunicodechar{ℝ}{\ensuremath{\mathbb{R}}}
\newunicodechar{ℤ}{\ensuremath{\mathbb{Z}}}
\newunicodechar{ℂ}{\ensuremath{\mathbb{C}}}
\newunicodechar{ℕ}{\ensuremath{\mathbb{N}}}
\newunicodechar{ℚ}{\ensuremath{\mathbb{Q}}}
\newunicodechar{𝔻}{\ensuremath{\mathbb{D}}}
\newunicodechar{α}{\ensuremath{\alpha}}
\newunicodechar{β}{\ensuremath{\beta}}
\newunicodechar{γ}{\ensuremath{\gamma}}
\newunicodechar{δ}{\ensuremath{\delta}}
\newunicodechar{ε}{\ensuremath{\varepsilon}}
\newunicodechar{θ}{\ensuremath{\theta}}
\newunicodechar{λ}{\ensuremath{\lambda}}
\newunicodechar{μ}{\ensuremath{\mu}}
\newunicodechar{σ}{\ensuremath{\sigma}}
\newunicodechar{τ}{\ensuremath{\tau}}
\newunicodechar{φ}{\ensuremath{\varphi}}
\newunicodechar{ψ}{\ensuremath{\psi}}
\newunicodechar{ω}{\ensuremath{\omega}}
\newunicodechar{Σ}{\ensuremath{\Sigma}}
% majuscules produites par \MakeUppercase dans les titres (α-aminés -> Α)
\newunicodechar{Α}{\ensuremath{\alpha}}
\newunicodechar{Β}{\ensuremath{\beta}}
\newunicodechar{Γ}{\ensuremath{\Gamma}}
\newunicodechar{Δ}{\ensuremath{\Delta}}
\newunicodechar{Ε}{\ensuremath{\varepsilon}}
\newunicodechar{Θ}{\ensuremath{\Theta}}
\newunicodechar{Λ}{\ensuremath{\Lambda}}
\newunicodechar{Μ}{\ensuremath{\mu}}
\newunicodechar{Τ}{\ensuremath{\tau}}
\newunicodechar{Φ}{\ensuremath{\Phi}}
\newunicodechar{Ψ}{\ensuremath{\Psi}}
\newunicodechar{Ω}{\ensuremath{\Omega}}
\newunicodechar{↦}{\ensuremath{\mapsto}}
\newunicodechar{∈}{\ensuremath{\in}}
\newunicodechar{∉}{\ensuremath{\notin}}
\newunicodechar{∧}{\ensuremath{\wedge}}
\newunicodechar{⇔}{\ensuremath{\Leftrightarrow}}
\newunicodechar{⟺}{\ensuremath{\Longleftrightarrow}}
\newunicodechar{⇒}{\ensuremath{\Rightarrow}}
\newunicodechar{⟹}{\ensuremath{\Longrightarrow}}
\newunicodechar{∘}{\ensuremath{\circ}}
\newunicodechar{∪}{\ensuremath{\cup}}
\newunicodechar{∩}{\ensuremath{\cap}}
\newunicodechar{≡}{\ensuremath{\equiv}}
\newunicodechar{⊂}{\ensuremath{\subset}}
\newunicodechar{⊥}{\ensuremath{\perp}}
\newunicodechar{∃}{\ensuremath{\exists}}
\newunicodechar{∀}{\ensuremath{\forall}}
\newunicodechar{∛}{\ensuremath{\sqrt[3]{\,}}}
\newunicodechar{∜}{\ensuremath{\sqrt[4]{\,}}}
\newunicodechar{⃗}{\ensuremath{{}^{\to}}}
\newunicodechar{ⁿ}{\ifmmode^{n}\else\textsuperscript{\ensuremath{n}}\fi}
\newunicodechar{ˣ}{\ifmmode^{x}\else\textsuperscript{\ensuremath{x}}\fi}
\newunicodechar{ᵇ}{\ifmmode^{b}\else\textsuperscript{\ensuremath{b}}\fi}
\newunicodechar{ʳ}{\ifmmode^{r}\else\textsuperscript{\ensuremath{r}}\fi}
\newunicodechar{ᵘ}{\ifmmode^{u}\else\textsuperscript{\ensuremath{u}}\fi}
\newunicodechar{ʸ}{\ifmmode^{y}\else\textsuperscript{\ensuremath{y}}\fi}
\newunicodechar{ᵃ}{\ifmmode^{a}\else\textsuperscript{\ensuremath{a}}\fi}
\newunicodechar{ᵉ}{\ifmmode^{e}\else\textsuperscript{\ensuremath{e}}\fi}
\newunicodechar{ᵏ}{\ifmmode^{k}\else\textsuperscript{\ensuremath{k}}\fi}
\newunicodechar{ᵖ}{\ifmmode^{p}\else\textsuperscript{\ensuremath{p}}\fi}
\newunicodechar{ᴺ}{\ifmmode^{N}\else\textsuperscript{\ensuremath{N}}\fi}
\newunicodechar{ᶜ}{\ifmmode^{c}\else\textsuperscript{\ensuremath{c}}\fi}
\newunicodechar{ʰ}{\ifmmode^{h}\else\textsuperscript{\ensuremath{h}}\fi}
\newunicodechar{ᵠ}{\ifmmode^{q}\else\textsuperscript{\ensuremath{q}}\fi}
\newunicodechar{ₙ}{\ifmmode_{n}\else\textsubscript{\ensuremath{n}}\fi}
\newunicodechar{ₐ}{\ifmmode_{a}\else\textsubscript{\ensuremath{a}}\fi}
\newunicodechar{ᵢ}{\ifmmode_{i}\else\textsubscript{\ensuremath{i}}\fi}
\newunicodechar{ᵣ}{\ifmmode_{r}\else\textsubscript{\ensuremath{r}}\fi}
\newunicodechar{ₖ}{\ifmmode_{k}\else\textsubscript{\ensuremath{k}}\fi}
\newunicodechar{ᵦ}{\ifmmode_{\beta}\else\textsubscript{\ensuremath{\beta}}\fi}
\newunicodechar{ₓ}{\ifmmode_{x}\else\textsubscript{\ensuremath{x}}\fi}
\newfontfamily\popdisplay{ArchivoBlack-Regular.ttf}[Path=../../../fonts/]
\newfontfamily\poplogo{SpaceGrotesk-Bold.ttf}[Path=../../../fonts/]
\newfontfamily\popsemi{PlusJakartaSans-SemiBold.ttf}[Path=../../../fonts/]
\newfontfamily\popxbold{PlusJakartaSans-ExtraBold.ttf}[Path=../../../fonts/]
\newfontfamily\popxboldls{PlusJakartaSans-ExtraBold.ttf}[Path=../../../fonts/,LetterSpace=3.0]

\setlist[enumerate]{leftmargin=1.7em,itemsep=5pt,topsep=4pt}
\setlist[itemize]{leftmargin=1.7em,itemsep=4pt,topsep=4pt,label={\color{poporange}\textbullet}}
\setlength{\parindent}{0pt}
\setlength{\parskip}{5pt}
\renewcommand{\arraystretch}{1.25}

% pgfplots : style Pop par défaut
\pgfplotsset{
  every axis/.append style={axis line style={popink,line width=1pt},tick style={popink},
    grid style={popline,line width=0.5pt},label style={font=\small},tick label style={font=\footnotesize}},
  popcurve/.style={poporange,line width=1.8pt,no markers,smooth},
}
% Même style disponible dans un \draw TikZ simple (hors pgfplots)
\tikzset{popcurve/.style={poporange,line width=1.8pt}}
\tikzset{popgrid/.style={popline,very thin}}
\colorlet{popcurve}{poporange} % le nom du style est parfois utilisé comme couleur

% ============================================================
% Pill / squiggle
% ============================================================
\newcommand{\poppill}[3]{%
  \tikz[baseline=-0.55ex]\node[draw=popink,line width=1.2pt,rounded corners=2.6pt,%
    fill=#2,inner xsep=6.5pt,inner ysep=3pt]{%
    \popxboldls\fontsize{7}{7}\selectfont\color{#3}\MakeUppercase{#1}};%
}
\newcommand{\popsquiggle}[1]{%
  \tikz[baseline=-0.4ex]{
    \draw[#1,line width=2.6pt,line cap=round]
      plot [smooth] coordinates {(0,0.09) (0.34,0.01) (0.68,0.21) (1.02,0.01) (1.36,0.10)};
  }%
}
% Mot-clé surligné (marqueur orange sous le texte)
\newcommand{\cle}[1]{{\popxbold\color{popdark}#1}}

% ============================================================
% En-tête / pied
% ============================================================
\pagestyle{fancy}
\fancyhf{}
\renewcommand{\headrulewidth}{0pt}
\renewcommand{\footrulewidth}{0pt}
\lhead{\raisebox{-0.15ex}{\poplogo\fontsize{13}{13}\selectfont\color{popink}OnBuch\color{poporange}+}}
\rhead{\poppill{\DOCMATIERE\ \textbullet\ \DOCNIVEAU\ \textbullet\ Leçon \DOCLECON}{white}{popink}}
\lfoot{\popxbold\fontsize{7.6}{7.6}\selectfont\color{popmuted}\MakeUppercase{Cours signé OnBuch+}\ \textbullet\ {\popsemi\DOCTITRE}}
\rfoot{\tikz[baseline=-0.6ex]\node[rounded corners=8.5pt,fill=popink,minimum height=17pt,minimum width=17pt,inner xsep=5pt,inner sep=0]{\popxbold\fontsize{8}{8}\selectfont\color{popcream}\thepage\,/\,\pageref*{LastPage}};}

% ============================================================
% Titres
% ============================================================
\newcommand{\chaptitre}[2]{%
  \begin{tcolorbox}[enhanced,colback=poporange,colframe=popink,boxrule=2.6pt,arc=11pt,
      drop shadow=popink,left=15pt,right=15pt,top=12pt,bottom=11pt,before skip=3pt,after skip=18pt]
    \poppill{#2}{popink}{popcream}\par\vspace{9pt}
    {\raggedright\hyphenpenalty=10000\exhyphenpenalty=10000\popdisplay\fontsize{22}{24}\selectfont\color{popcream}\MakeUppercase{#1}\par}
  \end{tcolorbox}
}
% Section numérotée : pastille orange avec le numéro + titre Archivo
\newcounter{popsec}
\newcommand{\coursec}[1]{%
  \stepcounter{popsec}\setcounter{popsub}{0}%
  \par\vspace{16pt}\noindent
  \begin{minipage}{\linewidth}
  \tikz[baseline=-0.7ex]\node[draw=popink,line width=1.6pt,fill=poporange,rounded corners=4pt,minimum size=22pt,inner sep=0]{\popdisplay\fontsize{12}{12}\selectfont\color{popcream}\thepopsec};%
  \hspace{8pt}{\popdisplay\fontsize{15}{17}\selectfont\color{popink}\MakeUppercase{#1}}\par
  \vspace{4pt}\noindent{\color{poporange}\rule{36pt}{3.6pt}}
  \end{minipage}\par\nopagebreak\vspace{8pt}
}
\newcounter{popsub}
\newcommand{\courssub}[1]{%
  \stepcounter{popsub}%
  \par\vspace{9pt}\noindent{\popxbold\fontsize{11.5}{13}\selectfont\color{popink}{\color{poporange}\thepopsec.\thepopsub}\hspace{6pt}#1}\par\nopagebreak\vspace{4pt}
}

% ============================================================
% Boîtes pédagogiques — même recette : fond blanc, bordure épaisse colorée,
% ombre dure encre, badge plein. Titre optionnel [#1].
% ============================================================
\tcbset{popbase/.style={breakable,enhanced,colback=white,boxrule=2.3pt,arc=9pt,
  shadow={3pt}{-3pt}{0pt}{popink},left=14pt,right=14pt,top=11pt,bottom=11pt,before skip=11pt,after skip=15pt,
  fontupper=\popsemi\fontsize{10.6}{14.5}\selectfont\color{popink},
  fontlower=\popsemi\fontsize{10.6}{14.5}\selectfont\color{popink}}}
% Coche dessinée (le glyphe ✓ n'existe pas dans les polices du document)
\newcommand{\popcheck}[1]{\tikz[baseline=-0.5ex]\draw[#1,line width=1.6pt,line cap=round,line join=round](-0.13,0.0)--(-0.04,-0.09)--(0.13,0.1);}
\providecommand{\coche}{\popcheck{popgreen}}
\providecommand{\resultat}[1]{\par\smallskip\noindent\fcolorbox{popgreen}{popgreenL}{\parbox{\dimexpr\linewidth-2\fboxsep-2\fboxrule\relax}{\textbf{Résultat.} #1}}\par\smallskip}
\newcommand{\popboxhead}[3]{\poppill{#1}{#2}{white}\ifx\relax#3\relax\else\hspace{6pt}{\popxbold\fontsize{10.6}{12}\selectfont\color{popink}#3}\fi\par\vspace{7pt}}

\newtcolorbox{aretenir}[1][]{popbase,colframe=popgreen,before upper={\popboxhead{À retenir}{popgreen}{#1}}}
% Texte en espagnol (document de lecture, dialogue, poème, extrait) : cadre
% d'étude. Un titre optionnel (source, auteur) s'affiche dans l'en-tête.
\newtcolorbox{textebox}[1][]{popbase,colframe=popblue,before upper={\popboxhead{Texto}{popblue}{#1}\begin{otherlanguage}{spanish}},after upper={\end{otherlanguage}}}
% Marques ✓ / ✗ (forme correcte / fautive) souvent utilisées par les modèles
\providecommand{\cross}{\textcolor{poppink}{\ensuremath{\times}}}
\providecommand{\xmark}{\cross}\providecommand{\crossmark}{\cross}\providecommand{\cmark}{\ensuremath{\checkmark}}
% Ligne de réponse / justification à compléter (inventée par les modèles)
\providecommand{\justification}[1]{\par\noindent\textit{Justification~:} \dotfill\par}
% \textipa (package tipa non chargé) : la police ne couvre pas l'API -> simple passage
\providecommand{\textipa}[1]{#1}
% Soulignement (ulem non chargé) : \uline, \ul
\providecommand{\uline}[1]{\underline{#1}}\providecommand{\ul}[1]{\underline{#1}}
% \glossaire{\item ...} inventée par les modèles : liste de glossaire
\providecommand{\glossaire}[1]{\begin{itemize}#1\end{itemize}}
% \alert (beamer) inventée par les modèles : texte mis en évidence
\providecommand{\alert}[1]{\textbf{\textcolor{poppink}{#1}}}
% variantes ASCII d'ordinaux écrites en macro
\providecommand{\iere}{\textsuperscript{re}}\providecommand{\ieme}{\textsuperscript{e}}\providecommand{\ere}{\textsuperscript{re}}\providecommand{\eme}{\textsuperscript{e}}\providecommand{\er}{\textsuperscript{er}}\providecommand{\ie}{\textsuperscript{e}}
% Ordinaux français écrits en macro par les modèles : 1\ière, 3\ième
\newcommand{\ière}{\textsuperscript{re}}\newcommand{\ième}{\textsuperscript{e}}
% \exemple (inventée par les modèles) : étiquette « Exemple : »
\providecommand{\exemple}{\textbf{Exemple~:}\ }
% Mot ou phrase en espagnol à l'intérieur d'un paragraphe français
\newcommand{\esp}[1]{\ifmmode\text{\textit{#1}}\else\begingroup\selectlanguage{spanish}\itshape #1\endgroup\fi}
\newtcolorbox{exemplebox}[1][]{popbase,colframe=poporange,before upper={\popboxhead{Exemple}{poporange}{#1}}}
\newtcolorbox{definition}[1][]{popbase,colframe=popblue,before upper={\popboxhead{Définition}{popblue}{#1}}}
\newtcolorbox{propriete}[1][]{popbase,colframe=poppurple,before upper={\popboxhead{Propriété}{poppurple}{#1}}}
\newtcolorbox{methode}[1][]{popbase,colframe=popgold,before upper={\popboxhead{Méthode}{popgold}{#1}}}
\newtcolorbox{attention}[1][]{popbase,colframe=poppink,before upper={\popboxhead{Attention}{poppink}{#1}}}
\newtcolorbox{experience}[1][]{popbase,colframe=popblue,colback=popblueL,before upper={\popboxhead{Expérience}{popblue}{#1}}}
% « Le savais-tu ? » : carte violette pleine (contexte, culture, Cameroun)
\newtcolorbox{savaistu}[1][]{popbase,colback=poppurple,colframe=popink,
  fontupper=\popsemi\fontsize{10.4}{14.2}\selectfont\color{white},
  before upper={\poppill{Le savais-tu\,?}{white}{poppurple}\ifx\relax#1\relax\else\hspace{6pt}{\popxbold\color{white}#1}\fi\par\vspace{7pt}}}
% Exercice résolu : énoncé en haut, \tcblower puis la solution sur fond crème
\newcounter{popexo}\newcounter{popexr}
\newtcolorbox{exoresolu}[1][]{popbase,colframe=poporange,colbacklower=popcream,
  segmentation style={popink,line width=1.2pt,dashed},
  before upper={\stepcounter{popexr}\popboxhead{Exercice résolu \thepopexr}{poporange}{#1}},
  before lower={\poppill{Solution}{popink}{popcream}\par\vspace{6pt}}}
% Exercice (non résolu en place) — la correction est dans la partie Corrigés
\newtcolorbox{exercice}[1][]{popbase,colframe=popink,
  before upper={\stepcounter{popexo}\popboxhead{Exercice \thepopexo}{popink}{#1}}}
% Corrigé
\newtcolorbox{corrige}[1][]{popbase,colframe=popgreen,colback=popgreenL,
  before upper={\popboxhead{Corrigé}{popgreen}{#1}}}
% Objectifs / prérequis / bilan
\newtcolorbox{objectifs}{popbase,colframe=popink,colback=poporangeL,
  before upper={\popboxhead{Objectifs de la leçon}{popink}{}}}
\newtcolorbox{prerequis}{popbase,colframe=popink,colback=popblueL,
  before upper={\popboxhead{Prérequis}{popblue}{}}}
\newtcolorbox{bilan}{popbase,colframe=popink,colback=white,boxrule=2.8pt,
  before upper={\popboxhead{Fiche bilan}{poporange}{L'essentiel à connaître pour l'examen}}}
% Figure encadrée avec légende
\newtcolorbox{popfigure}[1][]{enhanced,breakable=false,colback=white,colframe=popline,boxrule=1.4pt,arc=8pt,
  left=10pt,right=10pt,top=10pt,bottom=8pt,before skip=10pt,after skip=12pt,halign=center,
  lower separated=false,halign lower=center,
  fontlower=\popsemi\fontsize{9}{11.5}\selectfont\color{popmuted},
  #1}
\newcommand{\legende}[1]{\par\vspace{6pt}{\popsemi\fontsize{9}{11.5}\selectfont\color{popmuted}#1}}

% ============================================================
% Signature OnBuch+
% ============================================================
\newcommand{\poplogoinline}[1]{{\poplogo\fontsize{#1}{#1}\selectfont\color{popink}OnBuch\color{poporange}+}}
\newcommand{\signatureonbuch}{%
  \begin{tcolorbox}[enhanced,colback=popink,colframe=popink,boxrule=2.2pt,arc=10pt,
      drop shadow=poporange,left=14pt,right=14pt,top=10pt,bottom=10pt,before skip=0pt,after skip=0pt]
    \tikz[baseline=-0.7ex]\node[circle,fill=popgreen,draw=popcream,line width=1.3pt,minimum size=22pt,inner sep=0]{\popcheck{white}};%
    \hspace{9pt}\begin{minipage}[c]{0.62\linewidth}
      {\popxbold\fontsize{12}{13}\selectfont\color{popcream}Signé\ \textbullet\ {\poplogo OnBuch\color{poporange}+}}\par\vspace{2pt}
      {\raggedright\popsemi\fontsize{8}{10.5}\selectfont\color{popline}\MakeUppercase{Équipe pédagogique OnBuch+}\\\MakeUppercase{Conforme au programme officiel MINESEC}\par}
    \end{minipage}\hfill
    {\poplogo\fontsize{22}{22}\selectfont\color{popcream}OnBuch\color{poporange}+}
  \end{tcolorbox}%
}

% ============================================================
% Page de garde
% #1 titre · #2 matière · #3 niveau · #4 module · #5 n° leçon · #6 durée ·
% #7 description / objectifs (texte ou itemize)
% ============================================================
\newcommand{\pagegarde}[7]{%
  \thispagestyle{empty}
  \newgeometry{a4paper,margin=1.7cm,top=1.6cm,bottom=1.4cm}%
  \begin{tikzpicture}[remember picture,overlay]
    \fill[poporange!18] ([xshift=-3.2cm,yshift=-3.6cm]current page.north east) circle (4.6cm);
    \fill[poppurple!14] ([xshift=2.4cm,yshift=7.6cm]current page.south west) circle (3.8cm);
  \end{tikzpicture}%
  {\poplogo\fontsize{30}{30}\selectfont\color{popink}OnBuch\color{poporange}+}\hfill
  \raisebox{0.6ex}{\poppill{Cours signé \textbullet\ Premium}{popink}{popcream}}\par
  \vspace{1pt}\popsquiggle{poppurple}\par
  \vspace{8pt}
  \begin{tcolorbox}[enhanced,colback=poporange,colframe=popink,boxrule=2.8pt,arc=13pt,
      drop shadow=popink,left=18pt,right=18pt,top=12pt,bottom=13pt,after skip=10pt]
    \poppill{#2 \textbullet\ #3}{popink}{popcream}\hspace{5pt}\poppill{#4}{white}{popink}\par\vspace{8pt}
    {\popdisplay\fontsize{14}{14}\selectfont\color{popink}LEÇON #5}\par\vspace{4pt}
    {\raggedright\hyphenpenalty=10000\exhyphenpenalty=10000\popdisplay\fontsize{25}{27}\selectfont\color{popcream}\MakeUppercase{#1}\par}
    \vspace{8pt}
    {\popxbold\fontsize{9.5}{9.5}\selectfont\color{popink}\MakeUppercase{Durée conseillée : #6}}
  \end{tcolorbox}
  \begin{tcolorbox}[enhanced,colback=white,colframe=popink,boxrule=2.2pt,arc=10pt,
      drop shadow=popink,left=16pt,right=16pt,top=10pt,bottom=10pt,after skip=8pt]
    \poppill{Dans cette leçon}{poppurple}{white}\par\vspace{5pt}
    \popsemi\fontsize{9.3}{12.6}\selectfont\color{popink}#7
  \end{tcolorbox}
  \noindent\begin{minipage}[t]{0.487\linewidth}
    \begin{tcolorbox}[enhanced,colback=popgreen,colframe=popink,boxrule=2.2pt,arc=10pt,
        drop shadow=popink,left=11pt,right=11pt,top=10pt,bottom=10pt,height=3.1cm,valign=center]
      \tcbox[colback=white,colframe=popink,boxrule=1.3pt,arc=5pt,boxsep=0pt,nobeforeafter,
        left=3pt,right=3pt,top=3pt,bottom=3pt,valign=center]{\qrcode[height=1.9cm]{https://play.google.com/store/apps/details?id=com.luvvix.onbuch&hl=fr}}%
      \hspace{8pt}\begin{minipage}[c]{0.5\linewidth}
        {\popdisplay\fontsize{11}{12.5}\selectfont\color{white}\MakeUppercase{Télécharge OnBuch}}\par\vspace{4pt}
        {\raggedright\popsemi\fontsize{8.4}{11}\selectfont\color{white}Gratuit \textbullet\ Google Play\\Tuteur IA, annales, concours, résultats\par}
      \end{minipage}
    \end{tcolorbox}
  \end{minipage}\hfill
  \begin{minipage}[t]{0.487\linewidth}
    \begin{tcolorbox}[enhanced,colback=poppurple,colframe=popink,boxrule=2.2pt,arc=10pt,
        drop shadow=popink,left=11pt,right=11pt,top=10pt,bottom=10pt,height=3.1cm,valign=center]
      \tikz[baseline=-0.7ex]\node[circle,fill=white,draw=popink,line width=1.3pt,minimum size=1.6cm,inner sep=0]{\popdisplay\fontsize{24}{24}\selectfont\color{poppurple}+};%
      \hspace{8pt}\begin{minipage}[c]{0.6\linewidth}
        {\popdisplay\fontsize{11}{12.5}\selectfont\color{white}\MakeUppercase{Passe à OnBuch+}}\par\vspace{4pt}
        {\raggedright\popsemi\fontsize{8.4}{11}\selectfont\color{white}Tous les cours signés, Tuteur IA illimité, fiches PDF — onglet OnBuch+ de l'app\par}
      \end{minipage}
    \end{tcolorbox}
  \end{minipage}
  \vspace{8pt}
  \popcontenu
  \vfill
  \signatureonbuch
  \restoregeometry
  \newpage
}

% Rangée « Ce cours contient » (page de garde)
\newcommand{\poptile}[3]{% #1 couleur #2 gros texte #3 légende
  \begin{tcolorbox}[enhanced,colback=#1,colframe=popink,boxrule=2pt,arc=9pt,shadow={2.5pt}{-2.5pt}{0pt}{popink},
      left=6pt,right=6pt,top=9pt,bottom=9pt,halign=center,valign=center,height=1.75cm,width=\linewidth,nobeforeafter]
    {\popdisplay\fontsize{12.5}{13}\selectfont\color{white}\MakeUppercase{#2}}\par\vspace{3pt}
    {\centering\popsemi\fontsize{7.8}{9.5}\selectfont\color{white}#3\par}
  \end{tcolorbox}}
\newcommand{\popcontenu}{%
  \par\noindent{\popxboldls\fontsize{7.5}{7.5}\selectfont\color{popmuted}CE COURS SIGNÉ CONTIENT}\par\vspace{7pt}
  \noindent\begin{minipage}[t]{0.235\linewidth}\poptile{popblue}{Cours}{complet, illustré, pas à pas}\end{minipage}\hfill
  \begin{minipage}[t]{0.235\linewidth}\poptile{poporange}{Méthodes}{et exercices résolus}\end{minipage}\hfill
  \begin{minipage}[t]{0.235\linewidth}\poptile{poppink}{Exercices}{type examen + corrigés}\end{minipage}\hfill
  \begin{minipage}[t]{0.235\linewidth}\poptile{popgreen}{Fiche bilan}{l'essentiel à retenir}\end{minipage}\par}

% Sommaire
\newtcolorbox{sommaire}{enhanced,colback=white,colframe=popink,boxrule=2.2pt,arc=10pt,
  drop shadow=popink,left=18pt,right=18pt,top=13pt,bottom=13pt,before skip=6pt,after skip=20pt,
  before upper={\poppill{Sommaire}{poppurple}{white}\par\vspace{9pt}\popsemi\fontsize{10.6}{14.5}\selectfont\color{popink}}}

% Signature de fin de cours — poussée tout en bas de la dernière page
\newcommand{\findecours}{%
  \par\vfill\nopagebreak
  \begin{center}\popsquiggle{poporange}\end{center}\vspace{4pt}
  \signatureonbuch
}
\AtBeginDocument{\def\llbracket{\mathopen{[\mkern-3.5mu[}}\def\rrbracket{\mathclose{]\mkern-3.5mu]}}} % intervalles d'entiers (glyphe ⟦ absent de la police)
% Glyphes absents de Fira Math : redéfinis à partir de glyphes disponibles
\AtBeginDocument{%
  \def\setminus{\mathbin{\backslash}}\let\smallsetminus\setminus
  \def\vdots{\mathord{\vbox{\baselineskip3.2pt\lineskiplimit0pt\kern4pt\hbox{.}\hbox{.}\hbox{.}}}}%
  \def\ddots{\mathinner{\mkern1mu\raise6.5pt\hbox{.}\mkern2mu\raise3.6pt\hbox{.}\mkern2mu\raise0.7pt\hbox{.}\mkern1mu}}%
  \def\triangle{\mathord{\scalebox{0.9}{$\symup{\Delta}$}}}%
  \def\nabla{\mathord{\raisebox{\depth}{\scalebox{1}[-1]{$\symup{\Delta}$}}}}%
  \def\checkmark{\popcheck{popdark}}%
  \def\star{\mathbin{\ast}}\def\bigstar{\mathord{\ast}}%
  \def\diamond{\mathbin{\scalebox{0.6}{$\Diamond$}}}%
  \def\bigcup{\mathop{\mathchoice{\raisebox{-0.3em}{\scalebox{1.7}{$\cup$}}}{\scalebox{1.25}{$\cup$}}{\cup}{\cup}}\displaylimits}%
  \def\bigcap{\mathop{\mathchoice{\raisebox{-0.3em}{\scalebox{1.7}{$\cap$}}}{\scalebox{1.25}{$\cap$}}{\cap}{\cap}}\displaylimits}%
  \def\lesseqqgtr{\mathrel{\lessgtr}}\def\bigoplus{\mathop{\oplus}\displaylimits}%
}
\providecommand{\ssi}{si et seulement si} % abréviation fréquente
\AtBeginDocument{\providecommand{\wideparen}{\overparen}} % arc de cercle

%% FIGFIX-BEGIN v5 — correctifs globaux des figures (docs/onbuch-generation/tools/figfix)
\usepackage{adjustbox}\usepackage{etoolbox}\tcbuselibrary{raster}
% toute figure TikZ/pgfplots plus large que la ligne est réduite à la largeur disponible
\BeforeBeginEnvironment{tikzpicture}{\begin{adjustbox}{max width=\linewidth}}
\AfterEndEnvironment{tikzpicture}{\end{adjustbox}}
% légendes pgfplots lisibles par-dessus les courbes
\pgfplotsset{every axis/.append style={legend style={fill=white,fill opacity=0.92,text opacity=1,draw=black!15,inner sep=3pt}}}
% étiquettes d'arêtes (syntaxe edge["..."]) avec fond blanc
\tikzset{every edge quotes/.append style={fill=white,inner sep=1.2pt}}
%% FIGFIX-END
\begin{document}

\pagegarde{Traduction : phrase causale, finalité, futur dans la subordonnée, emphase (révision)}{Espagnol}{Tle A}{Módulo 4 : Activités économiques et monde du travail}{E21}{2 h}{Leçon de traduction-révision consacrée à trois structures fondamentales de la phrase complexe : l'expression de la cause, de la finalité et le subjonctif dans les subordonnées de futur. Elle s'achève par une révision ciblée des procédés d'emphase, avec entraînement au thème et à la version dans l'esprit du Bac.
\vspace{4pt}
\begin{multicols}{2}\small
\begin{itemize}
\item À la fin de cette leçon, je sais traduire la cause en espagnol avec porque, como, puesto que, ya que et dado que en respectant leur place dans la phrase.
\item Je sais exprimer la cause avec les prépositions nominales a causa de et debido a suivies d'un nom ou d'un infinitif.
\item Je sais traduire la finalité avec para + infinitif et para que / a fin de que / con el fin de que + subjonctif selon que le sujet est identique ou différent.
\item Je sais employer le subjonctif présent après cuando, en cuanto, tan pronto como, hasta que et si (conditionnel) quand la subordonnée renvoie au futur.
\item Je sais choisir entre indicatif et subjonctif dans la subordonnée temporelle selon que l'action est habituelle/accomplie ou à venir.
\item Je sais réviser et mobiliser les structures d'emphase : lo + adjectif, ce que... c'est (lo que... es), la mise en relief par dédoublement et l'ordre des mots.
\end{itemize}\end{multicols}}

\begin{sommaire}
\begin{multicols}{2}
\begin{enumerate}
\item Observation guidée : repérer les structures
\item La phrase causale
\item L'expression de la finalité
\item Le futur dans la subordonnée : le subjonctif
\item Révision de l'emphase
\item Entraînement à la traduction : thème et version
\item Activité d'intégration
\item Méthodes et exercices résolus
\item Exercices
\item Corrigés des exercices
\item Fiche bilan
\end{enumerate}
\end{multicols}
\end{sommaire}

\begin{objectifs}
\begin{itemize}
\item À la fin de cette leçon, je sais traduire la cause en espagnol avec porque, como, puesto que, ya que et dado que en respectant leur place dans la phrase.
\item Je sais exprimer la cause avec les prépositions nominales a causa de et debido a suivies d'un nom ou d'un infinitif.
\item Je sais traduire la finalité avec para + infinitif et para que / a fin de que / con el fin de que + subjonctif selon que le sujet est identique ou différent.
\item Je sais employer le subjonctif présent après cuando, en cuanto, tan pronto como, hasta que et si (conditionnel) quand la subordonnée renvoie au futur.
\item Je sais choisir entre indicatif et subjonctif dans la subordonnée temporelle selon que l'action est habituelle/accomplie ou à venir.
\item Je sais réviser et mobiliser les structures d'emphase : lo + adjectif, ce que... c'est (lo que... es), la mise en relief par dédoublement et l'ordre des mots.
\item Je sais repérer et éviter les pièges francophones : « pour que » + indicatif, futur après « quand », calque de « à cause de ».
\item Je sais traduire un court texte du français vers l'espagnol et inversement en combinant cause, finalité et futur dans la subordonnée.
\end{itemize}
\end{objectifs}

\begin{prerequis}
\begin{itemize}
\item Le subjonctif présent : formation régulière et verbes irréguliers (ser, ir, haber, tener, hacer, poder, querer) — Première
\item L'expression de l'emphase : lo + adjectif, lo que, ordre des mots (leçon de Première et début de Tle)
\item Les temps de l'indicatif : présent, futur simple, conditionnel
\item Les connecteurs simples déjà vus : porque, para, cuando
\item La concordance des temps entre principale et subordonnée (notions)
\end{itemize}
\end{prerequis}

\newpage

\coursec{Observation guidée : repérer les structures}

\courssub{Situation d'accroche}

Imaginez la scène : il est seize heures au lycée de Yaoundé, les cours viennent de se terminer. Assis sous un manguier de la cour, un élève de Terminale A écrit un message vocal à son correspondant mexicain de Guadalajara. Il veut lui raconter sa vie, ses projets, ses efforts. Voici ce qu'il dit :

\esp{Je travaille dur porque quiero entrar a la universidad. Como mi padre es comerciante, me ayuda para que yo estudie tranquilo. Cuando termine el bachillerato, buscaré trabajo.}

« Je travaille dur parce que je veux entrer à l'université. Comme mon père est commerçant, il m'aide pour que j'étudie tranquillement. Quand j'aurai terminé le bac, je chercherai du travail. »

Sans le savoir, en trois phrases seulement, cet élève vient d'utiliser les trois grandes structures de la subordination qui tombent presque chaque année à l'épreuve de traduction du Baccalauréat A : la \cle{cause} (\esp{porque}, \esp{como}), la \cle{finalité} (\esp{para que}), et le \cle{futur dans la subordonnée temporelle} (\esp{cuando termine}). Mieux encore : il a choisi le bon mode verbal à chaque fois, sans y réfléchir.

\textbf{Problématique :} comment l'espagnol exprime-t-il la cause, la finalité et le futur dans la subordonnée, et pourquoi le \cle{subjonctif} apparaît-il dans deux de ces trois cas ? Cette leçon de révision systématise ces structures pour ne plus jamais les rater. Mais avant d'énoncer la moindre règle, observons.

\courssub{Lecture du texte support : un jeune entrepreneur de Douala}

Lisons maintenant un paragraphe un peu plus long. Il raconte la vie d'Andréa, une jeune élève de Douala qui vend des beignets et des jus de fruits pour préparer son avenir. Toutes les structures cibles de la leçon s'y trouvent. Lisez-le d'abord en silence, puis à voix haute.

\begin{textebox}[Una joven empresaria de Douala]
Como no tengo capital, trabajo los fines de semana para ahorrar dinero. Mi tío me presta su tienda del mercado para que venda mis productos. Ya que mis padres no pueden pagar mis estudios, vendo beignets y zumos de frutas a fin de que mi familia viva mejor. Cuando termine el bachillerato, abriré mi propio negocio, porque quiero ser independiente. En cuanto tenga suficiente dinero, compraré una máquina para hacer más zumos, con el fin de emplear a otros jóvenes del barrio. Dado que los clientes vienen cada mañana, debo levantarme temprano para preparar todo antes de las seis.
\end{textebox}

\textbf{Glossaire :}

\begin{center}
\begin{tabularx}{\linewidth}{|X|X|}\hline
\rowcolor{popblueL} \textbf{Español} & \textbf{Français} \\ \hline
el capital & le capital (l'argent de départ) \\ \hline
ahorrar & économiser, épargner \\ \hline
prestar & prêter \\ \hline
la tienda & la boutique, le magasin \\ \hline
los beignets & les beignets (mot emprunté au français) \\ \hline
el zumo & le jus (de fruits) \\ \hline
a fin de que & afin que \\ \hline
el negocio & le commerce, l'affaire \\ \hline
en cuanto & dès que \\ \hline
la máquina & la machine \\ \hline
emplear & employer, embaucher \\ \hline
el barrio & le quartier \\ \hline
dado que & étant donné que \\ \hline
levantarse & se lever \\ \hline
\end{tabularx}
\end{center}

\textbf{Traduction intégrale :} « Comme je n'ai pas de capital, je travaille les week-ends pour économiser de l'argent. Mon oncle me prête sa boutique du marché pour que je vende mes produits. Étant donné que mes parents ne peuvent pas payer mes études, je vends des beignets et des jus de fruits afin que ma famille vive mieux. Quand j'aurai terminé le bac, j'ouvrirai mon propre commerce, parce que je veux être indépendante. Dès que j'aurai assez d'argent, j'achèterai une machine pour faire plus de jus, dans le but d'employer d'autres jeunes du quartier. Étant donné que les clients viennent chaque matin, je dois me lever tôt pour tout préparer avant six heures. »

\courssub{Méthode d'observation : quatre gestes de traducteur}

Avant de classer les structures, apprenons \emph{comment} observer un texte en vue de la traduction. Cette démarche en quatre étapes est celle que vous appliquerez le jour du Bac, sur votre brouillon, avant d'écrire la moindre phrase.

\begin{methode}[Observer une structure grammaticale en quatre étapes]
\begin{enumerate}
\item \textbf{Repérer le connecteur.} Soulignez le mot qui relie les deux propositions : \esp{porque}, \esp{como}, \esp{para que}, \esp{cuando}... C'est lui qui porte le sens logique (cause, but, temps).
\item \textbf{Identifier le mode et le temps du verbe qui suit.} Après le connecteur, le verbe est-il à l'\cle{indicatif} ou au \cle{subjonctif} ? À quel temps ? C'est l'observation capitale.
\item \textbf{Comparer avec le français.} Comment le français dit-il la même chose ? Y a-t-il un décalage (par exemple le futur français contre le subjonctif espagnol) ?
\item \textbf{Formuler la règle.} En une phrase simple : « Après tel connecteur, dans tel contexte, l'espagnol exige tel mode. »
\end{enumerate}
\end{methode}

\courssub{Repérage guidé dans le texte}

Appliquons la méthode. Reprenons le texte d'Andréa et effectuons le repérage demandé : \textbf{soulignons} les marqueurs de cause, \textbf{encadrons} ceux de finalité et \textbf{entourons} les subordonnées temporelles renvoyant au futur.

\textbf{1. Les marqueurs de cause (à souligner) :}

\begin{itemize}
\item \esp{Como no tengo capital, trabajo los fines de semana.} → \esp{como} en tête de phrase = « comme, puisque ».
\item \esp{Ya que mis padres no pueden pagar mis estudios...} → \esp{ya que} = « étant donné que ».
\item \esp{...abriré mi propio negocio, porque quiero ser independiente.} → \esp{porque} = « parce que ».
\item \esp{Dado que los clientes vienen cada mañana...} → \esp{dado que} = « vu que, étant donné que ».
\end{itemize}

Observation décisive : après ces quatre connecteurs, le verbe est à l'\cle{indicatif} (\esp{tengo}, \esp{pueden}, \esp{quiero}, \esp{vienen}). La cause est présentée comme un fait réel, constaté.

\textbf{2. Les marqueurs de finalité (à encadrer) :}

\begin{itemize}
\item \esp{...trabajo los fines de semana para ahorrar dinero.} → \esp{para} + infinitif (même sujet : c'est Andréa qui travaille et qui économise).
\item \esp{Mi tío me presta su tienda para que venda mis productos.} → \esp{para que} + verbe \esp{venda} : changement de sujet (le oncle prête / Andréa vend).
\item \esp{...a fin de que mi familia viva mejor.} → \esp{a fin de que} + verbe \esp{viva}.
\item \esp{...con el fin de emplear a otros jóvenes.} → \esp{con el fin de} + infinitif (même sujet).
\end{itemize}

Observation décisive : quand le sujet change, le connecteur est suivi du \cle{subjonctif} (\esp{venda}, \esp{viva}). Quand le sujet ne change pas, on trouve l'\cle{infinitif} (\esp{ahorrar}, \esp{emplear}).

\textbf{3. Les subordonnées temporelles au futur (à entourer) :}

\begin{itemize}
\item \esp{Cuando termine el bachillerato, abriré mi propio negocio.} → \esp{cuando} + \esp{termine} (subjonctif présent), puis futur \esp{abriré} dans la principale.
\item \esp{En cuanto tenga suficiente dinero, compraré una máquina.} → \esp{en cuanto} + \esp{tenga} (subjonctif présent), puis futur \esp{compraré}.
\end{itemize}

Observation décisive : l'action de la subordonnée (\esp{termine}, \esp{tenga}) n'est \emph{pas encore réalisée} ; elle est simplement attendue. L'espagnol la met au subjonctif présent, là où le français met un futur antérieur (« quand j'\emph{aurai terminé} », « dès que j'\emph{aurai} »).

\courssub{Les questions déclencheuses}

Deux questions suffisent à faire émerger la règle centrale de la leçon. Posez-les-vous devant chaque phrase.

\textbf{Question 1 : quel mode suit \esp{para que} ?}

Regardez : \esp{Mi tío me presta su tienda para que \textbf{venda} mis productos.} Le verbe \esp{venda} n'est pas l'indicatif (\esp{vende}) : c'est le \textbf{subjonctif présent}. Pourquoi ? Parce que la vente n'est pas un fait accompli : c'est un \cle{but à atteindre}, une action projetée, virtuelle. Tout ce qui est but, souhait, visée relève du subjonctif en espagnol.

\textbf{Question 2 : quel temps suit \esp{cuando} quand il s'agit du futur ?}

Regardez : \esp{Cuando \textbf{termine} el bachillerato, abriré mi negocio.} Le français dirait « quand j'\emph{aurai terminé} » avec un futur antérieur. L'espagnol, lui, refuse le futur après \esp{cuando} : il exige le \textbf{subjonctif présent} \esp{termine}. Pourquoi ? Parce que l'action n'est pas encore réalisée au moment où l'on parle : elle est incertaine, suspendue à l'avenir.

\begin{attention}[Le piège n° 1 des francophones : le futur après « quand »]
En français, on dit : « \textbf{Quand je finirai} mes études, je chercherai du travail » ou « Quand j'\textbf{aurai fini}... ». Le français utilise donc le futur (ou le futur antérieur) après « quand ».

L'espagnol l'\textbf{interdit absolument}. Après \esp{cuando} (et \esp{en cuanto}, \esp{tan pronto como}, \esp{hasta que}, \esp{mientras}) suivi d'une action \emph{future non réalisée}, il faut le \cle{subjonctif présent} :

\esp{Cuando termine el bachillerato, buscaré trabajo.} « Quand j'aurai terminé le bac, je chercherai du travail. »

Jamais : \esp{Cuando terminaré...} ni \esp{Cuando terminará...} — ces formes sont des fautes graves, immédiatement sanctionnées à l'épreuve de traduction. Retenez le réflexe : \textbf{« quand » + futur en français = \esp{cuando} + subjonctif en espagnol.}
\end{attention}

\courssub{Première généralisation : la loi du subjonctif}

Mettons côte à côte les deux observations précédentes. Dans un cas, le but à atteindre (\esp{para que venda}) ; dans l'autre, l'action future non encore réalisée (\esp{cuando termine}). Le point commun saute aux yeux : dans les deux cas, l'action de la subordonnée \textbf{n'existe pas encore dans la réalité}. Elle est imaginée, souhaitée, attendue.

\begin{aretenir}[La grande loi d'observation de cette leçon]
Le \cle{subjonctif} apparaît dès que la subordonnée exprime un \textbf{but à atteindre} (finalité : \esp{para que}, \esp{a fin de que}) ou une \textbf{action future non réalisée} (subordonnée temporelle : \esp{cuando}, \esp{en cuanto}, \esp{tan pronto como}).

En revanche, la \textbf{cause} est un fait réel, constaté, expliqué : elle se construit à l'\cle{indicatif} (\esp{porque tengo}, \esp{como no tengo}, \esp{ya que mis padres no pueden}).

Résumé mnémotechnique : \textbf{cause = réel = indicatif ; but ou futur = virtuel = subjonctif.}
\end{aretenir}

\begin{exemplebox}[Deux exemples fondateurs, traduits et commentés]
\esp{Como no tengo capital, trabajo.} = « Comme je n'ai pas de capital, je travaille. »

Cause réelle, constatée → indicatif \esp{tengo}. Notez que \esp{como} causal se place \emph{en tête de phrase} ; en milieu de phrase, on préfère \esp{porque}.

\medskip

\esp{Cuando termine el bachillerato, abriré mi negocio.} = « Quand j'aurai terminé le bac, j'ouvrirai mon commerce. »

Action future non réalisée → subjonctif présent \esp{termine} dans la subordonnée, futur \esp{abriré} dans la principale. Le français, lui, utilise le futur antérieur « aurai terminé » : c'est le décalage majeur entre les deux langues.
\end{exemplebox}

\courssub{Classement des exemples relevés : le tableau à trois colonnes}

Rassemblons maintenant tous les exemples du texte d'Andréa dans le tableau de classement demandé. Ce tableau est votre premier outil de révision : recopiez-le dans votre cahier et apprenez à le reproduire de mémoire.

\begin{center}
\begin{tabularx}{\linewidth}{|X|X|X|}\hline
\rowcolor{popblueL} \textbf{CAUSA (indicativo)} & \textbf{FINALIDAD (infinitivo o subjuntivo)} & \textbf{FUTURO EN LA SUBORDINADA (subjuntivo)} \\ \hline
\esp{Como no tengo capital, trabajo los fines de semana.} & \esp{Trabajo para ahorrar dinero.} (même sujet → infinitif) & \esp{Cuando termine el bachillerato, abriré mi negocio.} \\ \hline
\esp{Ya que mis padres no pueden pagar mis estudios, vendo beignets.} & \esp{Mi tío me presta su tienda para que venda mis productos.} (changement de sujet → subjonctif) & \esp{En cuanto tenga dinero, compraré una máquina.} \\ \hline
\esp{Abriré mi negocio, porque quiero ser independiente.} & \esp{Vendo zumos a fin de que mi familia viva mejor.} (subjonctif) & \\ \hline
\esp{Dado que los clientes vienen cada mañana, debo levantarme temprano.} & \esp{Compraré una máquina con el fin de emplear a otros jóvenes.} (infinitif) & \\ \hline
\end{tabularx}
\end{center}

Ce tableau fait apparaître visuellement la loi dégagée plus haut : la colonne de la cause ne contient que de l'indicatif ; les deux autres colonnes contiennent du subjonctif (ou de l'infinitif quand le sujet ne change pas, pour la finalité).

\courssub{Carte mentale de synthèse}

\begin{popfigure}
\begin{center}\tcbox[colback=popink,colframe=popink,coltext=white,arc=9pt,boxsep=4pt,left=6pt,right=6pt]{\bfseries\small Subordinadas clave}\end{center}
\vspace{-2pt}
\begin{tcbraster}[raster columns=2,raster equal height=rows,raster column skip=6pt,raster row skip=6pt,enhanced,blankest]
\begin{tcolorbox}[enhanced,colback=white,colframe=popblue,boxrule=0.9pt,arc=3pt,title={CAUSA indicativo},fonttitle=\bfseries\scriptsize,coltitle=white,colbacktitle=popblue,left=4pt,right=4pt,top=2pt,bottom=2pt,boxsep=1pt,toptitle=1.5pt,bottomtitle=1.5pt,halign=left]
\begin{itemize}[nosep,leftmargin=*,itemsep=1pt,font=\scriptsize]
\item porque como
\item ya que puesto que
\item dado que a causa de
\end{itemize}
\end{tcolorbox}
\begin{tcolorbox}[enhanced,colback=white,colframe=poporange,boxrule=0.9pt,arc=3pt,title={FINALIDAD para + inf. para que + subj.},fonttitle=\bfseries\scriptsize,coltitle=white,colbacktitle=poporange,left=4pt,right=4pt,top=2pt,bottom=2pt,boxsep=1pt,toptitle=1.5pt,bottomtitle=1.5pt,halign=left]
\begin{itemize}[nosep,leftmargin=*,itemsep=1pt,font=\scriptsize]
\item para para que
\item a fin de que con el fin de
\end{itemize}
\end{tcolorbox}
\begin{tcolorbox}[enhanced,colback=white,colframe=popgreen,boxrule=0.9pt,arc=3pt,title={FUTURO EN LA SUBORDINADA subjuntivo},fonttitle=\bfseries\scriptsize,coltitle=white,colbacktitle=popgreen,left=4pt,right=4pt,top=2pt,bottom=2pt,boxsep=1pt,toptitle=1.5pt,bottomtitle=1.5pt,halign=left]
\begin{itemize}[nosep,leftmargin=*,itemsep=1pt,font=\scriptsize]
\item cuando en cuanto
\item tan pronto como hasta que
\end{itemize}
\end{tcolorbox}
\end{tcbraster}

\legende{Carte mentale des trois structures de la leçon : sur chaque branche, le mode verbal exigé est indiqué.}
\end{popfigure}

\courssub{Comparaison français / espagnol : les trois décalages à connaître}

\begin{center}
\begin{tabularx}{\linewidth}{|X|X|X|}\hline
\rowcolor{popblueL} \textbf{Français} & \textbf{Español} & \textbf{Remarque} \\ \hline
Parce qu'il pleut, je reste. & \esp{Porque llueve, me quedo.} & Indicatif des deux côtés : aucun piège. \\ \hline
Je travaille pour réussir. & \esp{Trabajo para aprobar.} & Même sujet → infinitif dans les deux langues. \\ \hline
Je travaille pour que tu réussisses. & \esp{Trabajo para que apruebes.} & Subjonctif dans les deux langues : le français aide ici. \\ \hline
Quand je finirai, je partirai. & \esp{Cuando termine, me iré.} & Piège majeur : futur français, mais subjonctif espagnol. \\ \hline
Dès que j'aurai de l'argent, je paierai. & \esp{En cuanto tenga dinero, pagaré.} & Futur antérieur français = subjonctif présent espagnol. \\ \hline
\end{tabularx}
\end{center}

La leçon de cette comparaison est claire : pour la cause et la finalité, le français est un bon allié (les modes coïncident souvent) ; pour le futur dans la subordonnée temporelle, le français est un \emph{traître} : il faut combattre le calque.

\courssub{Exercice résolu d'observation}

\begin{exoresolu}[Repérage et analyse]
Dans les phrases suivantes, extraites d'une lettre d'un élève de Douala à son correspondant espagnol, (a) repérez le connecteur, (b) indiquez le mode du verbe qui suit, (c) justifiez ce choix, (d) traduisez en français.

1. \esp{Como mi hermano trabaja en el puerto, me ayuda con los deberes.}

2. \esp{Estudio mucho para que mis padres estén orgullosos de mí.}

3. \esp{Cuando acaben las clases, viajaremos a Kribi.}

\tcblower

\textbf{1.} (a) Connecteur : \esp{como} (en tête de phrase). (b) Mode : indicatif présent (\esp{trabaja}). (c) Justification : il s'agit d'une \textbf{cause réelle}, un fait constaté — le frère travaille effectivement au port — donc indicatif. (d) Traduction : « Comme mon frère travaille au port, il m'aide avec mes devoirs. »

\textbf{2.} (a) Connecteur : \esp{para que}. (b) Mode : subjonctif présent (\esp{estén}). (c) Justification : il s'agit d'une \textbf{finalité} avec changement de sujet (j'étudie / mes parents sont fiers) ; le but n'est pas encore atteint, donc subjonctif. (d) Traduction : « J'étudie beaucoup pour que mes parents soient fiers de moi. »

\textbf{3.} (a) Connecteur : \esp{cuando}. (b) Mode : subjonctif présent (\esp{acaben}). (c) Justification : la fin des cours est une \textbf{action future non réalisée} ; l'espagnol interdit le futur après \esp{cuando}. (d) Traduction : « Quand les cours seront terminés, nous voyagerons à Kribi. » Notez le futur français « seront terminés » rendu par le subjonctif espagnol.
\end{exoresolu}

\courssub{À retenir de cette observation}

\begin{aretenir}[Bilan de l'observation guidée]
\begin{itemize}
\item La \textbf{cause} s'exprime par \esp{porque}, \esp{como} (en tête de phrase), \esp{ya que}, \esp{puesto que}, \esp{dado que}, \esp{a causa de}, \esp{debido a} : le verbe est à l'\cle{indicatif}, car la cause est un fait réel.
\item La \textbf{finalité} s'exprime par \esp{para} + infinitif (même sujet) ou \esp{para que}, \esp{a fin de que} + \cle{subjonctif} (changement de sujet), car le but est virtuel.
\item Le \textbf{futur dans la subordonnée temporelle} (\esp{cuando}, \esp{en cuanto}, \esp{tan pronto como}) exige le \cle{subjonctif présent} : le futur y est interdit, contrairement au français.
\item Loi générale : \textbf{réel → indicatif ; but ou futur non réalisé → subjonctif.}
\end{itemize}
\end{aretenir}

Ces observations posent les fondations de toute la leçon. Dans les sections suivantes, nous étudierons chaque structure en détail — forme, emploi, exceptions — avant de nous entraîner au thème et à la version. Mais dès maintenant, gardez en tête l'image d'Andréa à Douala : \esp{como} pour expliquer, \esp{para que} pour viser, \esp{cuando} + subjonctif pour attendre l'avenir.

\coursec{La phrase causale}

Dans la section précédente, nous avons observé un texte et repéré plusieurs structures : des propositions introduites par \esp{porque}, \esp{como}, \esp{ya que}, des compléments introduits par \esp{a causa de}, \esp{debido a}. Toutes ces structures répondaient à la même question : \emph{pourquoi ?} Nous allons maintenant les étudier une par une, de façon systématique, car la traduction de la cause est l'un des points les plus fréquents — et les plus piégeux — de l'épreuve de thème au Baccalauréat.

\begin{definition}[La phrase causale]
La \cle{phrase causale} est une phrase complexe dont la subordonnée (ou le complément circonstanciel) exprime la \cle{raison}, le \cle{motif}, la \cle{cause} de l'action exprimée dans la proposition principale. En espagnol, la cause peut être introduite par une \cle{conjonction} suivie d'un verbe conjugué (\esp{porque}, \esp{como}, \esp{puesto que}, \esp{ya que}, \esp{dado que}) ou par une \cle{locution prépositive} suivie d'un nom ou d'un infinitif (\esp{a causa de}, \esp{debido a}, \esp{gracias a}, \esp{por}).
\end{definition}

\courssub{Observation : un même fait, plusieurs formulations}

Lisons d'abord ce court texte original, inspiré de la vie scolaire à Yaoundé, et soulignons mentalement tout ce qui exprime une cause.

\begin{textebox}[Un lunes difícil en Yaoundé]
Como llovía mucho esta mañana, miles de alumnos llegaron tarde al liceo. Los bendskins no pasaban a causa de las inundaciones en las calles del barrio, y debido al tráfico bloqueado, muchos profesores tampoco pudieron llegar a tiempo. El director suspendió la primera hora de clase porque los alumnos estaban completamente mojados. Ya que el examen de español estaba previsto para las diez, los profesores decidieron mantenerlo, puesto que los alumnos lo habían preparado desde hacía dos semanas. Gracias a la calma del personal, la jornada terminó bien. Dado que la lluvia continuaba, el director autorizó a los alumnos a salir antes de las cuatro. Por falta de transporte, algunos estudiantes volvieron a pie hasta su casa.
\end{textebox}

\textbf{Glossaire :} \esp{llovía} : il pleuvait ; \esp{los bendskins} : les motos-taxis ; \esp{las inundaciones} : les inondations ; \esp{mojados} : trempés ; \esp{previsto} : prévu ; \esp{mantenerlo} : le maintenir ; \esp{a pie} : à pied ; \esp{la jornada} : la journée.

\textbf{Questions de repérage :}
\begin{enumerate}
\item Relevez tous les mots ou groupes de mots qui introduisent une cause.
\item Lesquels sont suivis d'un verbe conjugué ? Lesquels sont suivis d'un nom ?
\item Lesquels se placent en tête de phrase ? Lesquels après la proposition principale ?
\end{enumerate}

\begin{corrige}[Réponses de l'observation]
Conjonctions + verbe conjugué : \esp{como}, \esp{porque}, \esp{ya que}, \esp{puesto que}, \esp{dado que} ; locutions + nom (ou infinitif) : \esp{a causa de}, \esp{debido a}, \esp{gracias a}, \esp{por falta de}. On observe que \esp{porque} suit la principale, alors que \esp{como} ouvre la phrase. \esp{Ya que}, \esp{puesto que} et \esp{dado que} apparaissent ici en tête, mais acceptent les deux positions.
\end{corrige}

\courssub{Porque : la conjonction de cause par excellence}

\begin{propriete}[Porque]
\esp{Porque} est la conjonction de cause la plus courante. Elle répond à la question \esp{¿por qué ?} (pourquoi ?) et se place \cle{après la proposition principale}, comme « parce que » en français.

\esp{No salgo porque llueve.} « Je ne sors pas parce qu'il pleut. »

\esp{No fue a la escuela porque estaba enfermo.} « Il n'est pas allé à l'école parce qu'il était malade. »
\end{propriete}

La structure est donc : \textbf{[principale] + porque + [cause]}. C'est la structure la plus sûre et la plus fréquente à l'écrit comme à l'oral. En thème, « parce que » se traduit presque toujours par \esp{porque}.

Autres exemples :
\begin{itemize}
\item \esp{María no vino a clase porque tenía fiebre.} « Marie n'est pas venue en classe parce qu'elle avait de la fièvre. »
\item \esp{Los cameruneses aman el fútbol porque une a todo el país.} « Les Camerounais aiment le football parce qu'il unit tout le pays. »
\item \esp{No compré cacao porque el precio había subido.} « Je n'ai pas acheté de cacao parce que le prix avait augmenté. »
\end{itemize}

\begin{attention}[Porque ou por qué ?]
Ne confondez jamais :
\begin{itemize}
\item \esp{porque} (en un mot, sans accent) = « parce que », conjonction de cause : \esp{No vino porque estaba cansado.}
\item \esp{por qué} (en deux mots, avec accent) = « pourquoi », interrogation directe ou indirecte : \esp{¿Por qué no vino?} « Pourquoi n'est-il pas venu ? » ; \esp{No sé por qué no vino.} « Je ne sais pas pourquoi il n'est pas venu. »
\end{itemize}
C'est l'une des fautes d'orthographe les plus sanctionnées au Bac.
\end{attention}

\courssub{Como : la cause en tête de phrase}

\begin{propriete}[Como causal]
Placé \cle{en tête de phrase} (ou après une incise), \esp{como} exprime une cause présentée comme \cle{évidente}, déjà connue de l'interlocuteur. Il correspond au « comme » causal du français. La structure est : \textbf{Como + [cause], [conséquence]}.

\esp{Como llueve, no salgo.} « Comme il pleut, je ne sors pas. »

\esp{Como estaba enfermo, no fue a la escuela.} « Comme il était malade, il n'est pas allé à l'école. »
\end{propriete}

\textbf{Règle d'or :} \esp{como} causal ne se place \cle{jamais en fin de phrase}. On ne peut pas dire *\esp{No salgo como llueve} pour « Je ne sors pas parce qu'il pleut » : dans cette position, il faut \esp{porque}.

Autres exemples :
\begin{itemize}
\item \esp{Como no tenía dinero, tomó un bendskin.} « Comme il n'avait pas d'argent, il a pris une moto-taxi. »
\item \esp{Como el examen es mañana, los alumnos repasan toda la noche.} « Comme l'examen est demain, les élèves révisent toute la nuit. »
\item \esp{Como eres mi amigo, te ayudo.} « Comme tu es mon ami, je t'aide. »
\end{itemize}

Remarquez la correspondance parfaite avec le français : le français « comme » causal, lui aussi, se place en tête de phrase. C'est une chance pour les francophones : il suffit de calquer la structure.

\courssub{Puesto que, ya que, dado que : le registre soutenu}

\begin{propriete}[Les conjonctions soutenues]
\esp{Puesto que}, \esp{ya que} et \esp{dado que} introduisent une cause présentée comme un \cle{fait connu}, une évidence que l'on rappelle. Elles appartiennent à un registre plus soutenu que \esp{porque} et sont très appréciées dans les rédactions et les commentaires de texte au Bac. Elles acceptent \cle{les deux positions} (avant ou après la principale), mais privilégient le début de phrase.

\esp{Dado que el examen es difícil, hay que prepararse.} « Étant donné que l'examen est difficile, il faut se préparer. »

\esp{Ya que estás aquí, ayúdame.} « Puisque tu es là, aide-moi. »
\end{propriete}

Nuances de traduction en français :
\begin{itemize}
\item \esp{puesto que} = « puisque », « étant donné que » ;
\item \esp{ya que} = « puisque », « du moment que » ;
\item \esp{dado que} = « étant donné que », « vu que ».
\end{itemize}

Exemples :
\begin{itemize}
\item \esp{Puesto que el cacao es la riqueza de la región, el gobierno apoya a los campesinos.} « Puisque le cacao est la richesse de la région, le gouvernement soutient les paysans. »
\item \esp{No pudimos jugar, ya que el campo estaba inundado.} « Nous n'avons pas pu jouer, puisque le terrain était inondé. »
\item \esp{Dado que los alumnos habían trabajado mucho, el profesor les felicitó.} « Étant donné que les élèves avaient beaucoup travaillé, le professeur les a félicités. »
\end{itemize}

\begin{aretenir}[Positions des conjonctions causales]
\begin{itemize}
\item \esp{Porque} : uniquement \textbf{après} la principale.
\item \esp{Como} : uniquement \textbf{avant} la principale (tête de phrase).
\item \esp{Puesto que}, \esp{ya que}, \esp{dado que} : les deux positions, avec une préférence pour le début.
\end{itemize}
\end{aretenir}

\courssub{A causa de, debido a, gracias a, por : la cause en un complément}

Quand la cause est exprimée par un \cle{nom} (et non par une proposition avec verbe conjugué), l'espagnol utilise des locutions prépositives. Le choix dépend de la \cle{valeur} de la cause : négative, neutre ou positive.

\begin{propriete}[Les locutions prépositives de cause]
\begin{itemize}
\item \esp{A causa de} + nom ou infinitif : cause \cle{négative ou fâcheuse} (= « à cause de »).

\esp{A causa de la huelga, no hubo clases.} « À cause de la grève, il n'y a pas eu cours. »

\item \esp{Debido a} + nom ou infinitif : cause \cle{neutre}, factuelle, registre administratif ou journalistique (= « en raison de », « dû à »).

\esp{Debido a la lluvia, el partido se aplazó.} « En raison de la pluie, le match a été reporté. »

\item \esp{Gracias a} + nom ou infinitif : cause \cle{positive} (= « grâce à »).

\esp{Gracias a su esfuerzo, aprobó el examen.} « Grâce à son effort, il a réussi l'examen. »

\item \esp{Por} + nom ou infinitif : cause concise, très fréquente (= « par », « faute de », « à cause de » selon le contexte).

\esp{Por falta de dinero, no viajó.} « Faute d'argent, il n'a pas voyagé. »
\end{itemize}
\end{propriete}

Autres exemples :
\begin{itemize}
\item \esp{A causa de un accidente, la carretera de Douala estaba cerrada.} « À cause d'un accident, la route de Douala était fermée. »
\item \esp{El vuelo se canceló debido a la niebla.} « Le vol a été annulé en raison du brouillard. »
\item \esp{Gracias a los entrenamientos, el equipo ganó el partido.} « Grâce aux entraînements, l'équipe a gagné le match. »
\item \esp{Por la lluvia, no fuimos al mercado.} « À cause de la pluie, nous ne sommes pas allés au marché. »
\item \esp{A causa de no estudiar, suspendió el examen.} « Pour ne pas avoir étudié, il a échoué à l'examen. » (locution + infinitif)
\end{itemize}

\begin{attention}[Debido a, pas debido de]
On écrit toujours \esp{debido a} (avec la préposition \esp{a}), jamais *\esp{debido de}. De même, méfiez-vous de la construction *\esp{a causa de que} + indicatif : elle existe dans la langue parlée mais est à éviter à l'écrit ; préférez \esp{porque} ou \esp{ya que} devant une proposition complète. Enfin, rappelez-vous que le français « à cause de » implique une nuance négative : pour une cause positive, l'espagnol exige \esp{gracias a} — on ne dira jamais *\esp{a causa de su esfuerzo, aprobó}.
\end{attention}

\courssub{Schéma récapitulatif de la phrase causale}

\begin{popfigure}
\begin{center}
\begin{tikzpicture}[
  box/.style={draw=popblue, fill=popblueL, rounded corners, align=center, minimum height=0.9cm, inner sep=5pt, font=\small},
  cbox/.style={draw=poporange, fill=poporangeL, rounded corners, align=center, minimum height=0.9cm, inner sep=5pt, font=\small},
  lbl/.style={font=\small\itshape, text=popdark, midway, above},
  arr/.style={-{Stealth[length=3mm]}, thick, poppurple},
  node distance=1.8cm
]
\node[box, align=center] (cons){Conséquence\\ \esp{no salgo}};
\node[cbox, left=of cons, align=center] (causaL){Cause\\ \esp{llueve}};
\node[cbox, right=of cons, align=center] (causaR){Cause\\ \esp{llueve}};
\draw[arr] (causaL.east) -- node[lbl] {\esp{como / ya que / dado que}} (cons.west);
\draw[arr] (causaR.west) -- node[lbl] {\esp{porque}} (cons.east);
\node[font=\small, text=popmuted, align=center, below=1cm of cons] {La cause peut précéder (\esp{como}, \esp{ya que}\ldots)\\ ou suivre (\esp{porque}) la proposition principale.};
\end{tikzpicture}
\end{center}
\legende{La place du connecteur dans la phrase causale : \esp{como} ouvre la phrase, \esp{porque} la ferme.}
\end{popfigure}

\courssub{Tableau récapitulatif}

\begin{center}
\begin{tabularx}{\linewidth}{|l|l|l|X|}
\hline
\rowcolor{popblueL} \textbf{Connecteur} & \textbf{Place} & \textbf{Registre / valeur} & \textbf{Exemple traduit} \\
\hline
\esp{porque} & après la principale & courant, neutre & \esp{No vino porque estaba enfermo.} « Il n'est pas venu parce qu'il était malade. » \\
\hline
\esp{como} & tête de phrase & cause évidente & \esp{Como estaba enfermo, no vino.} « Comme il était malade, il n'est pas venu. » \\
\hline
\esp{puesto que} & les deux (début de préférence) & soutenu & \esp{Puesto que llueve, nos quedamos.} « Puisqu'il pleut, nous restons. » \\
\hline
\esp{ya que} & les deux (début de préférence) & soutenu & \esp{Ya que estás aquí, entra.} « Puisque tu es là, entre. » \\
\hline
\esp{dado que} & les deux (début de préférence) & soutenu, formel & \esp{Dado que es tarde, salimos.} « Étant donné qu'il est tard, nous partons. » \\
\hline
\esp{a causa de} + nom/inf. & les deux & cause négative & \esp{A causa de la huelga, no hubo clases.} « À cause de la grève, il n'y a pas eu cours. » \\
\hline
\esp{debido a} + nom/inf. & les deux & cause neutre, factuelle & \esp{Debido a la lluvia, se aplazó el partido.} « En raison de la pluie, le match a été reporté. » \\
\hline
\esp{gracias a} + nom/inf. & les deux & cause positive & \esp{Gracias a ti, aprobé.} « Grâce à toi, j'ai réussi. » \\
\hline
\esp{por} + nom/inf. & les deux & cause concise & \esp{Por falta de tiempo, no terminé.} « Faute de temps, je n'ai pas fini. » \\
\hline
\end{tabularx}
\end{center}

\courssub{Comparaison français / espagnol}

\begin{center}
\begin{tabularx}{\linewidth}{|X|X|X|}
\hline
\rowcolor{popblueL} \textbf{Français} & \textbf{Espagnol} & \textbf{Remarque} \\
\hline
parce que & \esp{porque} & après la principale, dans les deux langues \\
\hline
comme (cause) & \esp{como} & en tête de phrase, dans les deux langues \\
\hline
puisque / étant donné que & \esp{puesto que}, \esp{ya que}, \esp{dado que} & registre soutenu \\
\hline
à cause de + nom & \esp{a causa de} + nom & nuance négative dans les deux langues \\
\hline
à cause de + proposition & \esp{porque}, \esp{ya que} & éviter *\esp{a causa de que} à l'écrit \\
\hline
grâce à & \esp{gracias a} & cause positive uniquement \\
\hline
en raison de / dû à & \esp{debido a} & jamais *\esp{debido de} \\
\hline
faute de / par manque de & \esp{por falta de} & très utile en thème \\
\hline
\end{tabularx}
\end{center}

\begin{savaistu}[« À cause de » au Bac]
À l'épreuve de thème, « à cause de » suivi d'une \emph{proposition} (sujet + verbe conjugué) se traduit presque toujours par \esp{porque} ou \esp{ya que}, et non par \esp{a causa de}. Ainsi : « Il a échoué à cause de sa paresse » peut se traduire de deux façons : \esp{Fracasó a causa de su pereza} (avec un nom) ou \esp{Fracasó porque era perezoso} (avec une proposition). Les deux sont correctes : à vous de choisir la structure que vous maîtrisez le mieux. Savoir passer du nom à la proposition, et inversement, est une compétence précieuse de traducteur.
\end{savaistu}

\courssub{Exercice résolu}

\begin{exoresolu}[Thème : traduisez en espagnol]
Énoncé : Traduisez les phrases suivantes en espagnol, en choisissant le connecteur causal approprié.

1. Comme il pleuvait, les élèves sont restés dans la salle de classe.

2. Je n'ai pas acheté de mangues parce qu'elles étaient trop chères.

3. À cause de la grève des transports, beaucoup d'élèves sont arrivés en retard.

4. Grâce à son travail, elle a obtenu une bourse pour étudier en Espagne.

5. Puisque tu connais bien Douala, montre-nous le chemin.

\tcblower
Solution détaillée :

1. La cause est en tête de phrase, présentée comme évidente : \esp{como}. « Il pleuvait » : imparfait \esp{llovía}. \esp{Como llovía, los alumnos se quedaron en el aula.}

2. « Parce que » après la principale : \esp{porque}. « Elles étaient trop chères » : imparfait \esp{eran} (ou \esp{costaban}). \esp{No compré mangos porque eran demasiado caros.}

3. Cause négative suivie d'un nom : \esp{a causa de}. \esp{A causa de la huelga de transportes, muchos alumnos llegaron tarde.}

4. Cause positive : \esp{gracias a}. « Une bourse » : \esp{una beca}. \esp{Gracias a su trabajo, obtuvo una beca para estudiar en España.}

5. « Puisque » : \esp{ya que} ou \esp{puesto que}, en tête de phrase. « Montre-nous » : impératif \esp{muéstranos}. \esp{Ya que conoces bien Douala, muéstranos el camino.}
\end{exoresolu}

\begin{attention}[Les trois pièges du francophone]
\begin{enumerate}
\item \textbf{La place de \esp{como}} : *\esp{No salgo como llueve} est une faute ; en fin de phrase, utilisez \esp{porque}.
\item \textbf{L'orthographe} : \esp{porque} (cause) / \esp{por qué} (pourquoi) ; \esp{debido a}, jamais *\esp{debido de}.
\item \textbf{La valeur de la cause} : \esp{a causa de} (négatif), \esp{debido a} (neutre), \esp{gracias a} (positif). Traduire « grâce à » par \esp{a causa de} inverse le sens de la phrase.
\end{enumerate}
\end{attention}

Maintenant que nous savons exprimer \emph{pourquoi} une action se produit, nous allons apprendre, dans la section suivante, à exprimer \emph{dans quel but} elle se produit : ce sera l'expression de la finalité avec \esp{para}, \esp{para que} et le subjonctif.

\coursec{L'expression de la finalité}

Après avoir étudié comment exprimer la \textbf{cause} (le « pourquoi » rétrospectif : \esp{porque, como, ya que...}), nous abordons maintenant la \textbf{finalité} (le « pour quoi » prospectif : le but visé). En espagnol, le choix de la structure dépend presque exclusivement de l'identité ou de la différence des \textbf{sujets} entre la proposition principale et la proposition de but. C'est un point cardinal du programme de Terminale et une source majeure d'erreurs au Bac : la confusion entre \esp{para + infinitif} et \esp{para que + subjonctif}.

\courssub{Observation guidée : même sujet ou sujet différent ?}

Lisez attentivement les deux phrases suivantes, extraites d'un reportage sur l'entrepreneuriat jeune à Yaoundé et à Madrid :

\begin{textebox}[Reportage : jeunes entrepreneurs]
Marta \textbf{estudia} ingeniería informática \textbf{para crear} su propia \emph{startup}.
Marta \textbf{estudia} ingeniería informática \textbf{para que sus padres estén} tranquilos.
\end{textebox}

\noindent\textbf{Traduction :}
\begin{itemize}
    \item \textit{Marta étudie l'ingénierie informatique pour créer sa propre startup.}
    \item \textit{Marta étudie l'ingénierie informatique pour que ses parents soient tranquilles.}
\end{itemize}

\begin{experience}[Analyse en binôme]
En binôme, répondez à ces questions :
\begin{enumerate}
    \item Qui est le sujet du verbe \esp{estudia} ? Qui est le sujet de \esp{crear} ? Sont-ils identiques ?
    \item Qui est le sujet du verbe \esp{estudia} ? Qui est le sujet de \esp{estén} ? Sont-ils identiques ?
    \item Quel mode verbal suit \esp{para} dans la première phrase ? Quel mode suit \esp{para que} dans la seconde ?
    \item Formulez la règle générale en français.
\end{enumerate}
\end{experience}

\courssub{Cas 1 — Même sujet : \texorpdfstring{\esp{para + infinitif}}{para + infinitif}}

\begin{propriete}[Règle fondamentale : identité des sujets]
Lorsque le \textbf{sujet de l'action principale} est \textbf{le même} que celui de l'action visée (le but), on utilise la préposition \cle{para} suivie de l'\textbf{infinitif}.
\medskip

\noindent\textbf{Structure :} \esp{Sujet + Verbe principal + para + Infinitif}
\medskip

\noindent\emph{Equivalent français :} « pour + infinitif » (sujet sous-entendu identique).
\end{propriete}

\begin{exemplebox}[Exemples variés - Même sujet]
\begin{itemize}
    \item \esp{Ahorro dinero \textbf{para comprar} una moto.} \rightarrow « J'économise de l'argent pour acheter une moto. »
    \item \esp{Los agricultores trabajan duro \textbf{para cosechar} buen cacao.} \rightarrow « Les agriculteurs travaillent dur pour récolter du bon cacao. »
    \item \esp{Vine a Douala \textbf{para buscar} trabajo.} \rightarrow « Je suis venu à Douala pour chercher du travail. »
    \item \esp{Estudiamos español \textbf{para comunicarnos} con Guinea Ecuatorial.} \rightarrow « Nous étudions l'espagnol pour communiquer avec la Guinée équatoriale. »
\end{itemize}
\end{exemplebox}

\begin{aretenir}[Variantes soutenues : \texorpdfstring{\esp{a fin de + infinitif}}{a fin de + infinitif} / \texorpdfstring{\esp{con el fin de + infinitif}}{con el fin de + infinitif}]
Dans un registre plus formel (rédaction administrative, discours officiel, presse, \textbf{épreuves de traduction au Bac}), on peut remplacer \esp{para + infinitif} par :
\begin{itemize}
    \item \esp{a fin de + infinitif} : \esp{El gobierno invierte \textbf{a fin de modernizar} las infraestructuras.}
    \item \esp{con el fin de + infinitif} : \esp{La empresa capacita a su personal \textbf{con el fin de mejorar} la productividad.}
\end{itemize}
\emph{Attention :} Ces locutions s'emploient \textbf{uniquement} si les sujets sont identiques. Elles ne s'utilisent \textbf{jamais} avec un sujet différent (pas de *« a fin de que » + infinitif).
\end{aretenir}

\courssub{Cas 2 — Sujets différents : \texorpdfstring{\esp{para que + subjonctif}}{para que + subjonctif}}

\begin{propriete}[Règle fondamentale : différence des sujets]
Lorsque le \textbf{sujet de l'action principale} est \textbf{différent} de celui de l'action visée (le but), on utilise la conjonction \cle{para que} suivie du \textbf{subjonctif} (généralement présent, parfois imparfait pour le passé).
\medskip

\noindent\textbf{Structure :} \esp{Sujet 1 + Verbe principal + para que + Sujet 2 + Subjonctif}
\medskip

\noindent\emph{Equivalent français :} « pour que + subjonctif ».
\end{propriete}

\begin{exemplebox}[Exemples variés - Sujets différents]
\begin{itemize}
    \item \esp{El profesor explica la lección \textbf{para que los alumnos la entiendan}.} \\
    (Sujet 1 = \esp{el profesor} ; Sujet 2 = \esp{los alumnos}) \\
    $\rightarrow$ « Le professeur explique la leçon pour que les élèves la comprennent. »
    \item \esp{Trabajo muchas horas \textbf{para que mis hijos puedan} estudiar en la universidad.} \\
    $\rightarrow$ « Je travaille beaucoup d'heures pour que mes enfants puissent étudier à l'université. »
    \item \esp{El Ministerio de Empleo crea programas \textbf{para que los jóvenes no emigren}.} \\
    $\rightarrow$ « Le Ministère de l'Emploi crée des programmes pour que les jeunes n'émigrent pas. »
    \item \esp{Te doy mi número \textbf{para que me llames} cuando llegues a Madrid.} \\
    $\rightarrow$ « Je te donne mon numéro pour que tu m'appelles quand tu arrives à Madrid. »
\end{itemize}
\end{exemplebox}

\begin{aretenir}[Variantes soutenues : \texorpdfstring{\esp{a fin de que}}{a fin de que} / \texorpdfstring{\esp{con el fin de que}}{con el fin de que} + subjonctif]
Pour varier le registre (style soutenu, traduction de textes officiels, essais), on utilise :
\begin{itemize}
    \item \esp{a fin de que + subjonctif} : \esp{Se aprueban leyes \textbf{a fin de que se respeten} los derechos humanos.}
    \item \esp{con el fin de que + subjonctif} : \esp{Organizamos talleres \textbf{con el fin de que las mujeres adquieran} autonomía económica.}
\end{itemize}
Ces formes sont \textbf{très appréciées des correcteurs du Bac} pour éviter la répétition de \esp{para que}.
\end{aretenir}

\courssub{Tableau récapitulatif : Choix de la structure}

\begin{center}
\begin{tabularx}{\linewidth}{|X|X|X|X|}
\hline
\rowcolor{popblueL}
\textbf{Situation} & \textbf{Structure espagnole} & \textbf{Mode / Forme} & \textbf{Équivalent français} \\
\hline
Même sujet & \esp{para + infinitif} & Infinitif & pour + infinitif \\
\hline
Même sujet (soutenu) & \esp{a fin de + infinitif} / \esp{con el fin de + infinitif} & Infinitif & dans le but de + infinitif / afin de + infinitif \\
\hline
Sujets différents & \esp{para que + subjonctif} & Subjonctif (présent/imparfait) & pour que + subjonctif \\
\hline
Sujets différents (soutenu) & \esp{a fin de que + subjonctif} / \esp{con el fin de que + subjonctif} & Subjonctif (présent/imparfait) & afin que + subjonctif / dans le but que + subjonctif \\
\hline
\end{tabularx}
\end{center}

\courssub{Organigramme de décision : comment choisir ?}

\begin{popfigure}
\begin{tikzpicture}[
    node distance=1.2cm and 2cm,
    decision/.style={diamond, draw=popblue, thick, fill=popblueL, text width=3.5cm, align=center, inner sep=4pt},
    process/.style={rectangle, draw=popgreen, thick, fill=popgreenL, text width=3.5cm, align=center, inner sep=4pt, rounded corners=4pt},
    start/.style={ellipse, draw=poporange, thick, fill=poporangeL, text width=2.5cm, align=center, inner sep=4pt},
    arrow/.style={-Stealth, thick, draw=popdark},
    yesno/.style={midway, fill=white, inner sep=2pt, font=\small\bfseries}
]
\node[start, align=center] (start){Je veux exprimer\\ un BUT / une FINALITÉ};
\node[decision, below of=start, align=center] (dec){Le sujet de la principale\\ est-il le MÊME que celui\\ du but ?};
\node[process, below left=1.5cm and 1cm of dec, align=center] (same){\esp{para + infinitif}\\\textit{(ou \esp{a fin de / con el fin de + inf.})}};
\node[process, below right=1.5cm and 1cm of dec, align=center] (diff){\esp{para que + subjonctif}\\\textit{(ou \esp{a fin de que / con el fin de que + subj.})}};

\draw[arrow] (start) -- (dec);
\draw[arrow] (dec) -- node[yesno, left] {OUI} (same);
\draw[arrow] (dec) -- node[yesno, right] {NON} (diff);
\end{tikzpicture}
\legende{Organigramme de décision pour l'expression de la finalité}
\end{popfigure}

\courssub{Point de grammaire avancée : \texorpdfstring{\esp{que + subjonctif}}{que + subjonctif} après verbes d'ordre / injonction (complément Bac)}

Dans la langue soutenue et littéraire, après des verbes exprimant un \textbf{ordre, une prière, une exigence, un conseil} (\esp{mandar, ordenar, pedir, rogar, exigir, aconsejar, sugerir}), la conjonction \esp{para que} peut s'omettre : on utilise directement \esp{que + subjonctif}. C'est l'\textbf{ellipse de la finalité}.

\begin{propriete}[Ellipse de \texorpdfstring{\esp{para que}}{para que} après verbes de volonté]
\begin{itemize}
    \item \esp{Le ordené \textbf{que se fuera} inmediatamente.} = \esp{Le ordené \textbf{para que se fuera}...} $\rightarrow$ « Je lui ai ordonné de partir / pour qu'il parte immédiatement. »
    \item \esp{Te ruego \textbf{que me escuches}.} = \esp{Te ruego \textbf{para que me escuches}.} $\rightarrow$ « Je t'en prie, écoute-moi / pour que tu m'écoutes. »
    \item \esp{El director sugirió \textbf{que contratáramos} más personal.} $\rightarrow$ « Le directeur a suggéré que nous embauchions plus de personnel. »
\end{itemize}
\emph{Remarque :} En français, on utilise souvent l'infinitif (« Je lui ai ordonné de partir ») là où l'espagnol maintient la subordonnée subjonctive (\esp{que se fuera}). Au Bac, en version, traduisez par « pour que + subjonctif » ou infinitif selon le contexte ; en thème, si le sujet change après \esp{mandar/pedir/rogar}, mettez \esp{que + subjonctif}.
\end{propriete}

\courssub{Comparaison systématique Français / Espagnol : le piège de l'indicatif}

C'est \textbf{la} erreur n°1 des francophones. En français, l'oreille hésite parfois : « Il travaille pour qu'il \textbf{peut} / \textbf{puisse} partir ». L'espagnol, lui, \textbf{n'hésite jamais} : après \esp{para que} (et ses variantes), c'est \textbf{obligatoirement le subjonctif}.

\begin{attention}[Erreur fatale : *« para que + indicatif / futur »]
\emph{Jamais} d'indicatif présent, de futur, de conditionnel après \esp{para que}, \esp{a fin de que}, \esp{con el fin de que}.

\begin{center}
\begin{tabularx}{\linewidth}{|X|X|}
\hline
\rowcolor{poporangeL}
\textbf{\color{poporange} INCORRECT (Calque français)} & \textbf{\color{popgreen} CORRECT (Espagnol standard)} \\
\hline
\esp{*Trabajo para que él \textbf{puede} irse.} & \esp{Trabajo para que él \textbf{pueda} irse.} \\
\esp{*Estudiamos para que \textbf{aprobaréis}.} & \esp{Estudiamos para que \textbf{aprobarais}.} \\
\esp{*Lo hago para que \textbf{será} más fácil.} & \esp{Lo hago para que \textbf{sea} más fácil.} \\
\esp{*Te llamo para que \textbf{sabrás} la verdad.} & \esp{Te llamo para que \textbf{sepas} la verdad.} \\
\hline
\end{tabularx}
\end{center}

\noindent\textbf{Règle d'or :} \cle{Para que} appelle \cle{le subjonctif}. Point final. Apprenez les formes du subjonctif présent des verbes irréguliers fréquents (\esp{tener $\rightarrow$ tenga, poder $\rightarrow$ pueda, poner $\rightarrow$ ponga, saber $\rightarrow$ sepa, haber $\rightarrow$ haya, ir $\rightarrow$ vaya, ser $\rightarrow$ sea, estar $\rightarrow$ esté}).
\end{attention}

\courssub{Méthode : Traduire « pour / pour que » du français vers l'espagnol (Thème)}

\begin{methode}[Algorithme de traduction « pour / pour que » $\rightarrow$ Espagnol]
\begin{enumerate}
    \item \textbf{Identifiez le verbe principal} (l'action faite) et le \textbf{verbe de but} (l'action visée).
    \item \textbf{Comparez les sujets} : \textit{Qui fait l'action principale ? Qui fait l'action de but ?}
    \item \textbf{Appliquez la règle :}
    \begin{itemize}
        \item \textbf{Même sujet} $\rightarrow$ \esp{para + infinitif} (ou \esp{a fin de / con el fin de + infinitif} si registre soutenu).
        \item \textbf{Sujets différents} $\rightarrow$ \esp{para que + subjonctif} (ou \esp{a fin de que / con el fin de que + subjonctif} si registre soutenu).
    \end{itemize}
    \item \textbf{Conjuguez le subjonctif} correctement (présent si principale au présent/futur/impératif ; imparfait si principale au passé/conditionnel).
    \item \textbf{Vérifiez :} Pas d'indicatif après \esp{que} ! Pas d'infinitif si sujets différents !
\end{enumerate}
\end{methode}

\begin{exoresolu}[Thème : Traduire en espagnol]
\textbf{Phrase :} « Le gouvernement camerounais construit des routes \textbf{pour que} les producteurs de cacao \textbf{puissent} transporter leur récolte vers Douala. »

\textbf{Résolution pas à pas :}
\begin{enumerate}
    \item Verbe principal : \esp{construir} (construire). Sujet : \esp{el gobierno camerunés}.
    \item Verbe de but : \esp{poder transportar} (pouvoir transporter). Sujet : \esp{los productores de cacao}.
    \item \textbf{Sujets différents ?} OUI (\esp{gobierno} $\neq$ \esp{productores}).
    \item Structure : \esp{para que + subjonctif}.
    \item Temps : Principale au présent (\esp{construye}) $\rightarrow$ Subjonctif \textbf{présent}.
    \item Subjonctif de \esp{poder} (3e pers. pl.) : \esp{puedan}. Subjonctif de \esp{transportar} : \esp{transporten}.
    \item Vocabulaire : \esp{carretera} (route), \esp{cosecha} (récolte).
\end{enumerate}

\textbf{Solution :}
\esp{El gobierno camerunés construye carreteras \textbf{para que los productores de cacao puedan transportar} su cosecha hacia Douala.}

\textbf{Variante soutenue (Bac) :}
\esp{El gobierno camerunés construye carreteras \textbf{a fin de que los productores de cacao puedan transportar} su cosecha hacia Douala.}
\end{exoresolu}

\courssub{Exercice d'application immédiate}

\begin{exercice}[Mettez le verbe à la forme correcte (Infinitif ou Subjonctif Présent)]
Complétez les phrases suivantes en conjuguant le verbe entre parenthèses.
\begin{enumerate}
    \item Los estudiantes leen el periódico \textbf{para \_\_\_\_\_\_\_\_} (estar) informados.
    \item El profesor repite la explicación \textbf{para que todos \_\_\_\_\_\_\_\_} (entender).
    \item Mi padre trabaja mucho \textbf{para \_\_\_\_\_\_\_\_} (pagar) la hipoteca.
    \item Te envío el correo \textbf{para que tú \_\_\_\_\_\_\_\_} (tener) la dirección.
    \item La ONG reparte semillas \textbf{a fin de que los campesinos \_\_\_\_\_\_\_\_} (sembrar) maíz.
    \item Ahorro cada mes \textbf{con el fin de \_\_\_\_\_\_\_\_} (comprar) un coche.
    \item El entrenador grita \textbf{para que los jugadores \_\_\_\_\_\_\_\_} (correr) más rápido.
    \item Voy al mercado \textbf{para \_\_\_\_\_\_\_\_} (comprar) plátanos y mango.
\end{enumerate}
\end{exercice}

\begin{corrige}[Exercice : Finalité]
\begin{enumerate}
    \item \textbf{estar} (Même sujet : \esp{los estudiantes} = sujet de \esp{estar}) $\rightarrow$ \esp{para estar}.
    \item \textbf{entiendan} (Sujets différents : \esp{el profesor} $\neq$ \esp{todos}) $\rightarrow$ \esp{para que todos entiendan}.
    \item \textbf{pagar} (Même sujet : \esp{mi padre} = sujet de \esp{pagar}) $\rightarrow$ \esp{para pagar}.
    \item \textbf{tengas} (Sujets différents : \esp{yo} (sujeito de \esp{envío}) $\neq$ \esp{tú}) $\rightarrow$ \esp{para que tú tengas}.
    \item \textbf{siembren} (Sujets différents : \esp{la ONG} $\neq$ \esp{los campesinos}) $\rightarrow$ \esp{a fin de que siembren}.
    \item \textbf{comprar} (Même sujet : \esp{yo} (sujeito de \esp{ahorro}) = sujet de \esp{comprar}) $\rightarrow$ \esp{con el fin de comprar}.
    \item \textbf{corran} (Sujets différents : \esp{el entrenador} $\neq$ \esp{los jugadores}) $\rightarrow$ \esp{para que corran}.
    \item \textbf{comprar} (Même sujet : \esp{yo} (sujeito de \esp{voy}) = sujet de \esp{comprar}) $\rightarrow$ \esp{para comprar}.
\end{enumerate}
\end{corrige}

\courssub{Le saviez-vous ? : Registre et stratégie au Bac}

\begin{savaistu}[Variation lexicale et réussite au Bac]
Dans les épreuves de \textbf{thème} et de \textbf{version} du Baccalauréat camerounais (série A), l'utilisation exclusive de \esp{para / para que} est perçue comme un niveau « seuil » (niveau B1). Pour atteindre le niveau « avancé » (B2/C1) et valoriser votre copie :
\begin{itemize}
    \item Remplacez systématiquement \textbf{un} \esp{para que} sur deux ou trois par \esp{a fin de que} ou \esp{con el fin de que}.
    \item Remplacez \textbf{un} \esp{para + infinitif} par \esp{a fin de + infinitif} ou \esp{con el fin de + infinitif}.
    \item Exemple : \esp{El Estado invierte en educación \textbf{para que} los jóvenes encuentren trabajo.} $\rightarrow$ \esp{El Estado invierte en educación \textbf{a fin de que} los jóvenes encuentren trabajo.}
    \item Cela montre votre maîtrise de la \textbf{variation stylistique}, critère explicite des grilles d'évaluation MINESEC.
\end{itemize}
\emph{Note culturelle :} En Guinée équatoriale (seul pays africain hispanophone), l'administration utilise massivement \esp{a fin de que} et \esp{con el fin de que} dans les décrets et communiqués officiels (Boletín Oficial del Estado). Les connaître, c'est aussi comprendre la francophonie... hispanophone !
\end{savaistu}

\courssub{Texte d'entraînement complet : Version / Thème}

\begin{textebox}[Article d'opinion : La formation professionnelle au Cameroun]
\textbf{Invertir en formación: una necesidad urgente}

El Gobierno camerunés ha lanzado un ambicioso plan de formación profesional \textbf{para que} miles de jóvenes adquieran competencias técnicas reales. Actualmente, muchos titulados universitarios vagan por las calles de Yaoundé y Douala \textbf{para buscar} un empleo que no existe, \textbf{mientras que} las empresas extranjeras contratan mano de obra cualificada del extranjero \textbf{para ejecutar} grandes proyectos de infraestructura.

Los centros de formación técnica se están modernizando \textbf{a fin de que} los aprendices dominen oficios como la soldadura, la electricidad industrial o la informática aplicada. Los formadores insisten en la disciplina \textbf{con el fin de que} los futuros técnicos sean competitivos en el mercado regional de la CEMAC.

Sin embargo, persisten obstáculos: las familias a veces presionan a sus hijos \textbf{para que} elijan carreras « prestigiosas » pero saturadas, \textbf{en lugar de} orientarlos hacia oficios manuales bien remunerados. Es vital que la sociedad cambie de mentalidad \textbf{para que} el país reduzca su dependencia exterior \textbf{y para} construir una economía soberana.
\end{textebox}

\noindent\textbf{Glossaire :}
\begin{itemize}
    \item \esp{vagar} : errer, chômer, chercher sans trouver.
    \item \esp{titulado} : diplômé.
    \item \esp{mano de obra cualificada} : main-d'œuvre qualifiée.
    \item \esp{aprendiz} : apprenti.
    \item \esp{soldadura} : soudure.
    \item \esp{saturado} : saturé (marché du travail).
    \item \esp{en lugar de} : au lieu de (+ infinitif).
    \item \esp{soberano} : souverain.
\end{itemize}

\begin{experience}[Activité de traduction guidée (Version $\rightarrow$ Thème)]
\begin{enumerate}
    \item \textbf{Version :} Traduisez le texte ci-dessus en français soigné (registre journalistique).
    \item \textbf{Analyse :} Relevez \textbf{toutes} les structures de finalité du texte (il y en a 7). Pour chacune, justifiez le choix (même sujet / sujets différents ; registre courant / soutenu).
    \item \textbf{Thème inverse :} Re-traduisez votre version française vers l'espagnol \textbf{sans regarder l'original}, en variant les structures : si l'original a \esp{para que}, mettez \esp{a fin de que} ; si l'original a \esp{para + inf.}, mettez \esp{con el fin de + inf.}, etc.
\end{enumerate}
\end{experience}

\courssub{Synthèse visuelle : Carte mentale de la finalité}

\begin{popfigure}
\begin{center}\tcbox[colback=popink,colframe=popink,coltext=white,arc=9pt,boxsep=4pt,left=6pt,right=6pt]{\bfseries\small FINALITÉ (EL PROPÓSITO)}\end{center}
\vspace{-2pt}
\begin{tcbraster}[raster columns=2,raster equal height=rows,raster column skip=6pt,raster row skip=6pt,enhanced,blankest]
\begin{tcolorbox}[enhanced,colback=popgreen,colframe=popgreen,coltext=white,boxrule=0.9pt,arc=3pt,left=4pt,right=4pt,top=3pt,bottom=3pt,boxsep=1pt,halign=left]\bfseries\scriptsize MÊME SUJET \esp{para + INF} \esp{a fin de + INF} \esp{con el fin de + INF}\end{tcolorbox}
\begin{tcolorbox}[enhanced,colback=poporange,colframe=poporange,coltext=white,boxrule=0.9pt,arc=3pt,left=4pt,right=4pt,top=3pt,bottom=3pt,boxsep=1pt,halign=left]\bfseries\scriptsize SUJETS DIFFÉRENTS \esp{para que + SUBJ} \esp{a fin de que + SUBJ} \esp{con el fin de que + SUBJ}\end{tcolorbox}
\begin{tcolorbox}[enhanced,colback=poppurple,colframe=poppurple,coltext=white,boxrule=0.9pt,arc=3pt,left=4pt,right=4pt,top=3pt,bottom=3pt,boxsep=1pt,halign=left]\bfseries\scriptsize ELLIPSE APRÈS VOLONTÉ \esp{que + SUBJ} (\esp{mandar, pedir, rogar...})\end{tcolorbox}
\begin{tcolorbox}[enhanced,colback=poppink,colframe=poppink,coltext=white,boxrule=0.9pt,arc=3pt,left=4pt,right=4pt,top=3pt,bottom=3pt,boxsep=1pt,halign=left]\bfseries\scriptsize PIÈGES FRANCOPHONES \textcolor{poporange}{$\times$ para que + IND} \textcolor{popgreen}{$\checkmark$ para que + SUBJ} \textcolor{poporange}{$\times$ a fin de + QUE + INF}\end{tcolorbox}
\end{tcbraster}

\legende{Carte mentale récapitulative : L'expression de la finalité}
\end{popfigure}

\begin{aretenir}[Check-list avant l'épreuve (Bac / DS)]
\begin{itemize}
    \item [$\square$] Je sais identifier le sujet de la principale et le sujet du but.
    \item [$\square$] Je conjugue sans hésiter le \textbf{subjonctif présent} des verbes irréguliers clés (\esp{tener, poder, poner, saber, haber, ir, ser, estar, dar}).
    \item [$\square$] Je n'écris \textbf{jamais} d'indicatif après \esp{para que / a fin de que / con el fin de que}.
    \item [$\square$] J'utilise \esp{a fin de que} / \esp{con el fin de que} au moins une fois dans ma copie de thème/version pour varier le registre.
    \item [$\square$] Je distingue \esp{para + infinitif} (même sujet) et \esp{para que + subjonctif} (sujets différents) à 100\%.
\end{itemize}
\end{aretenir}

\coursec{Le futur dans la subordonnée : le subjonctif}

Après avoir étudié l'expression de la finalité (\emph{para, para que + subjonctif}), nous abordons maintenant un point crucial de la syntaxe espagnole : l'emploi du \textbf{subjonctif dans les subordonnées temporelles à valeur de futur}. En français, on utilise le futur simple (ou le futur antérieur) après « quand », « dès que », « jusqu'à ce que » ; l'espagnol, lui, bascule systématiquement au \textbf{subjonctif présent} (ou \textbf{subjonctif parfait} pour l'antériorité) dès que l'action n'est pas encore réalisée. C'est une différence structurelle majeure, source de nombreuses erreurs au Baccalauréat.

\courssub{Observation guidée : le déclencheur du subjonctif}

Lisez attentivement les phrases suivantes, extraites d'un dialogue entre deux étudiants camerounais, Paul et Marie, qui préparent leur avenir après le Baccalauréat. Repérez les connecteurs temporels et le mode du verbe dans la subordonnée.

\begin{textebox}[Dialogue : Projets d'avenir]
— \esp{¿Qué harás \textbf{cuando termines} el Bachillerato, Paul?}
— \esp{Me iré a Yaoundé \textbf{en cuanto salgan} las notas. \textbf{Hasta que encuentre} trabajo, me quedaré con mi tío.}
— \esp{Y yo, \textbf{tan pronto como tenga} el título, iré a Guinea Ecuatorial. \textbf{Siempre que pueda}, enviaré dinero a mi familia.}
— \esp{\textbf{Mientras estudies}, no te faltará nada. \textbf{Una vez que apruebes}, todo cambiará.}
\end{textebox}

\textbf{Glossaire :} \esp{termines} (tu finiras, subj. prés.), \esp{salgan} (elles sortiront, subj. prés.), \esp{encuentre} (je trouverai, subj. prés.), \esp{tenga} (j'aurai, subj. prés.), \esp{pueda} (je pourrai, subj. prés.), \esp{estudies} (tu étudieras, subj. prés.), \esp{apruebes} (tu réussiras, subj. prés.).

\emph{Observation :} Dans toutes ces phrases, l'action principale est au \textbf{futur indicatif} (\esp{harás, iré, me quedaré, iré, enviaré, cambiará, faltará}). Or, les verbes des subordonnées introduites par \esp{cuando, en cuanto, hasta que, tan pronto como, siempre que, mientras, una vez que} sont tous au \textbf{subjonctif présent} (\esp{termines, salgan, encuentre, tenga, pueda, estudies, apruebes}). Pourquoi ? Parce qu'elles expriment une action \textbf{future, non encore réalisée, incertaine ou hypothétique} au moment de l'énonciation.

\courssub{La règle fondamentale : connecteurs temporels + futur = subjonctif}

\begin{propriete}[Règle d'or : Subordonnée temporelle à valeur de futur]
Lorsque la proposition principale est au \textbf{futur indicatif} (ou à l'impératif), et que la subordonnée est introduite par un connecteur temporel marquant l'antériorité, la simultanéité ou la limite temporelle d'une action \textbf{future}, le verbe de la subordonnée se met au \textbf{subjonctif présent} (action simultanée ou ultérieure) ou au \textbf{subjonctif parfait} (\esp{haya + participe}, action antérieure accomplie).
\end{propriete}

\begin{center}
\begin{tabularx}{\linewidth}{|X|X|X|}
\hline
\rowcolor{popblueL} \textbf{Connecteur temporel} & \textbf{Sens / Valeur} & \textbf{Mode requis (si futur)} \\ \hline
\esp{cuando} & Quand, au moment où & Subjonctif présent / parfait \\ \hline
\esp{en cuanto} / \esp{tan pronto como} & Dès que, aussitôt que & Subjonctif présent / parfait \\ \hline
\esp{hasta que} & Jusqu'à ce que & Subjonctif présent / parfait \\ \hline
\esp{mientras} & Tant que, pendant que & Subjonctif présent \\ \hline
\esp{siempre que} / \esp{cuandoquiera que} & Chaque fois que, pourvu que & Subjonctif présent \\ \hline
\esp{una vez que} / \esp{después de que} & Une fois que, après que & Subjonctif présent / parfait \\ \hline
\esp{antes de que} & Avant que & \textbf{Toujours} Subjonctif (présent/parfait) \\ \hline
\end{tabularx}
\end{center}

\begin{exemplebox}[Exemples traduits : Futur principal $\rightarrow$ Subjonctif subordonné]
\begin{itemize}
    \item \esp{Cuando \textbf{termine} el examen, \textbf{saldré} de vacaciones.} \\
    « Quand j'\textbf{aurai fini} l'examen, je \textbf{partirai} en vacances. » (Antériorité : subj. présent \esp{termine} = futur antérieur FR).
    \item \esp{Te llamaré en cuanto \textbf{llegue} a Douala.} \\
    « Je t'appellerai dès que j'\textbf{aurai arrivé} / j'\textbf{arriverai} à Douala. »
    \item \esp{No me iré hasta que \textbf{vuelvas}.} \\
    « Je ne partirai pas avant que tu \textbf{reviennes} / jusqu'à ce que tu \textbf{reviennes}. »
    \item \esp{Una vez que \textbf{hayas terminado}, \textbf{podrás} descansar.} \\
    « Une fois que tu \textbf{auras fini}, tu \textbf{pourras} te reposer. » (Antériorité marquée : subj. parfait \esp{hayas terminado}).
\end{itemize}
\end{exemplebox}

\emph{Remarque sur la traduction :} Le français utilise le \textbf{futur antérieur} (\emph{j'aurai fini}) pour marquer l'antériorité sur le futur simple (\emph{je partirai}). L'espagnol utilise le \textbf{subjonctif présent} (\esp{termine}) pour la simultanéité ou l'antériorité simple, et le \textbf{subjonctif parfait} (\esp{haya terminado}) pour insister sur l'accomplissement total avant l'action principale.

\begin{popfigure}
\begin{tikzpicture}[
    scale=0.9,
    every node/.style={font=\small},
    axis/.style={->, >=Stealth, thick, popdark},
    event/.style={circle, draw=popblue, fill=popblueL, thick, minimum width=1.5cm, align=center},
    main/.style={rectangle, draw=poporange, fill=poporangeL, thick, rounded corners, minimum width=2.5cm, align=center},
    label/.style={text=popmuted, font=\tiny}
]
% Axe temporel
\draw[axis] (-0.5,0) -- (12,0) node[right, label] {Temps $\rightarrow$};

% Point de parole (Maintenant)
\draw[dashed, poppink, thick] (0, -0.5) -- (0, 2.5) node[above, poppink] {\textbf{MAINTENANT (Énonciation)}};

% Subordonnée (Action future, antérieure ou simultanée)
\node[event, align=center] (sub) at (3, 1.5){Action subordonnée\\ \textbf{(Subjonctif Présent)}\\\esp{cuando termine}};
\draw[popblue, thick, ->] (sub) -- (3, 0.2) node[below, label] {Non réalisée / Incertaine};

% Principale (Action future, postérieure)
\node[main, align=center] (main) at (8, 1.5){Action principale\\ \textbf{(Futur Indicatif)}\\\esp{saldré}};
\draw[poporange, thick, ->] (main) -- (8, 0.2) node[below, label] {Réalisée plus tard};

% Flèche de relation
\draw[popgreen, thick, <->, >=Stealth, dashed] (sub.east) -- node[above, popgreen, font=\tiny] {Antériorité / Simultanéité} (main.west);

% Subjonctif Parfait (Antériorité forte)
\node[event, draw=poppurple, fill=poppurpleL, align=center] (subperf) at (3, -1.5){Action subordonnée\\ \textbf{(Subjonctif Parfait)}\\\esp{una vez que haya terminado}};
\draw[poppurple, thick, ->] (subperf) -- (3, -0.2) node[below, label] {Accomplie avant la principale};
\draw[poppurple, thick, <->, >=Stealth, dashed] (subperf.east) -- node[below, poppurple, font=\tiny] {Antériorité forte} (main.west);
\end{tikzpicture}
\legende{Frise temporelle : L'action de la subordonnée (subjonctif) précède ou coïncide avec l'action future de la principale (futur indicatif). Le point « MAINTENANT » montre que les deux sont futures au moment de parler.}
\end{popfigure}

\courssub{Contraste Indicatif / Subjonctif : Le piège de l'habitude et du passé}

La clé pour ne pas se tromper est de se demander : \textbf{« L'action de la subordonnée est-elle déjà réalisée / habituelle, ou est-elle à venir ? »}

\begin{propriete}[Choix du mode : Réalité vs Projection]
\begin{itemize}
    \item \textbf{Indicatif (Présent, Passé composé, Imparfait, Futur) :} Action \textbf{réelle, accomplie} (passé) ou \textbf{habituelle} (présent/général). Le locuteur constate un fait.
    \item \textbf{Subjonctif (Présent, Parfait) :} Action \textbf{future, non réalisée, hypothétique} ou \textbf{non encore vécue}. Le locuteur projette une action à venir.
\end{itemize}
\end{propriete}

\begin{exemplebox}[Le même connecteur, deux modes, deux sens]
\begin{tabular}{@{}p{0.48\linewidth}p{0.48\linewidth}@{}}
\textbf{INDICATIF (Réalité / Habitude)} & \textbf{SUBJONCTIF (Futur / Projection)} \\
\esp{Cuando \textbf{era} niño, jugaba al fútbol.} & \esp{Cuando \textbf{sea} grande, seré médico.} \\
« Quand j'\textbf{étais} enfant, je jouais... » (Passé réel) & « Quand je \textbf{serai} grand, je serai médecin. » (Futur) \\
\esp{Cuando \textbf{llegó} ayer, cenamos.} & \esp{Cuando \textbf{llegue} mañana, cenaremos.} \\
« Quand il \textbf{est arrivé} hier, nous avons dîné. » (Passé accompli) & « Quand il \textbf{arrivera} demain, nous dînerons. » (Futur) \\
\esp{Siempre que \textbf{llueve}, no salgo.} & \esp{Siempre que \textbf{llueva} mañana, no saldré.} \\
« Chaque fois qu'il \textbf{pleut} (habitude), je ne sors pas. » & « S'il \textbf{pleut} demain (hypothèse), je ne sortirai pas. » \\
\esp{Mientras \textbf{estudiaba}, dormía.} & \esp{Mientras \textbf{estudies}, te espero.} \\
« Pendant que j'\textbf{étudiais} (passé), je dormais. » & « Tant que tu \textbf{étudieras} (futur), je t'attendrai. »
\end{tabular}
\end{exemplebox}

\begin{attention}[Piège n°1 du Bac : « Quand je serai grand »]
Ne traduisez \textbf{JAMAIS} une subordonnée temporelle future par l'indicatif futur.
\begin{itemize}
    \item \textbf{Faux :} *\esp{Cuando seré grande...} / *\esp{Cuando tendré dinero...}
    \item \textbf{Correct :} \esp{Cuando \textbf{sea} grande...} / \esp{Cuando \textbf{tenga} dinero...} (Subjonctif présent).
\end{itemize}
\emph{Astuce :} Remplacez mentalement « quand » par « si » ou « le jour où ». Si la phrase garde un sens futur, c'est du subjonctif. \esp{« Le jour où je serai grand » $\rightarrow$ \esp{Cuando sea grande}}.
\end{attention}

\begin{popfigure}
\begin{tikzpicture}[
    scale=0.85,
    every node/.style={font=\small},
    cell/.style={draw=popline, minimum height=0.9cm, align=center, inner sep=2pt},
    header/.style={cell, fill=popblue, text=white, font=\bfseries},
    ind/.style={cell, fill=popgreenL},
    subj/.style={cell, fill=poppinkL},
    conn/.style={cell, fill=popgoldL, font=\bfseries\itshape, text=popdark}
]
% Grille : Connecteur | Passé/Habitude (Indicatif) | Futur (Subjonctif)
\matrix (tab) [matrix of nodes, nodes in empty cells, column sep=-\pgflinewidth, row sep=-\pgflinewidth,
    column 1/.style={nodes={conn, minimum width=2.8cm}},
    column 2/.style={nodes={ind, minimum width=5.5cm}},
    column 3/.style={nodes={subj, minimum width=5.5cm}},
    row 1/.style={nodes={header, minimum height=1cm}}
]{
    \textbf{Connecteur} & \textbf{Contexte : Passé / Habitude / Fait accompli $\rightarrow$ INDICATIF} & \textbf{Contexte : Futur / Hypothèse / Non-réalisé $\rightarrow$ SUBJONCTIF} \\
    \esp{cuando} &
    \esp{Cuando \textbf{llovió}, no salimos.} (Passé)\newline
    \esp{Cuando \textbf{llueve}, me mojo.} (Habitude Présent) &
    \esp{Cuando \textbf{llueva} mañana, no saldré.} (Futur) \\
    \esp{en cuanto / tan pronto como} &
    \esp{En cuanto \textbf{llegó}, comimos.} (Passé) &
    \esp{En cuanto \textbf{llegue}, comeremos.} (Futur) \\
    \esp{hasta que} &
    \esp{Esperé hasta que \textbf{llegó}.} (Passé accompli) &
    \esp{Esperaré hasta que \textbf{llegue}.} (Futur) \\
    \esp{mientras} &
    \esp{Mientras \textbf{leía}, dormía.} (Passé simultané) &
    \esp{Mientras \textbf{leas}, yo espero.} (Futur simultané) \\
    \esp{siempre que} &
    \esp{Siempre que \textbf{podía}, iba.} (Habitude Passé) &
    \esp{Siempre que \textbf{pueda}, iré.} (Futur conditionnel) \\
    \esp{una vez que / después de que} &
    \esp{Una vez que \textbf{terminó}, se fue.} (Passé antérieur) &
    \esp{Una vez que \textbf{termine / haya terminado}, se irá.} (Futur antérieur) \\
    \esp{antes de que} &
    \multicolumn{2}{c|}{\textbf{TOUJOURS SUBJONCTIF} (\esp{antes de que \textbf{llegue / llegara / haya llegado}})} \\
};
% Légendes couleur
\node[below=0.3cm of tab.south west, anchor=west, font=\tiny, text=popgreen] {■ Indicatif : Réalité vécue ou habitude};
\node[below=0.7cm of tab.south west, anchor=west, font=\tiny, text=poppink] {■ Subjonctif : Projection future, incertitude};
\end{tikzpicture}
\legende{Tableau de décision : Le même connecteur change de mode selon que l'action est perçue comme acquise (indicatif) ou à venir (subjonctif). \esp{Antes de que} est l'exception qui impose toujours le subjonctif.}
\end{popfigure}

\courssub{Le cas particulier de « Si » : Conditionnel et Subjonctif imparfait}

Le connecteur \esp{si} (si) suit une logique proche mais distincte : il introduit une \textbf{condition}. La règle d'or est l'interdiction absolue du conditionnel ou du futur après \esp{si}.

\begin{propriete}[Les trois types de phrases conditionnelles avec \esp{si}]
\begin{center}
\begin{tabularx}{\linewidth}{|X|X|X|X|}
\hline
\rowcolor{popblueL} \textbf{Type} & \textbf{Si + Subordonnée} & \textbf{Principale} & \textbf{Sens / Exemple} \\ \hline
\textbf{1. Réelle / Possible} & \textbf{Indicatif Présent} & \textbf{Futur Indicatif} (ou Présent/Imperatif) & \esp{Si \textbf{estudias}, \textbf{aprobarás}.} \\
& \esp{Si estudias...} & \esp{aprobarás} & « Si tu étudies, tu réussiras. » \\ \hline
\textbf{2. Irréelle Présent/Futur} & \textbf{Subjonctif Imparfait} & \textbf{Conditionnel Présent} & \esp{Si \textbf{estudiara(s)}, \textbf{aprobaría(s)}.} \\
& \esp{Si estudiara(s)...} & \esp{aprobaría(s)} & « Si tu étudiais, tu réussirais. » \\ \hline
\textbf{3. Irréelle Passé} & \textbf{Subjonctif Plus-que-parfait} & \textbf{Conditionnel Passé} & \esp{Si \textbf{hubieras estudiado}, \textbf{habrías aprobado}.} \\
& \esp{Si hubieras estudiado...} & \esp{habrías aprobado} & « Si tu avais étudié, tu aurais réussi. » \\ \hline
\end{tabularx}
\end{center}
\end{propriete}

\begin{attention}[Piège n°2 : *« Si tendría » / *« Si estudiaré »]
\begin{itemize}
    \item \textbf{Jamais} de conditionnel après \esp{si} : *\esp{Si tendría dinero...} $\rightarrow$ \esp{Si \textbf{tuviera} dinero...}
    \item \textbf{Jamais} de futur après \esp{si} : *\esp{Si estudiaré...} $\rightarrow$ \esp{Si \textbf{estudio}...}
    \item \textbf{Jamais} d'indicatif imparfait après \esp{si} pour l'irréel : *\esp{Si estudiaba...} $\rightarrow$ \esp{Si \textbf{estudiara}...}
\end{itemize}
\emph{Rappel forme :} Subjonctif imparfait = 3ème pers. plur. P.S. (\esp{estudiaron}) $\rightarrow$ enlever \esp{-ron} $\rightarrow$ \esp{estudiara, estudiases, estudiara, estudiáramos, estudiaseis, estudiaran}. Deux séries possibles (\esp{-ra} / \esp{-se}), la série en \esp{-ra} est la plus courante.
\end{attention}

\courssub{Subjonctif imparfait et plus-que-parfait : antériorité dans le passé}

Lorsque la proposition principale est au \textbf{conditionnel} (ou à un temps du passé : imparfait, passé composé, plus-que-parfait), la subordonnée temporelle (ou conditionnelle) bascule au \textbf{subjonctif imparfait} (simultanéité/futur dans le passé) ou \textbf{subjonctif plus-que-parfait} (antériorité dans le passé). C'est la \textbf{concordance des temps} (consecutio temporum).

\begin{exemplebox}[Concordance des temps : Principale au Conditionnel / Passé]
\begin{itemize}
    \item \esp{Te \textbf{llamaría} en cuanto \textbf{supiera} la noticia.} \\
    « Je t'\textbf{appellerais} dès que je \textbf{saurais} la nouvelle. » (Futur dans le passé $\rightarrow$ Subj. Imparfait).
    \item \esp{Esperé hasta que \textbf{llegara}.} \\
    « J'\textbf{ai attendu} jusqu'à ce qu'il \textbf{arrive}. » (Passé principal $\rightarrow$ Subj. Imparfait).
    \item \esp{Me \textbf{habría ido} antes de que \textbf{llegaran}.} \\
    « Je \textbf{serais parti} avant qu'ils n'\textbf{arrivent}. » (Conditionnel passé $\rightarrow$ Subj. Plus-que-parfait \esp{llegaran}).
    \item \esp{No salí hasta que \textbf{hubieran terminado}.} \\
    « Je ne suis pas sorti avant qu'ils n'\textbf{aient fini}. » (Passé principal $\rightarrow$ Subj. Plus-que-parfait).
\end{itemize}
\end{exemplebox}

\courssub{Méthode : Choisir le bon mode en 3 étapes}

\begin{methode}[Comment choisir Indicatif ou Subjonctif après un connecteur temporel ?]
\begin{enumerate}
    \item \textbf{Identifiez le temps de la principale.} Est-elle au Futur / Impératif / Présent de projection ? $\rightarrow$ Risque de Subjonctif. Est-elle au Passé / Conditionnel ? $\rightarrow$ Subjonctif Imparfait / Plus-que-parfait. Est-elle au Présent / Passé (fait réel) ? $\rightarrow$ Indicatif probable.
    \item \textbf{Analysez la valeur de la subordonnée.} L'action est-elle \textbf{déjà accomplie} (passé) ou \textbf{habituelle} (généralité) ? $\rightarrow$ \textbf{Indicatif}. L'action est-elle \textbf{future, incertaine, non réalisée, hypothétique} ? $\rightarrow$ \textbf{Subjonctif}.
    \item \textbf{Vérifiez le connecteur.} \esp{Antes de que} $\rightarrow$ \textbf{Toujours Subjonctif}. \esp{Si} $\rightarrow$ Voir tableau des 3 types (Jamais indicatif futur/conditionnel après \esp{si}). Les autres (\esp{cuando, hasta que, en cuanto...}) $\rightarrow$ Appliquent la règle de l'étape 2.
\end{enumerate}
\end{methode}

\begin{savaistu}[Una vez que : le connecteur « académique »]
\esp{Una vez que} (une fois que) est très fréquent dans les textes formels, administratifs ou journalistiques (ex: communiqués du MINESUP, contrats de travail en Guinée équatoriale). Il suit \textbf{exactement} la même règle que \esp{cuando / después de que} :
\begin{itemize}
    \item Futur principal : \esp{Una vez que \textbf{firmes} el contrato, \textbf{empezarás} a trabajar.} (Subj. Présent)
    \item Passé principal : \esp{Una vez que \textbf{hubieran firmado}, \textbf{empezaron}.} (Subj. Plus-que-parfait)
\end{itemize}
\emph{Ne pas confondre avec :} \esp{Una vez} (adverbe) + infinitif/participe : \esp{Una vez \textbf{terminado} el trabajo, nos fuimos.} (Participe passé accordé / valeur passive).
\end{savaistu}

\courssub{Entraînement : Thème et Version guidés}

\begin{exoresolu}[Traduction : Thème (Français $\rightarrow$ Espagnol)]
\textbf{Énoncé :} Traduisez en espagnol les phrases suivantes en justifiant le mode choisi.
\begin{enumerate}
    \item Quand tu auras fini tes devoirs, nous irons au stade.
    \item J'attendrai ici jusqu'à ce qu'il revienne.
    \item Dès que j'aurai mon salaire, je t'achèterai un cadeau.
    \item Si j'avais le temps, je voyagerais en Espagne.
    \item Quand j'étais petit, je jouais dans la rue. (Habitude passée)
\end{enumerate}
\tcblower
\textbf{Solution détaillée :}
\begin{enumerate}
    \item \esp{\textbf{Cuando termines} tus deberes, iremos al estadio.}
    \emph{Justification :} Principale au futur (\esp{iremos}) + \esp{cuando} + action future non réalisée $\rightarrow$ \textbf{Subjonctif présent} (\esp{termines} de \esp{terminar}).
    \item \esp{Esperaré aquí \textbf{hasta que vuelva}.}
    \emph{Justification :} Futur principal (\esp{esperaré}) + \esp{hasta que} + limite future $\rightarrow$ \textbf{Subjonctif présent} (\esp{vuelva} de \esp{volver}, irrégulier \emph{o$\rightarrow$ue}).
    \item \esp{\textbf{En cuanto tenga} mi salario, \textbf{te compraré} un regalo.} (ou \esp{tan pronto como tenga} / \esp{una vez que tenga}).
    \emph{Justification :} Futur principal + antériorité immédiate $\rightarrow$ \textbf{Subjonctif présent} (\esp{tenga} de \esp{tener}, irrégulier \emph{g} + \emph{e$\rightarrow$ie}). Note : \esp{haya tenido} (subj. parfait) possible si insistance sur l'accomplissement.
    \item \esp{\textbf{Si tuviera} tiempo, \textbf{viajaría} a España.}
    \emph{Justification :} Conditionnel irréel présent (Type 2) $\rightarrow$ \esp{Si} + \textbf{Subjonctif imparfait} (\esp{tuviera} de \esp{tener}) + \textbf{Conditionnel} (\esp{viajaría}). \textbf{Interdit :} *\esp{Si tengo / Si tendría / Si tendría}.
    \item \esp{\textbf{Cuando era} pequeño, \textbf{jugaba} en la calle.}
    \emph{Justification :} Habitude passée (réalité vécue) $\rightarrow$ \textbf{Indicatif Imparfait} (\esp{era, jugaba}). Pas de subjonctif car l'action est accomplie/réelle.
\end{enumerate}
\end{exoresolu}

\begin{exoresolu}[Traduction : Version (Espagnol $\rightarrow$ Français) \& Analyse]
\textbf{Énoncé :} Traduisez et analysez le mode des verbes soulignés.
\esp{« Cuando \textbf{termine} la reunión, \textbf{saldremos} inmediatamente. \textbf{Hasta que no \textbf{tengamos}} la autorización, no \textbf{podremos} empezar. \textbf{Siempre que \textbf{haya}} plazas libres, \textbf{os inscribiréis}. »}
\tcblower
\textbf{Traduction :}
« Quand la réunion \textbf{sera finie} (ou : \textbf{finira}), nous \textbf{sortirons} immédiatement. \textbf{Jusqu'à ce que nous \textbf{ayons}} l'autorisation, nous ne \textbf{pourrons} pas commencer. \textbf{Chaque fois qu'il \textbf{y aura}} des places libres, vous \textbf{vous inscrirez}. »

\textbf{Analyse grammaticale :}
\begin{itemize}
    \item \esp{\textbf{termine}} (subj. prés. de \esp{terminar}) : Subordonnée \esp{cuando} + principale futur (\esp{saldremos}) $\rightarrow$ Futur antérieur FR / Subj. Prés. ES.
    \item \esp{\textbf{tengamos}} (subj. prés. de \esp{tener}, irrég. \emph{g} + \emph{e$\rightarrow$ie}) : \esp{Hasta que no} (jusqu'à ce que... non) + principale futur (\esp{podremos}) $\rightarrow$ Subj. Prés. \emph{Note :} \esp{hasta que no} renforce l'attente de la réalisation.
    \item \esp{\textbf{haya}} (subj. prés. de \esp{haver}, irrégulier) : \esp{Siempre que} + principale futur (\esp{os inscribiréis}) $\rightarrow$ Valeur conditionnelle future (« s'il y a » / « chaque fois qu'il y aura ») $\rightarrow$ Subj. Prés.
\end{itemize}
\end{exoresolu}

\begin{exercice}[Entraînement personnel : Thème / Version]
\textbf{A. Thème :} Mettez au bon temps (Indicatif ou Subjonctif, Présent/Imparfait/Parfait) les verbes entre parenthèses.
\begin{enumerate}
    \item Te avisaré en cuanto yo \_\_\_\_\_\_\_\_\_\_ (saber) la verdad.
    \item \_\_\_\_\_\_\_\_\_\_ (Cuando - ser) yo mayor, seré piloto.
    \item No compres el billete hasta que yo \_\_\_\_\_\_\_\_\_\_ (decir) que sí.
    \item Si yo \_\_\_\_\_\_\_\_\_\_ (tener) coche, iría a la playa.
    \item Cuando nosotros \_\_\_\_\_\_\_\_\_\_ (llegar) ayer, ya habían cenado.
    \item Esperaremos mientras tú \_\_\_\_\_\_\_\_\_\_ (hacer) la cola.
    \item Una vez que ustedes \_\_\_\_\_\_\_\_\_\_ (terminar / parfait), podréis salir.
    \item Si hubiéramos sabido, \_\_\_\_\_\_\_\_\_\_ (nosotros - venir) antes.
\end{enumerate}

\textbf{B. Version :} Traduisez en français.
\esp{« Cuando llegues a Malabo, llámame. Si no tienes crédito, envíame un mensaje. Hasta que no reciba tu noticia, no me iré. Si pudiera, iría yo mismo a buscarte. »}

\textbf{C. Piège à éviter :} Corrigez les erreurs suivantes.
\begin{enumerate}
    \item *Cuando \textbf{iré} a Yaoundé, te veré.
    \item *Si \textbf{lloverá} mañana, no voy.
    \item *Hasta que \textbf{llegará}, espero.
    \item *Antes de que \textbf{llega}, avísame.
\end{enumerate}
\end{exercice}

\begin{corrige}[Exercice Entraînement personnel]
\textbf{A. Thème :}
\begin{enumerate}
    \item \esp{sepa} (Subj. Prés. \esp{saber} irrégulier : \emph{sep-}) — Futur principal + \esp{en cuanto}.
    \item \esp{Cuando sea} mayor... (Subj. Prés. \esp{ser} irrégulier) — Futur principal (\esp{seré}) + \esp{cuando} futur.
    \item \esp{diga} (Subj. Prés. \esp{decir} irrégulier : \emph{dig-}) — Futur principal (\esp{compres} impératif/futur) + \esp{hasta que}.
    \item \esp{tuviera / tuviese} (Subj. Imparfait \esp{tener}) — Type 2 \esp{Si} + Conditionnel (\esp{iría}).
    \item \esp{llegamos} (Indicatif \textbf{Prétérit Indéfini} / Passé Simple) — Passé réel accompli (\esp{cuando} + passé ponctuel \emph{ayer}). \emph{Variante :} \esp{llegábamos} (Imparfait) si aspect duratif/habituel.
    \item \esp{hagas} (Subj. Prés. \esp{hacer} irrégulier : \emph{hag-}) — Futur principal (\esp{esperaremos}) + \esp{mientras} futur.
    \item \esp{hayan terminado} (Subj. Parfait \esp{terminar}) — Futur principal (\esp{podréis}) + \esp{una vez que} + antériorité marquée.
    \item \esp{habríamos venido} (Conditionnel Passé) — Type 3 \esp{Si} + Subj. Plus-que-parfait (\esp{hubiéramos sabido}).
\end{enumerate}

\textbf{B. Version :}
« Quand tu \textbf{arriveras} à Malabo, appelle-moi. Si tu n'\textbf{as} pas de crédit, envoie-moi un message. \textbf{Jusqu'à ce que je \textbf{reçoive}} ta nouvelle, je ne partirai pas. Si je \textbf{pouvais}, j'irai moi-même te chercher. »
\emph{Analyse :} \esp{llegues} (Subj. Prés.), \esp{tengas} (Indicatif Présent : \esp{Si no tienes} = Type 1, condition réelle possible), \esp{reciba} (Subj. Prés. \esp{recibir}) après \esp{hasta que no} + Futur, \esp{pudiera} (Subj. Imparfait) + Conditionnel \esp{iría} (sous-entendu).

\textbf{C. Corrections :}
\begin{enumerate}
    \item \esp{Cuando \textbf{vaya} a Yaoundé...} (Subj. Prés. \esp{ir}).
    \item \esp{Si \textbf{llueve} mañana...} (Indicatif Présent Type 1) ou \esp{Si \textbf{lloviera} mañana...} (Subj. Imparfait Type 2).
    \item \esp{Hasta que \textbf{llegue}, espero.} (Subj. Prés.).
    \item \esp{Antes de que \textbf{llegue}, avísame.} (Subj. Prés. \textbf{obligatoire} après \esp{antes de que}).
\end{enumerate}
\end{corrige}

\coursec{Révision de l'emphase et structures causales / finales}

\courssub{Transition et observation guidée}

Après avoir étudié l'emploi du subjonctif dans les subordonnées temporelles et conditionnelles exprimant le futur (\emph{cuando, en cuanto, si} + subjonctif), nous abordons maintenant un point transversal essentiel pour la traduction et l'expression écrite : \textbf{l'emphase}. En espagnol, comme en français, il ne suffit pas d'énoncer une cause (\emph{porque, ya que}) ou une finalité (\emph{para que} + subjonctif) ; il faut souvent \textbf{mettre en relief} un élément du message pour insister sur la responsabilité, la raison unique ou l'objectif principal. C'est le rôle des structures emphatiques que nous allons réviser systématiquement avant de les combiner avec les notions de cause et de finalité vues précédemment.

Observons d'abord un court texte original qui concentre plusieurs de ces structures dans un contexte professionnel camerounais.

\begin{textebox}[Témoignage : l'entrepreneuriat jeune à Douala]
\textbf{El reto de emprender en Douala}

Lo fundamental no es tener capital, \textbf{lo que cuenta} es la formación. \textbf{Es la disciplina lo que permite} sostener un negocio cuando \textbf{es porque} el mercado fluctúa \textbf{que} muchos abandonan. A los jóvenes \textbf{les digo} siempre: \textbf{no es el diploma lo que garantiza} el éxito, \textbf{sino la constancia}. \textbf{Para que} un proyecto prospere, \textbf{es necesario que} el emprendedor \textbf{se forme} continuamente. \textbf{¡Eso sí que es} la clave!
\end{textebox}

\begin{center}
\begin{tabularx}{\linewidth}{|X|X|}
\hline
\textbf{Espagnol (extrait)} & \textbf{Français} \\
\hline
\emph{Lo fundamental no es tener capital} & Ce qui est fondamental, ce n'est pas d'avoir du capital \\
\emph{lo que cuenta} & ce qui compte / c'est ce qui compte \\
\emph{Es la disciplina lo que permite} & C'est la discipline qui permet \\
\emph{es porque el mercado fluctúa que} & c'est parce que le marché fluctue que \\
\emph{A los jóvenes les digo} & Aux jeunes, je leur dis / Je dis aux jeunes \\
\emph{no es el diploma lo que garantiza} & ce n'est pas le diplôme qui garantit \\
\emph{¡Eso sí que es la clave!} & Ça, c'est vraiment la clé ! \\
\hline
\end{tabularx}
\end{center}

\courssub{Rappel : \texorpdfstring{\cle{lo + adjectif}}{lo + adjectif} — « Ce qui est [adj]… »}

\begin{propriete}[Structure \emph{lo + adjectif (singulier masculin)}]
\emph{Lo} + adjectif (au masculin singulier) + verbe (souvent \emph{ser}) agit comme un \textbf{nom neutre} signifiant « ce qui est [adjectif] », « la chose [adjectif] », « l'aspect [adjectif] ». L'adjectif reste \textbf{invariable} (masculin singulier) quel que soit le sujet réel de la phrase.
\end{propriete}

\begin{exemplebox}[Exemples variés]
\esp{Lo importante es la salud.} — « Ce qui est important, c'est la santé. » \\
\esp{Lo difícil será convencer al director.} — « Ce qui sera difficile, ce sera de convaincre le directeur. » \\
\esp{Lo mejor de todo es que ya tenemos el visado.} — « Le meilleur de tout, c'est que nous avons déjà le visa. » \\
\esp{Lo triste es que nadie escucha.} — « Ce qui est triste, c'est que personne n'écoute. »
\end{exemplebox}

\begin{attention}[Piège francophone : accord de l'adjectif]
En français, « ce qui est \emph{importante} » n'existe pas ; on dit « ce qui est \emph{important} » (masculin singulier neutre). En espagnol, \textbf{jamais} *\emph{la importante} ou *\emph{lo importada}. L'adjectif après \emph{lo} est \textbf{toujours masculin singulier}.
\begin{itemize}
    \item \checkmark \esp{Lo necesario es estudiar.}
    \item \cross \esp{La necesaria es estudiar.} / \esp{Lo necesarias son estudiar.}
\end{itemize}
\end{attention}

\courssub{Rappel : \texorpdfstring{\cle{lo que + verbe}}{lo que + verbe} — « Ce que / Ce qui… »}

\begin{propriete}[Structure \emph{lo que} + verbe conjugué]
\emph{Lo que} (pronom neutre « ce que / ce qui ») introduit une \textbf{proposition subordonnée nominale} qui fonctionne comme sujet ou COD. Le verbe qui suit \emph{lo que} s'accorde avec son propre sujet (explicite ou implicite).
\end{propriete}

\begin{exemplebox}[Sujet et COD]
\esp{Lo que \textbf{me preocupa} es el paro.} (Sujet : \emph{lo que me preocupa}) — « Ce qui m'inquiète, c'est le chômage. » \\
\esp{No entiendo \textbf{lo que \textbf{dices}}.} (COD : \emph{lo que dices}) — « Je ne comprends pas ce que tu dis. » \\
\esp{\textbf{Lo que \textbf{hicimos}} ayer fue un error.} — « Ce que nous avons fait hier était une erreur. »
\end{exemplebox}

\begin{aretenir}[Différence \emph{lo + adj} vs \emph{lo que + verbe}]
\begin{itemize}
    \item \emph{Lo + adj} = \textbf{qualité/valeur} abstraite (statique). \esp{Lo urgente} = l'urgence.
    \item \emph{Lo que + verbe} = \textbf{action/fait} concret (dynamique). \esp{Lo que hago} = ce que je fais.
\end{itemize}
\end{aretenir}

\courssub{Mise en relief du nom : \texorpdfstring{\cle{es X lo que / quien}}{es X lo que / quien} — « C'est X qui/que… »}

C'est la structure majeure pour traduire l'emphase française \emph{c'est... qui/que} (cleft sentence) sur un \textbf{nom} ou un \textbf{pronom tonique}.

\begin{propriete}[Structure \emph{Ser + élément mis en relief + lo que / quien + verbe}]
Pour mettre l'accent sur un \textbf{nom} (sujet, COD, COI, etc.), on utilise : \textbf{Es/Éramos/Fue/Será} + \textbf{[Élément]} + \textbf{lo que / quien} + \textbf{verbe}.
\begin{itemize}
    \item \emph{Lo que} : si l'élément mis en relief est une \textbf{chose} ou une \textbf{idée} (antécédent inanimé ou clause).
    \item \emph{Quien / Quienes} : si l'élément mis en relief est une \textbf{personne} (antécédent animé).
\end{itemize}
\end{propriete}

\begin{exemplebox}[Cas principaux]
\textbf{Sujet mis en relief (chose) :} \esp{Es el esfuerzo \textbf{lo que} da resultados.} — « C'est l'effort qui donne des résultats. » \\
\textbf{Sujet mis en relief (personne) :} \esp{Fue \textbf{mi hermano quien} me ayudó.} — « C'est mon frère qui m'a aidé. » \\
\textbf{COD mis en relief :} \esp{Es la paciencia \textbf{lo que} necesitas.} — « C'est la patience qu'il te faut / dont tu as besoin. » \\
\textbf{COI / Complément prépositionnel :} \esp{Es \textbf{para ti} \textbf{lo que} trabajo.} — « C'est pour toi que je travaille. »
\end{exemplebox}

\begin{attention}[Erreur majeure : *\emph{Es... que} au lieu de *\emph{Es... lo que}]
C'est \textbf{la} faute la plus fréquente des francophones par calque sur « C'est... que ».
\begin{itemize}
    \item \cross \esp{*Es el esfuerzo \textbf{que} cuenta.} (Incorrect : \emph{que} seul ne peut pas reprendre un nom antécédent comme pronom relatif neutre ici).
    \item \checkmark \esp{Es el esfuerzo \textbf{lo que} cuenta.}
    \item \cross \esp{*Es Juan \textbf{que} vino.}
    \item \checkmark \esp{Es Juan \textbf{quien} vino.} (Personne $\rightarrow$ \emph{quien}).
\end{itemize}
\emph{Nota} : \emph{Es... que} n'existe que pour l'emphase sur \textbf{adverbe} ou \textbf{complément circonstanciel} de temps/lieu : \esp{Es \textbf{mañana que} vienen} (C'est demain qu'ils viennent).
\end{attention}

\courssub{Dédoublement pronominal (Topicalisation) : \texorpdfstring{\cle{A + nom + pronom}}{A + nom + pronom}}

Très courant à l'oral et dans l'écrit journalistique pour \textbf{thématiser} (mettre en thème de départ) un complément d'objet (souvent animé) ou un complément circonstanciel.

\begin{propriete}[Redoublement du COD/COI par pronom atone]
Structure : \textbf{Preposition \emph{a} + Nom/Pronom tonique} + \textbf{Pronom atone (le, la, lo, les, les, se, me, te, nos, os)} + \textbf{Verbe}.
Permet de mettre l'accent sur \textbf{qui} reçoit l'action ou \textbf{à qui} elle s'adresse.
\end{propriete}

\begin{exemplebox}[Exemples contextuels]
\esp{\textbf{A Juan} \textbf{lo} vi ayer en el mercado.} — « Jean, je l'ai vu hier au marché. » (COD animé topicalisé) \\
\esp{\textbf{A los estudiantes} \textbf{les} explicaremos la beca.} — « Aux étudiants, nous leur expliquerons la bourse. » (COI topicalisé) \\
\esp{\textbf{A mí} \textbf{me} parece bien.} — « Moi, cela me semble bien. » (Insistance sur la personne) \\
\esp{\textbf{Ese libro} \textbf{lo} he leído tres veces.} — « Ce livre, je l'ai lu trois fois. » (COD inanimé sans \emph{a}, redoublé par \emph{lo})
\end{exemplebox}

\begin{savaistu}[Le \emph{leísmo}, \emph{loísmo}, \emph{laísmo} au Cameroun et en Guinée équatoriale]
En Espagne (péninsulaire), on distingue \emph{le} (COI) et \emph{lo/la} (COD). En Amérique latine et en Guinée équatoriale (donc contexte camerounais proche), l'usage standard tend vers le \textbf{\emph{leísmo}} pour les personnes (COD personne = \emph{le}) : \esp{A Juan \textbf{le} vi}. Cependant, la norme académique internationale (RAE) pour l'écrit formel (examens, administration) recommande \textbf{\emph{lo/la}} pour le COD : \esp{A Juan \textbf{lo} vi}. \textbf{Conseil :} Pour le Baccalauréat, utilisez \emph{lo/la} (COD) et \emph{le} (COI).
\end{savaistu}

\courssub{Le futur dans la subordonnée temporelle et conditionnelle}

\begin{propriete}[Subjonctif présent après conjonctions de temps/condition à valeur future]
Lorsque l'action de la subordonnée est \textbf{postérieure} à l'énonciation (futur), et non réalisée au moment de la parole, on utilise le \textbf{subjonctif présent} après :
\begin{itemize}
    \item \textbf{Conjonctions de temps :} \emph{cuando, en cuanto, tan pronto como, hasta que, mientras, después de que, antes de que}.
    \item \textbf{Conjonction de condition :} \emph{si} (si + subjonctif = hypothèse future peu probable ou non réalisée ; \emph{si} + indicatif = condition réelle ou hypothèse ouverte).
\end{itemize}
La principale est au futur simple, à l'impératif, ou au présent de valeur future.
\end{propriete}

\begin{exemplebox}[Temporelles et conditionnelles]
\esp{\textbf{Cuando termines} (subj.), te llamo.} — « Quand tu auras fini, je t'appellerai. » (Futur antérieur $\rightarrow$ Subj.) \\
\esp{\textbf{En cuanto llegue} (subj.), empezamos.} — « Dès qu'il arrivera, nous commencerons. » \\
\esp{\textbf{No saldré} \textbf{hasta que} \textbf{vengas} (subj.).} — « Je ne sortirai pas avant que tu viennes. » \\
\esp{\textbf{Si tengo} (ind.) \textbf{tiempo, iré.}} — « Si j'ai le temps, j'irai. » (Condition réelle/ouverte $\rightarrow$ Indicatif). \\
\esp{\textbf{Si tuviera} (subj. imparfait) \textbf{tiempo, iría.}} — « Si j'avais le temps, j'irais. » (Irréel présent $\rightarrow$ Subj. imparfait).
\end{exemplebox}

\begin{attention}[Indicatif vs Subjonctif : Habitude vs Futur]
Si l'action est \textbf{habituelle} ou \textbf{passée/accomplie}, on met l'\textbf{indicatif}.
\esp{\textbf{Cuando termino} (ind.) \textbf{trabajo, voy al gimnasio.}} — « Quand je finis le travail (habitude), je vais à la salle. » \\
\esp{\textbf{Cuando terminé} (ind.) \textbf{el trabajo, fui al gimnasio.}} — « Quand j'ai fini le travail (passé), je suis allé... »
\end{attention}

\courssub{Inversion sujet-verbe et adverbes d'emphase}

\begin{propriete}[Inversion et particules \emph{sí}, \emph{ya}, \emph{bien}]
\begin{enumerate}
    \item \textbf{Inversion VS} : Place le verbe avant le sujet pour insister sur l'action ou le verbe lui-même. \esp{\textbf{Vino} Juan.} (C'est Juan qui est venu / Juan EST venu).
    \item \textbf{\emph{Sí} + verbe} : Affirmation forte, contredit une négation ou un doute. \esp{\textbf{Sí} estudio.} (J'étudie \emph{bien sûr} / Si, j'étudie).
    \item \textbf{\emph{Ya} + indicatif} : « Déjà / Maintenant / Enfin » avec insistance sur l'accomplissement. \esp{\textbf{Ya} terminó.} (Il a \emph{enfin} fini).
    \item \textbf{\emph{Qué} + nom / \emph{Cómo} + verbe} : Exclamatif emphatique. \esp{\textbf{Qué} paciencia tienes!} \esp{\textbf{Cómo} has crecido!}
\end{enumerate}
\end{propriete}

\begin{exemplebox}[Le \emph{« Eso sí que... »} du texte d'observation]
\esp{\textbf{¡Eso sí que es} la clave!} — « Ça, c'est \emph{vraiment} la clé ! » (\emph{sí} renforce l'affirmation, \emph{eso} est topicalisé).
\esp{\textbf{¡Eso no que no!}} — « Ça, non, certainement pas ! » (Refus catégorique).
\end{exemplebox}

\courssub{Emphase combinée avec la cause et la finalité (Cœur de la leçon E21)}

C'est ici que la révision prend tout son sens pour la traduction : il faut savoir emphatiser \textbf{la cause} (\emph{porque}) ou \textbf{la finalité} (\emph{para que}).

\begin{propriete}[Emphase sur la cause : \emph{Es porque... que} / \emph{Es por... que}]
Pour traduire « C'est \textbf{parce que}... \textbf{que}... » ou « C'est \textbf{à cause de}... \textbf{que}... » :
\begin{itemize}
    \item \textbf{Es porque} + indicatif + \textbf{que} + indicatif. \esp{\textbf{Es porque estudia} \textbf{que} aprueba.} = « C'est \emph{parce qu'il étudie} qu'il réussit. »
    \item \textbf{Es por} + nom/infinitif + \textbf{que} + indicatif. \esp{\textbf{Es por su esfuerzo} \textbf{que} ganó.} = « C'est \emph{grâce à son effort} qu'il a gagné. »
\end{itemize}
\end{propriete}

\begin{propriete}[Emphase sur la finalité : \emph{Es para... que} + subjonctif / \emph{Es con el fin de que} + subjonctif]
Pour traduire « C'est \textbf{pour que}... / \textbf{afin que}... » :
\begin{itemize}
    \item \textbf{Es para que} + \textbf{subjonctif}. \esp{\textbf{Es para que aprendas} \textbf{que} te corrijo.} = « C'est \emph{pour que tu apprennes} que je te corrige. »
    \item \textbf{Es con el fin de que} + \textbf{subjonctif} (plus formel). \esp{\textbf{Es con el fin de que todos participen} \textbf{que} organizamos la reunion.} = « C'est \emph{afin que tous participent} que nous organisons la réunion. »
\end{itemize}
\end{propriete}

\begin{exemplebox}[Synthèse : Cause vs Finalité en emphase]
\begin{center}
\begin{tabularx}{\linewidth}{|X|X|X|}
\hline
\rowcolor{popblueL} \textbf{Français (Emphase)} & \textbf{Espagnol (Structure)} & \textbf{Mode du verbe} \\
\hline
C'est \textbf{parce qu'il travaille} qu'il réussit. & \esp{Es \textbf{porque trabaja} \textbf{que} aprueba.} & Indicatif (cause réelle) \\
\hline
C'est \textbf{à cause de la pluie} qu'on annule. & \esp{Es \textbf{por la lluvia} \textbf{que} cancelamos.} & Indicatif \\
\hline
C'est \textbf{pour qu'il réussisse} qu'on l'aide. & \esp{Es \textbf{para que apruebe} \textbf{que} lo ayudamos.} & \textbf{Subjonctif} (finalité) \\
\hline
C'est \textbf{afin que tout le monde sache} qu'on publie. & \esp{Es \textbf{con el fin de que todos sepan} \textbf{que} publicamos.} & \textbf{Subjonctif} \\
\hline
\end{tabularx}
\end{center}
\end{exemplebox}

\courssub{Méthode : Traduire une phrase emphatique français $\rightarrow$ espagnol (Thème)}

\begin{methode}[Étapes pour le thème emphatique]
\begin{enumerate}
    \item \textbf{Identifier l'élément mis en relief} (ce qui suit « C'est... qui/que » ou « Ce qui/ce que »).
    \item \textbf{Classifier cet élément} :
        \begin{itemize}
            \item Adjectif / Qualité abstraite $\rightarrow$ \cle{Lo + adjectif}.
            \item Action / Fait / Proposition entière $\rightarrow$ \cle{Lo que + verbe}.
            \item Nom (chose) $\rightarrow$ \cle{Es + Nom + lo que}.
            \item Nom (personne) $\rightarrow$ \cle{Es + Nom + quien}.
            \item Cause (\emph{parce que}) $\rightarrow$ \cle{Es porque... que}.
            \item Finalité (\emph{pour que}) $\rightarrow$ \cle{Es para que... que} + \textbf{Subjonctif}.
        \end{itemize}
    \item \textbf{Construire la phrase espagnole} en respectant l'ordre : Verbe \emph{ser} (temps requis) + Élément mis en relief + Relatif (\emph{lo que/quien/porque/para que}) + Verbe principal (mode requis).
    \item \textbf{Vérifier} : Accords (sujet/verbe), mode (subjonctif si finalité/futur hypothétique), pronoms redoublés si COD/COI topicalisé.
\end{enumerate}
\end{methode}

\courssub{Tableau récapitulatif des correspondances français / espagnol}

\begin{center}
\begin{tabularx}{\linewidth}{|>{\bfseries}X|X|X|}
\hline
\rowcolor{popblueL} \textbf{Tournure française} & \textbf{Structure espagnole canonique} & \textbf{Exemple traduit} \\
\hline
Ce qui est [adj], c'est... & \esp{Lo + adj (masc. sg.) + ser} & \esp{Lo urgente es actuar.} (L'urgent, c'est d'agir.) \\
\hline
Ce que / Ce qui [verbe]... & \esp{Lo que + verbe} & \esp{Lo que quiero es paz.} (Ce que je veux, c'est la paix.) \\
\hline
C'est [Nom chose] qui/que... & \esp{Es + Nom + lo que + verbe} & \esp{Es el trabajo lo que dignifica.} (C'est le travail qui dignifie.) \\
\hline
C'est [Nom personne] qui/que... & \esp{Es + Nom + quien + verbe} & \esp{Es mi padre quien decide.} (C'est mon père qui décide.) \\
\hline
C'est [Prénom/Pronom] que j'ai vu. & \esp{A + Prénom + lo/la + verbe} & \esp{A María la vi ayer.} (Marie, je l'ai vue hier.) \\
\hline
C'est \textbf{parce que} [cause] \textbf{que}... & \esp{Es porque + ind. + que + ind.} & \esp{Es porque llueve que no voy.} (C'est parce qu'il pleut que je n'y vais pas.) \\
\hline
C'est \textbf{pour que} [but] \textbf{que}... & \esp{Es para que + subj. + que + ind.} & \esp{Es para que estudies que te doy el libro.} (C'est pour que tu étudies que je te donne le livre.) \\
\hline
Ça, c'est [adj] ! / Vraiment ! & \esp{¡Eso sí que es + adj!} & \esp{¡Eso sí que es caro!} (Ça, c'est vraiment cher !) \\
\hline
\end{tabularx}
\end{center}

\courssub{Carte mentale : L'ÉMPHASE EN ESPAGNOL}

\begin{popfigure}
\begin{center}\tcbox[colback=popdark,colframe=popdark,coltext=white,arc=9pt,boxsep=4pt,left=6pt,right=6pt]{\bfseries\small ÉNFASIS}\end{center}
\vspace{-2pt}
\begin{tcbraster}[raster columns=2,raster equal height=rows,raster column skip=6pt,raster row skip=6pt,enhanced,blankest]
\begin{tcolorbox}[enhanced,colback=white,colframe=popblue,boxrule=0.9pt,arc=3pt,title={Lo + Adj},fonttitle=\bfseries\scriptsize,coltitle=white,colbacktitle=popblue,left=4pt,right=4pt,top=2pt,bottom=2pt,boxsep=1pt,toptitle=1.5pt,bottomtitle=1.5pt,halign=left]
\begin{itemize}[nosep,leftmargin=*,itemsep=1pt,font=\scriptsize]
\item Lo importante
\item Lo difícil
\item Invariable M.Sg.
\end{itemize}
\end{tcolorbox}
\begin{tcolorbox}[enhanced,colback=white,colframe=poporange,boxrule=0.9pt,arc=3pt,title={Lo que + Verbe},fonttitle=\bfseries\scriptsize,coltitle=white,colbacktitle=poporange,left=4pt,right=4pt,top=2pt,bottom=2pt,boxsep=1pt,toptitle=1.5pt,bottomtitle=1.5pt,halign=left]
\begin{itemize}[nosep,leftmargin=*,itemsep=1pt,font=\scriptsize]
\item Sujet: Lo que quiero
\item COD: No sé lo que dices
\item Action / Fait
\end{itemize}
\end{tcolorbox}
\begin{tcolorbox}[enhanced,colback=white,colframe=poppurple,boxrule=0.9pt,arc=3pt,title={Ser + X + lo que/quien},fonttitle=\bfseries\scriptsize,coltitle=white,colbacktitle=poppurple,left=4pt,right=4pt,top=2pt,bottom=2pt,boxsep=1pt,toptitle=1.5pt,bottomtitle=1.5pt,halign=left]
\begin{itemize}[nosep,leftmargin=*,itemsep=1pt,font=\scriptsize]
\item Chose: Es el esfuerzo lo que...
\item Personne: Es Juan quien...
\item C'est X qui/que...
\end{itemize}
\end{tcolorbox}
\begin{tcolorbox}[enhanced,colback=white,colframe=popgreen,boxrule=0.9pt,arc=3pt,title={Dédoublement / Inversion},fonttitle=\bfseries\scriptsize,coltitle=white,colbacktitle=popgreen,left=4pt,right=4pt,top=2pt,bottom=2pt,boxsep=1pt,toptitle=1.5pt,bottomtitle=1.5pt,halign=left]
\begin{itemize}[nosep,leftmargin=*,itemsep=1pt,font=\scriptsize]
\item A Juan lo vi
\item Vino Juan (VS)
\item ¡Sí que es!
\end{itemize}
\end{tcolorbox}
\end{tcbraster}

\legende{Carte mentale des quatre piliers de l'emphase espagnole (niveau Terminale).}
\end{popfigure}

\courssub{Entraînement : Thème et Version résolus}

\begin{exoresolu}[Thème : Emphase sur cause, finalité et « ce qui/ce que »]
\textbf{Énoncé :} Traduire en espagnol :
\begin{enumerate}
    \item Ce qui est essentiel, c'est de ne pas abandonner.
    \item C'est parce qu'il a de l'expérience qu'il a obtenu le poste.
    \item C'est pour que vous compreniez que je le répète.
    \item Ce que le patron veut, c'est la ponctualité.
    \item C'est mon oncle qui m'a prêté l'argent.
    \item Aux employés, nous leur verserons la prime demain.
\end{enumerate}

\tcblower
\textbf{Solution détaillée :}

\begin{enumerate}
    \item \esp{\textbf{Lo esencial} es no rendirse.}
    \begin{itemize}
        \item « Ce qui est essentiel » $\rightarrow$ Qualité abstraite $\rightarrow$ \textbf{Lo + adj (masc. sg.)}.
        \item « C'est de ne pas abandonner » $\rightarrow$ Infinitif sujet réel.
    \end{itemize}
    \item \esp{\textbf{Es porque tiene experiencia} \textbf{que} ha conseguido el puesto.}
    \begin{itemize}
        \item Emphase sur cause $\rightarrow$ \textbf{Es porque + indicatif (tiene) + que + indicatif (ha conseguido)}.
        \item Attention : pas *\emph{por que}, mais \emph{porque} (conjonction).
    \end{itemize}
    \item \esp{\textbf{Es para que lo entiendan} \textbf{que} lo repito.}
    \begin{itemize}
        \item Emphase sur finalité $\rightarrow$ \textbf{Es para que + SUBJONCTIF (entiendan, 3e pers. pl. \emph{ustedes}) + que + indicatif (repito)}.
        \item \emph{Entender} $\rightarrow$ \emph{entiendan} (subj. présent). \emph{Nota :} En contexte camerounais (norme \emph{ustedes}), on évite \emph{entendáis} (\emph{vosotros}).
    \end{itemize}
    \item \esp{\textbf{Lo que quiere el jefe} es la puntualidad.}
    \begin{itemize}
        \item « Ce que le patron veut » $\rightarrow$ Proposition nominale sujet $\rightarrow$ \textbf{Lo que + verbe (quiere)}.
        \item Ordre possible : \esp{Es la puntualidad lo que quiere el jefe.} (Emphase sur \emph{la puntualidad}).
    \end{itemize}
    \item \esp{\textbf{Es mi tío quien} me ha prestado el dinero.}
    \begin{itemize}
        \item Personne mise en relief $\rightarrow$ \textbf{Es + Nom + QUIEN}.
        \item Pas *\emph{lo que} (réservé aux choses).
    \end{itemize}
    \item \esp{\textbf{A los empleados} \textbf{les} pagaremos la prima mañana.}
    \begin{itemize}
        \item Topicalisation COI animé $\rightarrow$ \textbf{A + Nom + pronom atone COI (les)}.
        \item Futur simple \emph{pagaremos}.
    \end{itemize}
\end{enumerate}
\end{exoresolu}

\begin{exoresolu}[Version : Repérer et rendre l'emphase]
\textbf{Énoncé :} Traduire en français le paragraphe suivant en respectant les nuances emphatiques :
\esp{« No es el título universitario lo que abre las puertas, \textbf{es porque} tienes contactos \textbf{que} consigues la entrevista. \textbf{Lo que realmente importa} es la actitud. \textbf{Es para que} demuestres tu valor \textbf{que} te dan el periodo de prueba. \textbf{A ti} \textbf{te} toca aprovecharlo. ¡\textbf{Sí que} puedes! »}

\tcblower
\textbf{Traduction commentée :}
« Ce n'est \textbf{pas le diplôme universitaire qui} ouvre les portes, \textbf{c'est parce que} tu as des contacts \textbf{que} tu obtiens l'entretien. \textbf{Ce qui compte vraiment}, c'est l'attitude. \textbf{C'est pour que} tu démontres ta valeur \textbf{qu'}on te donne la période d'essai. \textbf{Toi}, c'est à toi d'en profiter. \textbf{Sí que} tu peux ! / Mais oui, tu peux ! »

\begin{itemize}
    \item \esp{No es el título... lo que} $\rightarrow$ Négation de l'emphase sur nom chose : « Ce n'est pas X qui... ».
    \item \esp{Es porque... que} $\rightarrow$ Emphase causale : « C'est parce que... que... ».
    \item \esp{Lo que realmente importa} $\rightarrow$ Sujet nominalisé : « Ce qui compte vraiment ».
    \item \esp{Es para que... que} + \textbf{subj. \emph{demuestres}} $\rightarrow$ Emphase finalité : « C'est pour que... que... ».
    \item \esp{A ti te toca} $\rightarrow$ Topicalisation + redoublement : « Toi, ça te revient / À toi de... ».
    \item \esp{¡Sí que puedes!} $\rightarrow$ \emph{Sí} emphatique : « Mais oui tu peux / Si, tu peux ! ».
\end{itemize}
\end{exoresolu}

\courssub{Exercices d'application}

\begin{exercice}[Thème : Structures emphatiques mixtes]
Traduire en espagnol :
\begin{enumerate}
    \item Ce qui est injuste, c'est qu'ils ne paient pas les heures sup.
    \item C'est la corruption qui freine le développement.
    \item C'est pour que l'équipe gagne que l'entraîneur change de tactique.
    \item Ce que j'ai dit hier n'était pas vrai.
    \item C'est ma sœur qui m'a appris l'espagnol.
    \item Ce livre, je ne l'ai pas lu.
    \item C'est parce qu'il pleut que le match est annulé.
    \item À mon avis, c'est une erreur. (Utiliser \emph{A mí me parece}).
\end{enumerate}
\end{exercice}

\begin{exercice}[Version : Analyse et traduction]
Traduire en français et \textbf{surligner (ou mettre en gras dans votre copie)} la structure emphatique utilisée dans chaque phrase :
\begin{enumerate}
    \item \esp{Lo prioritario es reducir la deuda externa.}
    \item \esp{Es la educación lo que transforma la sociedad.}
    \item \esp{Es para que no haya errores que revisamos todo.}
    \item \esp{A los candidatos los evaluaremos mañana.}
    \item \esp{¡Eso sí que es una buena noticia!}
    \item \esp{Es porque llegó tarde que perdió el tren.}
    \item \esp{Lo que no entiendo es su indiferencia.}
    \item \esp{Es mi abuela quien hace los mejores pasteles.}
\end{enumerate}
\end{exercice}

\begin{corrige}[Exercice 1 - Thème]
\begin{enumerate}
    \item \esp{Lo injusto es que no paguen las horas extra.} (Subjonctif \emph{paguen} car opinion/valeur sur fait non réalisé ou subjectif ; indicatif \emph{no pagan} possible si fait constaté).
    \item \esp{Es la corrupción lo que frena el desarrollo.}
    \item \esp{Es para que gane el equipo que el entrenador cambia de táctica.} (\emph{gane} = subj. finalité).
    \item \esp{Lo que dije ayer no era verdad.} (Prétérit indéfini \emph{dije} / Imparfait \emph{era}).
    \item \esp{Es mi hermana quien me enseñó español.}
    \item \esp{Ese libro no lo he leído.} (Ou \emph{A ese libro no lo he leído} moins courant pour inanimé).
    \item \esp{Es porque llueve que el partido se cancela / está cancelado.}
    \item \esp{A mí me parece un error.}
\end{enumerate}
\end{corrige}

\begin{corrige}[Exercice 2 - Version]
\begin{enumerate}
    \item \textbf{Lo prioritario es...} — « Ce qui est prioritaire, c'est... » (\cle{Lo + adj}).
    \item \textbf{Es la educación lo que...} — « C'est l'éducation qui... » (\cle{Es + Nom + lo que}).
    \item \textbf{Es para que no haya errores que...} — « C'est pour qu'il n'y ait pas d'erreurs que... » (\cle{Es para que + subj. \emph{haya} + que}).
    \item \textbf{A los candidatos los evaluaremos...} — « Les candidats, nous les évaluerons... » (\cle{Dédoublement COD animé}).
    \item \textbf{¡Eso sí que es...!} — « Ça, c'est vraiment... ! » (\cle{Adverbe emphatique \emph{sí}}).
    \item \textbf{Es porque llegó tarde que...} — « C'est parce qu'il est arrivé en retard que... » (\cle{Es porque + ind. + que}).
    \item \textbf{Lo que no entiendo...} — « Ce que je ne comprends pas... » (\cle{Lo que + verbe}).
    \item \textbf{Es mi abuela quien...} — « C'est ma grand-mère qui... » (\cle{Es + Nom personne + quien}).
\end{enumerate}
\end{corrige}

\begin{experience}[Débat oral : « Mérito vs Contactos »]
\textbf{Consigne :} En binômes, préparez un dialogue de 2 minutes entre un \emph{recién graduado} (jeune diplômé) et un \emph{empresario} (chef d'entreprise) à Yaoundé ou Douala.
\textbf{Contraintes linguistiques obligatoires} (chacun doit en utiliser au moins 3 différentes) :
\begin{enumerate}
    \item Une phrase avec \emph{Lo + adjectif}.
    \item Une phrase avec \emph{Lo que + verbe}.
    \item Une emphase \emph{Es... lo que / quien}.
    \item Une emphase causale \emph{Es porque... que}.
    \item Une emphase finalité \emph{Es para que... que} + \textbf{subjonctif}.
    \item Un dédoublement \emph{A + nom + pronom}.
\end{enumerate}
\textbf{Exemple d'amorce :}
\emph{— Lo fundamental no es el título, es la experiencia.}
\emph{— Es porque no tienes contactos que no encuentras trabajo.}
\emph{— A mí me parece que es para que aprendan que exigen experiencia.}
\end{experience}

\begin{savaistu}[L'emphase en Guinée équatoriale : particularités phonétiques et syntaxiques]
En Guinée équatoriale (seul pays africain hispanophone officiel), l'emphase orale est marquée par une \textbf{intonation très montante} sur l'élément topicalisé (souvent plus expressive qu'en Espagne). On observe aussi une fréquence accrue du \textbf{dédoublement pronominal} pour les COD inanimés (\esp{Ese libro lo leí} $\rightarrow$ \esp{A ese libro lo leí} par analogie avec l'animé) et l'usage de \emph{es que} explicatif en tête de phrase (\esp{Es que... lo que pasa es que...}) comme marqueur discursif très fort, calqué sur les structures bubi/fang de focalisation. Pour le Bac, restez sur la norme panhispanique standard décrite ci-dessus.
\end{savaistu}

\coursec{Entraînement à la traduction : thème et version}

Après avoir révisé les structures d'emphase (\emph{lo que, lo de, el hecho de, ser + infinitivo, cleft sentences}), nous passons maintenant à la mise en pratique intégrée. La traduction, au Baccalauréat comme dans la vie professionnelle, n'est pas un exercice de transposition mot à mot : c'est un travail de \textbf{réécriture} qui exige de repérer la \emph{fonction} de chaque segment (cause, finalité, condition temporelle, focalisation) pour choisir l'outil grammatical espagnol adéquat. Cette section vous guide pas à pas, de la version collective au thème individuel, avec une grille d'auto-évaluation calquée sur les attendus du jury.

\courssub{Observation guidée : repérer les structures}

\begin{textebox}[Texte support — Insertion professionnelle à Yaoundé]
¿Por qué tantos jóvenes aceptan trabajos precarios? \textbf{Porque} la necesidad es urgente. \textbf{Como} el mercado formal ofrece pocas plazas, \textbf{debido a} la saturación, muchos optan por el sector informal. \textbf{Para} sobrevivir, venden en los cruces. El Estado lanza programas \textbf{para que} accedan a microcréditos. \textbf{Cuando} terminen su formación, \textbf{podrán} formalizar su actividad. \textbf{Lo fundamental} \textbf{es} creer en su potencial.
\end{textebox}

\begin{methode}[Observation active : 5 questions pour décoder la phrase]
\begin{enumerate}
\item \textbf{Soulignez} tous les connecteurs logiques (\esp{porque, como, debido a, para, para que, cuando, lo fundamental es}).
\item \textbf{Classez}-les : Cause (\esp{porque, como, debido a}) / Finalité (\esp{para, para que}) / Temps futur (\esp{cuando}) / Emphase (\esp{lo fundamental es}).
\item \textbf{Identifiez le mode} du verbe qui suit : Indicatif (\esp{es, ofrece, optan, lanzan, podrán}) / Subjonctif (\esp{accedan, terminen}) / Infinitif (\esp{survivir, creer}).
\item \textbf{Justifiez} : Pourquoi \esp{para que accedan} (subj.) et pas \esp{*para acceder} ? (Sujets différents : \emph{El Estado} / \emph{ellos}). Pourquoi \esp{cuando terminen} (subj.) et pas \esp{*cuando terminan} ? (Action future non réalisée).
\item \textbf{Traduisez} mentalement en français en respectant la structure (ex. \esp{debido a la saturación} $\to$ \emph{en raison de la saturation}, nom + préposition).
\end{enumerate}
\end{methode}

\courssub{La phrase causale}

\begin{propriete}[Connecteurs de cause : conjugaison + nuance]
\begin{center}
\begin{tabularx}{\linewidth}{|l|X|X|}
\hline
\rowcolor{popblueL}
\textbf{Connecteur} & \textbf{Structure / Mode} & \textbf{Nuance / Registre} \\
\hline
\esp{porque} & + \textbf{Indicatif} (tous temps) & Réponse à \emph{¿por qué?}, cause nouvelle, insistance sur la raison. \\
\hline
\esp{como} & + \textbf{Indicatif} & Cause connue, évidente, en tête de phrase. \emph{« Comme il pleut, je reste. »} \\
\hline
\esp{ya que} & + \textbf{Indicatif} & Cause évidente, argumentative. \emph{« Puisque / vu que ».} \\
\hline
\esp{puesto que} & + \textbf{Indicatif} & Cause admise, formelle. \emph{« Attendu que / dans la mesure où ».} \\
\hline
\esp{dado que} & + \textbf{Indicatif} & Cause factuelle, écrite, administrative. \emph{« Étant donné que ».} \\
\hline
\esp{a causa de} / \esp{debido a} & + \textbf{Nom / Infinitif} & Locutions prépositionnelles (pas de verbe conjugué). \emph{« À cause de / en raison de ».} \\
\hline
\end{tabularx}
\end{center}
\end{propriete}

\begin{attention}[Pièges francophones : Cause]
\begin{itemize}
\item \textbf{\esp{Porque} vs \esp{Por qué} vs \esp{Por que} vs \esp{Porqué}} : \esp{Porque} (conjonction \emph{car}), \esp{Por qué} (interrogatif/exclamatif \emph{pourquoi}), \esp{Por que} (préposition + pronom relatif \emph{pour lequel}), \esp{Porqué} (nom masculin \emph{la raison}).
\item \textbf{\esp{Debido a} / \esp{A causa de}} exigent un \textbf{nom} ou un \textbf{infinitif}, jamais une proposition conjuguée. \emph{Faux :} \esp{*Debido a que llueve} $\to$ \textbf{Correct :} \esp{Debido a la lluvia} / \esp{Debido a llover}.
\item \textbf{\esp{Gracias a}} = cause positive (\emph{grace à}). Ne pas utiliser pour une cause négative. \emph{Faux :} \esp{*Gracias a la crisis, perdí mi trabajo} $\to$ \textbf{Correct :} \esp{A causa de la crisis...}
\end{itemize}
\end{attention}

\begin{exemplebox}[Cause : exemples traduits]
\begin{itemize}
\item \esp{No vine \textbf{porque} estaba enfermo.} $\to$ \emph{Je ne suis pas venu \textbf{car} j'étais malade.} (Réponse à pourquoi, indicatif).
\item \esp{\textbf{Como} no hay dinero, \textbf{no} podemos ir.} $\to$ \emph{\textbf{Comme} il n'y a pas d'argent, nous ne pouvons pas y aller.} (Cause connue, indicatif).
\item \esp{El proyecto fracasó \textbf{debido a} la falta de fondos.} $\to$ \emph{Le projet a échoué \textbf{en raison du} manque de fonds.} (Nom).
\item \esp{\textbf{Dado que} el mercado está saturado, \textbf{innovamos}.} $\to$ \emph{\textbf{Étant donné que} le marché est saturé, nous innovons.} (Formel, indicatif).
\end{itemize}
\end{exemplebox}

\courssub{L'expression de la finalité}

\begin{propriete}[Finalité : le choix du connecteur selon l'identité des sujets]
\begin{center}
\begin{tabularx}{\linewidth}{|l|X|X|}
\hline
\rowcolor{popgreenL}
\textbf{Structure} & \textbf{Condition} & \textbf{Traduction française} \\
\hline
\esp{Para} + \textbf{Infinitif} & \textbf{Même sujet} (sujet principal = sujet de la finalité) & \emph{Pour} + infinitif / \emph{Afin de} + infinitif \\
\hline
\esp{Para que} + \textbf{Subjonctif Présent} & \textbf{Sujets différents} & \emph{Afin que / Pour que} + subjonctif \\
\hline
\esp{A fin de que} + \textbf{Subjonctif Présent} & Sujets différents (soutenu, administratif) & \emph{Afin que / Dans le but que} + subjonctif \\
\hline
\esp{Con el fin de que} + \textbf{Subjonctif Présent} & Sujets différents (très formel, juridique) & \emph{Dans le but que} + subjonctif \\
\hline
\esp{Con el fin de} / \esp{A fin de} + \textbf{Infinitif} & Même sujet (formel) & \emph{Dans le but de} / \emph{Afin de} + infinitif \\
\hline
\end{tabularx}
\end{center}
\end{propriete}

\begin{attention}[Pièges francophones : Finalité]
\begin{itemize}
\item \textbf{Règle d'or :} \cle{Même sujet $\to$ Infinitif (\esp{para} + inf.)} | \cle{Sujets différents $\to$ Subjonctif (\esp{para que} + subj.)}. Ne jamais écrire \esp{*para que yo voy} (indicatif) ni \esp{*para ir} si les sujets diffèrent.
\item \textbf{Négation :} \esp{Para que no} + subjonctif. \esp{Trabajo para que no falte nada.} (\emph{Je travaille pour qu'il ne manque rien}).
\item \textbf{\esp{Para} + Infinitif passé} (\esp{para haber hecho}) : finalité antérieure (rare, soutenu).
\item \textbf{Faux ami :} \esp{Por} + infinitif = cause/motif (\emph{à cause de}), non finalité. \esp{Lo hice por ayudarte} (\emph{Je l'ai fait par aide pour toi / à cause de t'aider}) vs \esp{Lo hice para ayudarte} (\emph{Je l'ai fait pour t'aider}).
\end{itemize}
\end{attention}

\begin{exemplebox}[Finalité : exemples traduits]
\begin{itemize}
\item \esp{Estudio \textbf{para} ser médico.} (Même sujet) $\to$ \emph{J'étudie \textbf{pour} devenir médecin.}
\item \esp{El profesor explica \textbf{para que} los alumnos \textbf{entiendan}.} (Sujets diff.) $\to$ \emph{Le prof explique \textbf{afin que} les élèves \textbf{comprennent}.}
\item \esp{El MINJEC lanza becas \textbf{a fin de que} los jóvenes \textbf{emprendan}.} (Soutenu) $\to$ \emph{Le MINJEC lance des bourses \textbf{afin que} les jeunes \textbf{entreprennent}.}
\item \esp{Los emprendedores invierten \textbf{con el fin de} \textbf{crecer}.} (Même sujet, formel) $\to$ \emph{Les entrepreneurs investissent \textbf{dans le but de} \textbf{grandir}.}
\end{itemize}
\end{exemplebox}

\courssub{Le futur dans la subordonnée : le subjonctif}

\begin{propriete}[Subordonnées temporelles / conditionnelles à valeur future]
Lorsque l'action de la subordonnée est \textbf{postérieure} à l'énonciation (futur, hypothèse non réalisée), l'espagnol utilise le \textbf{Subjonctif Présent}. Le principal est au \textbf{Futur Simple} ou à l'\textbf{Impératif}.
\begin{center}
\begin{tabularx}{\linewidth}{|l|X|X|}
\hline
\rowcolor{poporangeL}
\textbf{Connecteur} & \textbf{Mode / Temps} & \textbf{Équivalent français} \\
\hline
\esp{Cuando, En cuanto, Tan pronto como, Una vez que} & \textbf{Subjonctif Présent} & \emph{Quand, Dès que, Aussitôt que} + \textbf{Futur Antérieur} \\
\hline
\esp{Hasta que} & \textbf{Subjonctif Présent} & \emph{Jusqu'à ce que} + \textbf{Subjonctif / Futur Antérieur} \\
\hline
\esp{Si} (condition réelle) & \textbf{Indicatif Présent} & \emph{Si} + \textbf{Présent} (principal au Futur) \\
\hline
\esp{Si} (condition hypothétique/irréelle) & \textbf{Imparfait Subjonctif} & \emph{Si} + \textbf{Imparfait} (principal au Conditionnel) \\
\hline
\end{tabularx}
\end{center}
\end{propriete}

\begin{attention}[Piège majeur : \esp{Si} vs \esp{CUANDO}]
\begin{itemize}
\item \textbf{\esp{Si} + INDICATIF PRÉSENT} = Condition réelle, possible. \esp{Si llueve, no iré.} (\emph{S'il pleut, j'irai pas}).
\item \textbf{\esp{Cuando} + SUBJONCTIF PRÉSENT} = Certitude temporelle future. \esp{Cuando llueva, me mojaré.} (\emph{Quand il pleuvra / Dès qu'il pleuvra, je serai mouillé}).
\item \textbf{Jamais} \esp{*Cuando lloverá} (futur indicatif en subordonnée). \textbf{Jamais} \esp{*Si lloveré}.
\end{itemize}
\end{attention}

\begin{exemplebox}[Futur dans la subordonnée : exemples traduits]
\begin{itemize}
\item \esp{\textbf{Cuando termines} el Bachillerato, \textbf{buscarás} trabajo.} $\to$ \emph{\textbf{Quand tu auras terminé} le Bac, tu \textbf{chercheras} du travail.}
\item \esp{\textbf{En cuanto} el ministro \textbf{firme} el decreto, \textbf{empezaremos}.} $\to$ \emph{\textbf{Dès que} le ministre \textbf{aura signé} le décret, nous \textbf{commencerons}.}
\item \esp{No me voy \textbf{hasta que} \textbf{llegues}.} $\to$ \emph{Je ne pars pas \textbf{jusqu'à ce que} tu \textbf{arrives} (tu auras arrivé).}
\item \esp{\textbf{Si} apruebas el examen, \textbf{te regalaré} un móvil.} $\to$ \emph{\textbf{Si} tu \textbf{réussis} l'examen, je \textbf{t'offrirai} un portable.} (Indicatif !)
\end{itemize}
\end{exemplebox}

\begin{experience}[Atelier : Conjugaison éclair — Subjonctif Présent (5 min)]
L'enseignant dicte des verbes (réguliers et irréguliers fréquents), les élèves écrivent la forme du subjonctif présent aux 3 personnes du singulier et 3ème pluriel.
\begin{center}
\begin{tabular}{|l|l|l|l|l|l|l|}
\hline
\rowcolor{poppurpleL}
\textbf{Infinitif} & \textbf{Yo} & \textbf{Tú} & \textbf{Él/Ella/Ud.} & \textbf{Nosotros} & \textbf{Vosotros} & \textbf{Ellos/Uds.} \\
\hline
\esp{Hablar} & \esp{hable} & \esp{hables} & \esp{hable} & \esp{hablemos} & \esp{habléis} & \esp{hablen} \\
\hline
\esp{Comer} & \esp{coma} & \esp{comas} & \esp{coma} & \esp{comamos} & \esp{comáis} & \esp{coman} \\
\hline
\esp{Vivir} & \esp{viva} & \esp{vivas} & \esp{viva} & \esp{vivamos} & \esp{viváis} & \esp{vivan} \\
\hline
\esp{Tener} & \esp{tenga} & \esp{tengas} & \esp{tenga} & \esp{tengamos} & \esp{tengáis} & \esp{tengan} \\
\hline
\esp{Poder} & \esp{pueda} & \esp{puedas} & \esp{pueda} & \esp{podamos} & \esp{podáis} & \esp{puedan} \\
\hline
\esp{Saber} & \esp{sepa} & \esp{sepas} & \esp{sepa} & \esp{sepamos} & \esp{sepáis} & \esp{sepan} \\
\hline
\esp{Ir} & \esp{vaya} & \esp{vayas} & \esp{vaya} & \esp{vayamos} & \esp{vayáis} & \esp{vayan} \\
\hline
\esp{Ser/Estar} & \esp{sea/esté} & \esp{seas/estés} & \esp{sea/esté} & \esp{seamos/estemos} & \esp{seáis/estéis} & \esp{sean/estén} \\
\hline
\end{tabular}
\end{center}
\emph{Consigne :} Mémorisez les irréguliers (\esp{tenga, pueda, sepa, vaya, sea, esté, dé, haya, vaya}) : ils tombent à tous les Bacs.
\end{experience}

\courssub{Révision de l'emphase}

\begin{propriete}[Structures d'emphase (Focalisation) : formes et emplois]
\begin{center}
\begin{tabularx}{\linewidth}{|l|X|X|}
\hline
\rowcolor{poppinkL}
\textbf{Structure espagnole} & \textbf{Français (Cleft sentence)} & \textbf{Élément focalisé} \\
\hline
\esp{Lo que} + verbe + \esp{ser} + nom/inf. & \emph{Ce qui / Ce que ... c'est ...} & Proposition entière (sujet/objet) \\
\hline
\esp{Lo de} + nom / infinitif + \esp{ser} + adj/nom & \emph{L'affaire du / Le fait de ... c'est ...} & Nom / Action (nominalisée) \\
\hline
\esp{El hecho de que} + \textbf{Subjonctif} + \esp{ser} + adj/nom & \emph{Le fait que ... c'est ...} & Fait accompli (réel, subj.) \\
\hline
\esp{Lo} + \textbf{Adjectif} + \esp{ser} + infinitif/nom & \emph{L'... / Ce qui est ... c'est ...} & Qualité / Valeur (adj. neutre) \\
\hline
\esp{Ser} + \textbf{sujet focalisé} + \esp{quien/el que/lo que} + verbe & \emph{C'est ... qui/que ...} & Sujet / Objet / Complément \\
\hline
\end{tabularx}
\end{center}
\end{propriete}

\begin{attention}[Pièges francophones : Emphase]
\begin{itemize}
\item \textbf{\esp{Lo que} $\neq$ \emph{Ce qui} seul.} Il faut le verbe \esp{ser} derrière. \emph{Faux :} \esp{*Lo que me gusta es la música} (Correct car \esp{es} présent). \emph{Faux :} \esp{*Lo que quiero ir} $\to$ \textbf{Correct :} \esp{Lo que quiero \textbf{es} ir}.
\item \textbf{Accord du verbe \esp{ser}} : Singulier si l'attribut est singulier/infinitif (\esp{Lo importante \textbf{es} estudiar}). Pluriel si attribut pluriel (\esp{Lo importante \textbf{son} los estudios}).
\item \textbf{\esp{Lo de}} : Très utile pour nominaliser. \esp{Lo de Juan \textbf{es} un misterio.} (\emph{L'affaire Juan / Ce qui concerne Juan est un mystère}).
\item \textbf{\esp{El hecho de que}} impose le \textbf{Subjonctif} (fait réel mais subjectivisé). \esp{El hecho de que \textbf{haya} venido \textbf{es} importante.}
\end{itemize}
\end{attention}

\begin{exemplebox}[Emphase : exemples traduits]
\begin{itemize}
\item \esp{\textbf{Lo que} me preocupa \textbf{es} el desempleo.} $\to$ \emph{\textbf{Ce qui} m'inquiète, \textbf{c'est} le chômage.}
\item \esp{\textbf{Lo de} la corrupción \textbf{es} vergonzoso.} $\to$ \emph{\textbf{L'affaire de} la corruption / \textbf{Le fait de} la corruption \textbf{est} honteux.}
\item \esp{\textbf{El hecho de que} \textbf{haya} paro \textbf{es} grave.} $\to$ \emph{\textbf{Le fait qu'}il y \textbf{ait} du chômage \textbf{est} grave.}
\item \esp{\textbf{Lo fundamental} \textbf{es} perseverar.} $\to$ \emph{\textbf{L'essentiel / Ce qui est fondamental}, \textbf{c'est de} persévérer.}
\item \esp{\textbf{Es} el gobierno \textbf{quien} debe actuar.} $\to$ \emph{\textbf{C'est} le gouvernement \textbf{qui} doit agir.}
\end{itemize}
\end{exemplebox}

\courssub{Version guidée : jeunesse et travail au Cameroun}

\begin{textebox}[Version guidée — Jeunesse et travail au Cameroun]
Como la economía es difícil, muchos jóvenes trabajan para ayudar a sus familias. El gobierno crea programas para que encuentren empleo. Cuando terminen sus estudios, podrán contribuir al desarrollo del país. Lo importante es no rendirse.
\end{textebox}

\begin{methode}[Méthode de la version : du sens global au mot juste]
\begin{enumerate}
\item \textbf{Lire pour comprendre} : identifier le sujet, le ton, le registre (ici : constat social, ton engagé, actualité camerounaise).
\item \textbf{Repérer les connecteurs logiques} : souligner \esp{como}, \esp{para}, \esp{para que}, \esp{cuando}, \esp{lo ... es}.
\item \textbf{Choisir l'équivalent français} :
\begin{itemize}
\item \esp{como} $\to$ \emph{puisque / comme} (cause connue, indicatif).
\item \esp{para} + infinitif $\to$ \emph{pour} + infinitif (même sujet).
\item \esp{para que} + subjonctif $\to$ \emph{afin que / pour que} + subjonctif (sujets différents).
\item \esp{cuando} + subjonctif $\to$ \emph{quand / dès que} + \textbf{futur antérieur}.
\item \esp{lo ... es} $\to$ \emph{l'important c'est / ce qui compte c'est} + infinitif.
\end{itemize}
\item \textbf{Vérifier les modes et temps} : Indicatif cause réelle (\esp{es, trabajan, crea}), Subjonctif finalité/futur (\esp{encuentren, terminen}), Futur principal (\esp{podrán}).
\item \textbf{Relire en français naturel} : éviter le calque (\emph{« Comme l'économie est difficile »} $\to$ \emph{« Puisque l'économie est difficile »} ou \emph{« L'économie étant difficile »}).
\end{enumerate}
\end{methode}

\begin{exemplebox}[Traduction commentée]
\begin{itemize}
\item \esp{Como la economía es difícil} $\to$ \emph{Puisque l'économie est difficile} (cause avérée, indicatif présent).
\item \esp{muchos jóvenes trabajan para ayudar a sus familias} $\to$ \emph{Beaucoup de jeunes travaillent pour aider leur famille} (\esp{para} + infinitif = finalité, même sujet).
\item \esp{El gobierno crea programas para que encuentren empleo} $\to$ \emph{Le gouvernement crée des programmes afin qu'ils trouvent un emploi} (changement de sujet $\to$ \esp{para que} + subjonctif présent).
\item \esp{Cuando terminen sus estudios, podrán contribuir al desarrollo del país} $\to$ \emph{Quand ils auront terminé leurs études, ils pourront contribuer au développement du pays} (\esp{cuando} + subjonctif présent = valeur de futur antérieur ; principal au futur).
\item \esp{Lo importante es no rendirse} $\to$ \emph{L'important, c'est de ne pas baisser les bras} (emphase par \esp{lo + adjectif + ser + infinitif}).
\end{itemize}
\end{exemplebox}

\begin{exoresolu}[Version guidée — Correction collective justifiée]
\textbf{Énoncé :} Traduisez le texte ci-dessus en français en justifiant chaque choix grammatical.

\tcblower
\textbf{Solution détaillée :}

\emph{« Puisque l'économie est difficile, beaucoup de jeunes travaillent pour aider leur famille. Le gouvernement met en place des programmes afin qu'ils trouvent un emploi. Quand ils auront terminé leurs études, ils pourront contribuer au développement du pays. L'important, c'est de ne pas baisser les bras. »}

\begin{itemize}
\item \textbf{\esp{Como} + indicatif} : cause réelle, connue de l'énonciateur $\to$ \emph{puisque / comme / attendu que}. \emph{Distinction :} \esp{porque} répond à \emph{¿por qué?} (information nouvelle), \esp{como} introduit une cause présupposée.
\item \textbf{\esp{para} + infinitif} : finalité, même sujet (\emph{los jóvenes} travaillent \emph{et} aident) $\to$ \emph{pour} + infinitif.
\item \textbf{\esp{para que} + subjonctif} : finalité, sujets différents (\emph{el gobierno} crée $\neq$ \emph{ellos} trouvent) $\to$ \emph{afin que / pour que} + subjonctif présent.
\item \textbf{\esp{cuando} + subjonctif présent} : subordonnée temporelle se rapportant à l'avenir $\to$ français : \emph{quand / dès que / aussitôt que} + \textbf{futur antérieur} (\emph{auront terminé}). \emph{Règle :} \esp{cuando / en cuanto / tan pronto como / hasta que} + subjonctif si action future non encore réalisée.
\item \textbf{\esp{Lo importante es no rendirse}} : structure emphatique \emph{lo + adj. + ser + inf.} $\to$ \emph{L'important / Ce qui compte, c'est de} + infinitif. Variante : \emph{Il faut surtout ne pas baisser les bras}.
\end{itemize}
\end{exoresolu}

\courssub{Thème guidé : difficulté progressive}

Appliquez la \textbf{stratégie en 4 étapes} (voir organigramme p.~\pageref{fig:organigramme}) à chaque phrase. Écrivez d'abord le connecteur espagnol, puis le mode, puis la phrase complète.

\begin{exercice}[Thème guidé — De la phrase simple au paragraphe]
Traduisez en espagnol :
\begin{enumerate}
\item Comme il pleut, nous restons à la maison pour réviser.
\item Je t'expliquerai la leçon pour que tu comprennes.
\item Le professeur donne des devoirs afin que les élèves progressent.
\item Quand tu auras fini, appelle-moi.
\item Ce qui m'inquiète, c'est le chômage des jeunes.
\item Puisque le marché est saturé, les entrepreneurs créent des entreprises émergentes pour innover.
\item Le ministère lance des formations pour que les diplômés trouvent du travail.
\item Dès que le projet sera approuvé, nous commencerons.
\item L'essentiel est de persévérer.
\end{enumerate}
\end{exercice}

\begin{corrige}[Exercice : Thème guidé]
\begin{enumerate}
\item \esp{Como llueve, nos quedamos en casa para repasar.} (Cause \esp{como} + ind. ; finalité même sujet \esp{para} + inf.)
\item \esp{Te explicaré la lección para que entiendas.} (Finalité sujets différents $\to$ \esp{para que} + subj. présent \esp{entiendas}).
\item \esp{El profesor da deberes para que los alumnos progresen.} (\esp{para que} + subj. \esp{progresen}).
\item \esp{Cuando termines, llámame.} (\esp{cuando} + subj. \esp{termines} = futur antérieur ; impératif principal).
\item \esp{Lo que me preocupa es el desempleo juvenil.} (Emphase \esp{lo que} + verbe + \esp{ser} + nom). Variante : \esp{Lo que me inquieta...}
\item \esp{Puesto que el mercado está saturado, los emprendedores crean empresas emergentes para innovar.} (Variété connecteur cause : \esp{puesto que} + ind. ; \esp{para} + inf. même sujet).
\item \esp{El ministerio lanza formaciones para que los titulados encuentren trabajo.} (\esp{para que} + subj. \esp{encuentren}).
\item \esp{En cuanto se apruebe el proyecto, empezaremos.} (\esp{en cuanto} + subj. \esp{se apruebe} ; futur principal \esp{empezaremos}).
\item \esp{Lo fundamental es perseverar.} (Emphase \esp{lo + adj.} + \esp{ser} + inf.).
\end{enumerate}
\end{corrige}

\begin{attention}[Piège fatal au Bac : le mode après le connecteur]
En thème, \textbf{écrivez d'abord le connecteur + le mode du verbe} avant de compléter la phrase. C'est là que se perdent 80~\% des points de traduction.
\begin{itemize}
\item \esp{Para que} $\to$ \textbf{subjonctif présent} (jamais indicatif, jamais infinitif).
\item \esp{Cuando} (futur) $\to$ \textbf{subjonctif présent} (valeur futur antérieur).
\item \esp{Porque / como / puesto que / ya que / dado que} $\to$ \textbf{indicatif} (cause réelle).
\item \esp{Para} + infinitif $\to$ \textbf{même sujet} ; \esp{para que} + subjonctif $\to$ \textbf{sujets différents}.
\item \esp{Si} (réel) $\to$ \textbf{indicatif présent} ; \esp{Si} (irréel) $\to$ \textbf{imparfait subjonctif}.
\end{itemize}
\textit{Exemple :} \emph{« Je travaille pour que ma mère soit fière. »} $\to$ \esp{Trabajo} \underline{\textbf{para que}} \underline{\textbf{mi madre esté}} \esp{orgullosa.} (Pas \esp{*para que mi madre está} ; pas \esp{*para estar} car sujets différents).
\end{attention}

\courssub{Stratégie de traduction : l'organigramme « Lire $\to$ Repérer $\to$ Choisir $\to$ Vérifier $\to$ Relire »}

\begin{popfigure}
\begin{tikzpicture}[
  node distance=1.2cm and 1.5cm,
  box/.style={rectangle, rounded corners, draw=popblue, thick, fill=popblueL, text width=2.8cm, align=center, font=\small},
  arrow/.style={-Stealth, thick, popdark},
  decision/.style={diamond, aspect=2, draw=poporange, thick, fill=poporangeL, text width=2.5cm, align=center, font=\small}
]
\node[box, align=center] (lire){Lire le texte\\source (sens global)};
\node[box, below=of lire, align=center] (reperer){Repérer les\\connecteurs logiques};
\node[decision, below=of reperer, align=center] (choisir){Choisir l'équivalent\\espagnol + mode};
\node[box, below=of choisir, align=center] (verifier){Vérifier :\\accords, temps, accents};
\node[box, below=of verifier, align=center] (relire){Relire en espagnol\\naturel et fluide};
\draw[arrow] (lire) -- (reperer);
\draw[arrow] (reperer) -- (choisir);
\draw[arrow] (choisir) -- (verifier);
\draw[arrow] (verifier) -- (relire);
\node[align=center, font=\tiny, popmuted, right=0.3cm of choisir] {Cause : \esp{porque/como/puesto que} + IND\\
Finalité même sujet : \esp{para} + INF\\
Finalité sujets diff. : \esp{para que} + SUBJ\\
Futur sub. : \esp{cuando/en cuanto} + SUBJ\\
Emphase : \esp{lo que/lo de/lo + adj}};
\end{tikzpicture}
\legende{Organigramme de la méthode de traduction — À mémoriser pour l'épreuve du Bac.}
\label{fig:organigramme}
\end{popfigure}

\courssub{Grille d'auto-évaluation (à utiliser pour chaque phrase)}

\begin{center}
\begin{tabularx}{\linewidth}{|X|c|c|c|c|}
\hline
\rowcolor{popblueL}
\textbf{Critère} & \textbf{Connecteur correct} & \textbf{Mode correct} & \textbf{Temps correct} & \textbf{Emphase respectée} \\
\hline
Cause (\emph{puisque, comme, car}) & \esp{como / ya que / puesto que / dado que} + IND & Indicatif & Présent / Imparfait / Passé composé & — \\
\hline
Finalité même sujet (\emph{pour}) & \esp{para} + INF & Infinitif & — & — \\
\hline
Finalité sujets diff. (\emph{afin que}) & \esp{para que / a fin de que / con el fin de que} + SUBJ & Subjonctif présent & Présent (valeur futur) & — \\
\hline
Futur dans la sub. (\emph{quand, dès que}) & \esp{cuando / en cuanto / tan pronto como / hasta que} + SUBJ & Subjonctif présent & Valeur futur antérieur & — \\
\hline
Condition réelle (\emph{si}) & \esp{si} + IND & Indicatif présent & Présent (valeur futur) & — \\
\hline
Emphase (\emph{ce qui compte, l'important}) & \esp{lo que / lo de / lo + adj. + ser} & Indicatif (souvent) & Selon contexte & Structure cleft conservée \\
\hline
\end{tabularx}
\end{center}

\courssub{Correction collective avec justification par la règle}

\begin{experience}[Correction collective guidée — 20 min]
\begin{enumerate}
\item L'enseignant projette une phrase du thème (ex. n°6).
\item Un élève propose sa traduction au tableau.
\item La classe vote : \emph{correct / à améliorer / faux}.
\item L'enseignant demande : \textbf{« Quel connecteur as-tu choisi ? Pourquoi ce mode ? Quel temps ? L'emphase est-elle rendue ? »}
\item L'élève justifie en citant la règle (ex. : \emph{« \esp{Puesto que} introduit une cause connue, donc indicatif ; \esp{para} + infinitif car même sujet \emph{emprendedores} »}).
\item L'enseignant valide ou corrige en renvoyant au tableau récapitulatif de la leçon.
\end{enumerate}
\end{experience}

\begin{savaistu}[Le savais-tu ? — Variété connecteurs = points bonus]
À l'épreuve de traduction du Bac A, les examinateurs valorisent la \textbf{variété des connecteurs}. Un candidat qui alterne \esp{porque}, \esp{ya que}, \esp{dado que}, \esp{puesto que}, \esp{como} selon la nuance (cause objective, cause connue, cause insistée) obtient une meilleure note qu'un candidat qui répète \esp{porque} cinq fois. De même, varier \esp{para que}, \esp{a fin de que}, \esp{con el fin de que} montre une maîtrise lexicale supérieure. \textbf{Astuce :} préparez une « fiche connecteurs » avec 3 colonnes (Français / Espagnol courant / Espagnol soutenu) et révisez-la la veille de l'examen.
\end{savaistu}

\courssub{Bilan : les trois règles d'or de la leçon}

\begin{aretenir}[RÈGLES D'OR — À apprendre par cœur]
\begin{enumerate}
\item \textbf{Cause vs Finalité :} \cle{Connecteur + Mode} — Cause réelle $\to$ \esp{porque / como / ya que / puesto que / dado que} + \textbf{INDICATIF} ; Finalité même sujet $\to$ \esp{para} + \textbf{INFINITIF} ; Finalité sujets différents $\to$ \esp{para que / a fin de que} + \textbf{SUBJONCTIF PRÉSENT}.
\item \textbf{Futur dans la subordonnée :} \cle{Quand / Dès que / Jusqu'à ce que} + action future $\to$ \esp{cuando / en cuanto / tan pronto como / hasta que} + \textbf{SUBJONCTIF PRÉSENT} (valeur futur antérieur). Le principal est au \textbf{FUTUR} ou à l'\textbf{IMPÉRATIF}. \textbf{Attention :} \esp{Si} + indicatif présent (réel) / imparfait subjonctif (irréel).
\item \textbf{Emphase :} \cle{Structure cleft} — \emph{Ce qui compte, c'est...} $\to$ \esp{Lo que importa es...} / \esp{Lo importante es...} / \esp{Lo fundamental es...} + \textbf{INFINITIF} ou \textbf{NOM}. Ne traduisez jamais \emph{ce qui} par \esp{*lo que} sans verbe \esp{ser} derrière.
\end{enumerate}
\end{aretenir}

\begin{center}
\textbf{Fin de la leçon E21 — Module 4 : Activités économiques et monde du travail}
\end{center}

\newpage

\coursec{Activité d'intégration}

\coursec{Leçon E21 : Traduction — Cause, finalité, futur dans la subordonnée, emphase (révision)}

\courssub{Observation guidée : Repérage des structures cibles}
\begin{textebox}[Correo de Carlos a Amadou -- Contexto previo]
Querido Amadou:

Me alegra mucho saber que el \emph{Ngondo} fue todo un éxito. \textbf{Como} ya eres mayor de edad, \textbf{puedes decidir} tu futuro. \textbf{Dado que} las notas del primer trimestre son buenas, \textbf{tienes muchas opciones}. Estudio ADE \textbf{para} trabajar en una multinacional, \textbf{pero} mi hermana estudia Derecho \textbf{para que} yo pueda entender sus apuntes (broma). \textbf{Cuando termines} el bac, \textbf{cuéntamelo todo}. \textbf{Lo que más me interesa} es saber si vienes a Madrid. \textbf{No es la distancia lo que importa}, sino las ganas.

Un abrazo,

Carlos
\end{textebox}

\begin{experience}[Activité de repérage (10 min)]
\begin{enumerate}
    \item Souligne dans le texte les \textbf{connecteurs de cause} (\esp{como}, \esp{dado que}) et précise le registre de chacun.
    \item Identifie les \textbf{deux structures de finalité} : \esp{para} + infinitif (sujet identique ?) et \esp{para que} + subjonctif (sujet différent ?).
    \item Repère la \textbf{subordonnée temporelle future} : \esp{Cuando termines}. Quel mode ? Pourquoi ?
    \item Trouve les \textbf{deux structures d'emphase} : \esp{Lo que más me interesa es...} et \esp{No es... sino...}.
\end{enumerate}
\end{experience}

\courssub{Objectifs de l'activité}
\noindent Cette activité d'intégration vise à mobiliser, dans une production écrite authentique et personnelle, l'ensemble des notions grammaticales révisées au cours de la leçon E21 : l'expression de la \cle{cause} (connecteurs variés), de la \cle{finalité} (deux constructions distinctes), du \cle{futur dans la subordonnée temporelle} (subjonctif présent) et d'une structure d'\cle{emphase}. L'élève doit démontrer sa capacité à choisir la structure adaptée à son intention communicative, à contrôler la concordance des temps et des modes, et à s'auto-évaluer à l'aide d'une grille précise. Le contexte camerounais (projets post-bac, orientation, vie étudiante) ancre la tâche dans la réalité de l'apprenant.

\courssub{Situation}
\begin{textebox}[Contexte : Échange épistolaire avec un correspondant hispanophone]
Tu t'appelles \textbf{Amadou} (ou \textbf{Aïcha}), tu as 18 ans, tu es en Terminale A au \textbf{Lycée Général-Leclerc de Yaoundé}. Tu prépares le baccalauréat et tu réfléchis activement à ton avenir : études supérieures (droit, médecine, ingénierie, commerce), concours de la fonction publique, ou peut-être une année de césure pour travailler et voyager.

Depuis six mois, tu échanges des lettres et des messages avec \textbf{Carlos}, un étudiant madrilène de 20 ans qui étudie l'\emph{Administración y Dirección de Empresas} (ADE) à l'\emph{Universidad Complutense de Madrid}, ou avec \textbf{Ximena}, une jeune Mexicaine de ton âge qui termine son \emph{preparatoria} à \emph{Guadalajara} et rêve de devenir architecte. Votre dernier échange portait sur les fêtes de fin d'année au Cameroun (le \emph{Ngondo} à Douala, les danses \emph{bendskin} à Bafoussam) et sur la \emph{Navidad} en Espagne (\emph{Los Reyes Magos}, le \emph{turrón}) ou au Mexique (\emph{Las Posadas}, la \emph{rosca de reyes}).

Aujourd'hui, tu décides d'écrire une lettre plus longue et plus personnelle pour leur expliquer \textbf{tes projets concrets après le bac}, les raisons qui guident ton choix (\textbf{cause}), tes objectifs (\textbf{finalité}), les étapes qui t'attendent (\textbf{futur dans la subordonnée}) et ce qui compte vraiment pour toi (\textbf{emphase}). Tu veux que ton correspondant comprenne bien ta situation, ton environnement et tes motivations.
\end{textebox}

\courssub{Consignes}
\begin{enumerate}
    \item \textbf{Rédaction individuelle (30 min).} Rédige une lettre de \textbf{10 à 12 lignes} (environ 120--150 mots) adressée à Carlos ou à Ximena. Tu dois \textbf{impérativement} y intégrer les structures suivantes (coche chaque case une fois réalisée) :
    \begin{itemize}
        \item [☐] Au moins \textbf{deux connecteurs de cause différents}, dont obligatoirement \esp{como} \textbf{ou} \esp{dado que} (l'autre peut être \esp{porque}, \esp{puesto que}, \esp{ya que}, \esp{a causa de} + nom, \esp{debido a} + nom).
        \item [☐] Une finalité exprimée par \esp{para} + \textbf{infinitif}.
        \item [☐] Une finalité exprimée par \esp{para que} + \textbf{subjonctif présent} (sujet différent du principal).
        \item [☐] \textbf{Deux subordonnées temporelles} faisant référence au futur** (introduites par \esp{cuando}, \esp{en cuanto}, \esp{hasta que}, \esp{tan pronto como}) avec le verbe au \textbf{subjonctif présent}.
        \item [☐] \textbf{Une tournure emphatique} au choix : \esp{Lo importante es que} + subjonctif, \esp{Es... lo que}..., \esp{Lo que}... \esp{es}..., \esp{No es... sino...}, ou cleft sentence avec pronom tonique (\esp{A mí me gusta...}).
    \end{itemize}
    \item \textbf{Échange par binômes (15 min).} Échange ta lettre avec ton voisin. Avec un stylo de couleur, \textbf{souligne} dans sa lettre :
    \begin{itemize}
        \item les connecteurs de cause (trait continu \textcolor{popblue}{\rule{1cm}{0.5pt}}),
        \item les structures de finalité (trait pointillé \textcolor{poporange}{\rule{0.3cm}{0.5pt}} \textcolor{poporange}{\rule{0.3cm}{0.5pt}}),
        \item les subordonnées temporelles au futur (vague \textcolor{popgreen}{\rule{0.3cm}{0.5pt}} \textcolor{popgreen}{\rule{0.3cm}{0.5pt}} \textcolor{popgreen}{\rule{0.3cm}{0.5pt}}),
        \item la structure d'emphase (double trait \textcolor{poppurple}{\rule{1cm}{0.5pt}} \textcolor{poppurple}{\rule{1cm}{0.5pt}}).
    \end{itemize}
    \item \textbf{Auto-évaluation croisée (10 min).} Remplis la \textbf{grille d'auto-évaluation} ci-dessous pour la lettre de ton camarade. Coche \emph{Oui} / \emph{Non} / \emph{Partiel} et note une observation si besoin.
    \begin{center}
    \begin{tabularx}{\linewidth}{|X|c|c|c|X|}
    \hline
    \rowcolor{popblueL} \textbf{Critère} & \textbf{Oui} & \textbf{Non} & \textbf{Partiel} & \textbf{Observation / Correction proposée} \\
    \hline
    2 connecteurs de cause variés (dont \esp{como} ou \esp{dado que}) & & & & \\
    \hline
    \esp{para} + infinitif (sujet identique) & & & & \\
    \hline
    \esp{para que} + subjonctif (sujet différent) & & & & \\
    \hline
    2 temporelles futur : \esp{cuando/en cuanto/hasta que} + subjonctif & & & & \\
    \hline
    1 structure d'emphase correcte & & & & \\
    \hline
    Orthographe, accents, ponctuation (\esp{¿ ¡}) corrects & & & & \\
    \hline
    Cohérence globale, ton approprié (lettre amicale) & & & & \\
    \hline
    \end{tabularx}
    \end{center}
    \item \textbf{Réécriture et affichage (10 min).} Intègre les corrections de ton camarade. Les trois meilleures lettres (vote de la classe) sont lues à voix haute par leurs auteurs et affichées au tableau d'affichage \emph{« Rincón hispano »} de la classe.
\end{enumerate}

\courssub{Ressources à mobiliser}
\begin{definition}[Rappel : Connecteurs de cause -- Niveaux de langue]
\begin{tabularx}{\linewidth}{|l|X|X|}
\hline
\rowcolor{popblueL} \textbf{Connecteur} & \textbf{Construction} & \textbf{Niveau / Registre} \\
\hline
\esp{porque} & + indicatif (cause nouvelle, explication) & Neutre, courant \\
\hline
\esp{como} & + indicatif (cause connue, évidente, en tête de phrase) & Courant, explicatif \\
\hline
\esp{puesto que} & + indicatif (cause admise, argumentative) & Soutenu, écrit \\
\hline
\esp{ya que} & + indicatif (cause justificative, « puisque ») & Courant, argumentatif \\
\hline
\esp{dado que} & + indicatif (cause factuelle, « étant donné que ») & Formel, administratif \\
\hline
\esp{a causa de} / \esp{debido a} & + \textbf{nom / GN} (pas de verbe) & Formel, écrit \\
\hline
\end{tabularx}
\emph{Attention :} \esp{porque} (un mot) = « parce que » (conjonction) ; \esp{por qué} (deux mots, accent) = « pourquoi » (question) ; \esp{porqué} (nom masculin) = « le pourquoi, la raison ».
\end{definition}

\begin{propriete}[Finalité : \esp{para} + infinitif vs \esp{para que} + subjonctif]
\begin{itemize}
    \item \textbf{\esp{Para} + infinitif} : sujet de la finalité \textbf{identique} au sujet du principal. \esp{Estudio para ser médico.} (C'est moi qui étudie et c'est moi qui veux être médecin.)
    \item \textbf{\esp{Para que} + subjonctif présent} : sujet de la finalité \textbf{différent} du sujet du principal. \esp{Estudio para que mis padres estén orgullosos.} (J'étudie ; \emph{eux} sont fiers.)
    \item Autres connecteurs de finalité (plus formels) : \esp{a fin de que}, \esp{con el fin de que}, \esp{con objeto de que} + \textbf{subjonctif}. \esp{Con el fin de} + infinitif (sujet identique) est possible mais lourd.
\end{itemize}
\end{propriete}

\begin{propriete}[Futur dans la subordonnée temporelle : Subjonctif présent]
\begin{center}
\begin{tabularx}{\linewidth}{|l|X|X|}
\hline
\rowcolor{popblueL} \textbf{Conjonction} & \textbf{Règle} & \textbf{Exemple} \\
\hline
\esp{cuando}, \esp{en cuanto}, \esp{tan pronto como}, \esp{hasta que}, \esp{después de que}, \esp{antes de que} & Si l'action est \textbf{future / hypothétique / non réalisée} par rapport au moment de l'énonciation $\rightarrow$ \textbf{Subjonctif présent}. & \esp{Te llamaré \textbf{cuando llegue} a Madrid.} (Pas encore arrivé.) \\
\hline
\emph{Mêmes conjonctions} & Si l'action est \textbf{habituelle / passée / réalisée} $\rightarrow$ \textbf{Indicatif} (présent, passé composé, imparfait...). & \esp{Siempre como \textbf{cuando llego} a casa.} (Habitude.) \\
\hline
\end{tabularx}
\end{center}
\emph{Piège fréquent :} \esp{después de que} et \esp{antes de que} exigent \textbf{toujours} le subjonctif (même au passé : \esp{antes de que llegara} -- imparfait subjonctif). Pour le futur : \esp{antes de que} + subj. présent.
\end{propriete}

\begin{definition}[Structures d'emphase -- Rappel synthétique]
\begin{enumerate}
    \item \textbf{\esp{Lo + adj/adv + es que} + indicatif/subjonctif} : \esp{Lo importante es que estudies.} / \esp{Lo mejor es que vengas.}
    \item \textbf{\esp{Es} + élément mis en relief + \esp{lo que} + verbe} : \esp{Es el derecho \textbf{lo que} quiero estudiar.} (C'est le droit que je veux étudier.) \esp{Fue ayer \textbf{lo que} \textbf{recibí} la carta.}
    \item \textbf{\esp{Lo que} + verbe + \esp{es} + élément} : \esp{Lo que quiero es estudiar derecho.}
    \item \textbf{\esp{No es} A \esp{sino} B} : \esp{No es el dinero \textbf{sino} la vocación \textbf{lo que} importa.}
    \item \textbf{Pronoms toniques} : \esp{A mí me preocupa el futuro.} (À moi, moi, le futur m'inquiète.)
\end{enumerate}
\end{definition}

\begin{textebox}[Lexique utile : Orientation, études, projets (Cameroun / Hispanophone)]
\begin{description}
    \item[Les filières] \esp{Derecho} (Droit), \esp{Medicina} (Médecine), \esp{Ingeniería} (Ingénierie), \esp{Administración de Empresas / ADE} (Gestion), \esp{Arquitectura} (Architecture), \esp{Enfermería} (Soins infirmiers), \esp{Periodismo} (Journalisme), \esp{Informática} (Informatique).
    \item[Les examens / concours] \esp{El bachillerato} (Bac), \esp{la Selectividad / PAU / EBAU} (Concours entrée fac Espagne), \esp{el concurso} (Concours fonction publique Cameroun), \esp{la prueba de acceso} (Épreuve d'admission).
    \item[Verbes clés] \esp{matricularse en} (s'inscrire à), \esp{solicitar una beca} (demander une bourse), \esp{presentarse a} (passer un concours), \esp{optar por} (opter pour), \esp{decantarse por} (se décider pour), \esp{formarse} (se former), \esp{emprender} (entreprendre).
    \item[Expressions de cause / finalité] \esp{Por motivos económicos}, \esp{por vocación}, \esp{dado que mi familia...}, \esp{como no hay plazas...}, \esp{con el fin de que...}, \esp{a fin de poder...}.
    \item[Marqueurs de lettre amicale] \esp{Querido Carlos / Querida Ximena:}, \esp{Te escribo para contarte...}, \esp{Espero tus noticias.}, \esp{Un abrazo fuerte,}, \esp{Tu amigo, Amadou}.
\end{description}
\end{textebox}

\courssub{Proposition de production et corrigé commenté}
\begin{exoresolu}[Lettre modèle : Amadou (Yaoundé) → Carlos (Madrid)]
\textbf{Énoncé :} Rédige la lettre d'Amadou à Carlos en respectant toutes les contraintes grammaticales imposées.

\tcblower

\textbf{Production modèle (espagnol) :}

\begin{flushleft}
Querido Carlos:

Te escribo desde Yaoundé, donde el calor de marzo nos recuerda que el final del curso está cerca. \textbf{Como} ya sabes, estoy en Terminale A en el Lycée Leclerc y \textbf{dado que} las notas del primer trimestre fueron buenas, me siento optimista. \textbf{Lo que más me ilusiona es} estudiar Derecho en la Universidad de Yaoundé II-Soa, \textbf{para} defender luego los derechos de los trabajadores del cacao. \textbf{Para que} mi familia no tenga que pagar todo, \textbf{cuando} \textbf{obtenga} el bac, \textbf{solicitaré} una beca del Ministerio de Educación Superior. \textbf{En cuanto} \textbf{tenga} la respuesta, te la comunicaré inmediatamente. No es el dinero \textbf{sino} la vocación \textbf{lo que} me impulsa.

Un abrazo fuerte,

Amadou
\end{flushleft}

\textbf{Traduction française (version) :}

Cher Carlos,

Je t'écris depuis Yaoundé, où la chaleur de mars nous rappelle que la fin de l'année scolaire approche. \textbf{Comme} tu le sais déjà, je suis en Terminale A au Lycée Leclerc et \textbf{étant donné que} les notes du premier trimestre étaient bonnes, je me sens optimiste. \textbf{Ce qui m'enthousiasme le plus, c'est} d'étudier le Droit à l'Université de Yaoundé II-Soa, \textbf{afin de} défendre ensuite les droits des travailleurs du cacao. \textbf{Afin que} ma famille n'ait pas à tout payer, \textbf{quand} j'\textbf{aurai obtenu} le bac, je \textbf{demanderai} une bourse au Ministère de l'Enseignement Supérieur. \textbf{Dès que} j'\textbf{aurai} la réponse, je te la communiquerai immédiatement. Ce n'est pas l'argent \textbf{mais} la vocation \textbf{qui} me pousse.

Une forte accolade,

Amadou

\textbf{Commentaire linguistique détaillé (corrigé commenté) :}

\begin{enumerate}
    \item \textbf{\esp{Como ya sabes...}} : \esp{Como} en tête de phrase, cause connue/évidente (le correspondant sait qu'Amadou est en Terminale). Indicatif \esp{sabes} (fait réel, présent). \emph{Alternative :} \esp{Puesto que ya sabes...} (plus formel).
    \item \textbf{\esp{dado que las notas... fueron buenas}} : \esp{dado que} + indicatif \esp{fueron} (prétérit indéfini, fait passé accompli). Registre formel/administratif, adapté à une justification scolaire. \emph{Respecte la contrainte :} \esp{como} ET \esp{dado que} utilisés.
    \item \textbf{\esp{Lo que más me ilusiona es estudiar...}} : Emphase de type \textbf{\esp{Lo que} + verbe + \esp{es} + infinitif/nom}. Met en relief l'objet de l'enthousiasme. \esp{Ilusiona} (verbe affectif, « enthousiasmer, faire rêver ») + infinitif \esp{estudiar}. \emph{Variante :} \esp{Es el Derecho lo que quiero estudiar} (cleft sentence sur le nom).
    \item \textbf{\esp{para defender luego los derechos...}} : \textbf{\esp{Para} + infinitif}. Sujet identique : Amadou étudie \textbf{et} Amadou défend. Finalité professionnelle. \emph{Attention :} pas de \esp{para que} ici (même sujet).
    \item \textbf{\esp{Para que mi familia no tenga que pagar todo}} : \textbf{\esp{Para que} + subjonctif présent \esp{tenga}}. Sujet \textbf{différent} : Amadou sollicite la bourse (principal) ; \emph{sa famille} ne paie pas (finalité). \esp{Tenga} = subj. présent 3e pers. sg de \esp{tener} (irrégulier : \esp{teng-}). \esp{Que pagar} = infinitif après \esp{tener que} (obligation).
    \item \textbf{\esp{cuando obtenga el bac}} : \textbf{Temporelle futur} \esp{cuando} + \textbf{subjonctif présent \esp{obtenga}} (verbe \esp{obtener}, irrégulier \esp{obteng-}). L'obtention du bac est future, non réalisée au moment de l'énonciation. \emph{Erreur fréquente :} \esp{cuando obtendré} (futur indicatif) $\leftarrow$ \textbf{INCORRECT} en espagnol standard.
    \item \textbf{\esp{solicitaré una beca...}} : Principal au \textbf{futur simple indicatif} (action future certaine/volontaire). \esp{Solicitaré} = 1e pers. sg futur de \esp{solicitar}.
    \item \textbf{\esp{En cuanto tenga la respuesta}} : \textbf{2e temporelle futur} \esp{en cuanto} + \textbf{subjonctif présent \esp{tenga}} (verbe \esp{tener}). « Dès que j'aurai la réponse. » \emph{Respecte la contrainte :} deux subordonnées temporelles futures au subjonctif (\esp{cuando obtenga}, \esp{en cuanto tenga}).
    \item \textbf{\esp{te la comunicaré inmediatamente}} : Futur indicatif \esp{comunicaré} + pronoms enclitiques \esp{te la} (te = à toi, la = la réponse). Ordre : \esp{te la comunicaré} (pas \esp{la te}).
    \item \textbf{\esp{No es el dinero sino la vocación lo que me impulsa}} : \textbf{Double emphase} : structure \esp{No es A sino B} (correction/restriction) + \esp{lo que} (relatif neutre antéposé). \esp{Impulsa} = indicatif présent (fait réel, vérité générale pour le locuteur). \emph{Variante possible :} \esp{Lo importante no es el dinero, sino la vocación.}
\end{enumerate}

\textbf{Vérification des contraintes (grille) :}
\begin{center}
\begin{tabularx}{\linewidth}{|X|c|X|}
\hline
\rowcolor{popblueL} \textbf{Contrainte} & \textbf{OK ?} & \textbf{Occurrence dans le texte} \\
\hline
2 connecteurs cause (dont \esp{como} / \esp{dado que}) & \textcolor{popgreen}{OUI} & \esp{Como ya sabes} ; \esp{dado que las notas...} \\
\hline
\esp{para} + infinitif & \textcolor{popgreen}{OUI} & \esp{para defender luego...} \\
\hline
\esp{para que} + subjonctif (sujet diff.) & \textcolor{popgreen}{OUI} & \esp{Para que mi familia no tenga...} \\
\hline
2 temporelles futur + subjonctif & \textcolor{popgreen}{OUI} & \esp{cuando obtenga} ; \esp{en cuanto tenga} \\
\hline
1 emphase & \textcolor{popgreen}{OUI} & \esp{Lo que más me ilusiona es...} ET \esp{No es el dinero sino la vocación lo que...} (2 !) \\
\hline
Longueur (10-12 lignes / 120-150 mots) & \textcolor{popgreen}{OUI} & $\approx$ 135 mots \\
\hline
Ton lettre amicale, vocabulaire contextualisé & \textcolor{popgreen}{OUI} & \esp{Querido Carlos}, \esp{Un abrazo fuerte}, \esp{Yaoundé}, \esp{Soa}, \esp{cacao}, \esp{becario}... \\
\hline
\end{tabularx}
\end{center}
\end{exoresolu}

\courssub{Bilan de l'activité}
\begin{aretenir}[Ce qu'il faut retenir de cette intégration]
\begin{itemize}
    \item La \textbf{cause} se nuance selon le connecteur : \esp{como}/\esp{dado que} (cause connue, formel) vs \esp{porque} (cause nouvelle) vs \esp{a causa de}/\esp{debido a} + nom (style administratif). Le choix trahit le registre de langue.
    \item La \textbf{finalité} impose un choix binaire strict : \textbf{même sujet} $\rightarrow$ \esp{para} + \textbf{infinitif} ; \textbf{sujet différent} $\rightarrow$ \esp{para que} + \textbf{subjonctif}. Aucune exception.
    \item Le \textbf{futur dans la subordonnée temporelle} est \textbf{le marqueur majeur du niveau Terminale} : \esp{cuando/en cuanto/hasta que} + \textbf{subjonctif présent} = action future non réalisée. L'indicatif futur (\esp{cuando obtendré}) est une erreur grave (calque du français « quand j'obtiendrai »).
    \item L'\textbf{emphase} (\esp{Lo que... es}, \esp{Es... lo que}, \esp{No es... sino}) permet de structurer l'information, de hiérarchiser les arguments et d'adopter un ton personnel, engagé -- essentiel dans une lettre de motivation ou un essai argumentatif.
    \item L'\textbf{auto-évaluation par grille} développe la métacognition : l'élève devient correcteur, verbalise la règle, repère l'erreur (songe à l'accord \esp{tenga} vs \esp{tiene}, au choix \esp{para} vs \esp{para que}).
\end{itemize}
\end{aretenir}

\begin{experience}[Prolongements possibles -- Différenciation]
\begin{itemize}
    \item \textbf{Niveau soutenu / Spécialité :} Transformer la lettre en \textbf{lettre de motivation formelle} pour une bourse d'études (Convocatoria MAEC-AECID, Becas Fundación Carolina, Erasmus+). Contraintes ajoutées : \esp{usted}, vocabulaire administratif (\esp{por la presente solicito}, \esp{adjunto certificados}), emphases multiples, pas de contractions.
    \item \textbf{Niveau renforcé / Remédiation :} Fournir une \textbf{lettre à trous} (texte à clozes) où ne manquent que les verbes (choix mode/temps) et les connecteurs. L'élève ne gère que la grammaire ciblée.
    \item \textbf{Expression orale :} Enregistrement d'un \textbf{message vocal WhatsApp} (1 min 30) reprenant le contenu de la lettre. Évaluation sur la prononciation du subjonctif (\esp{obtenga, tenga, solicite}) et l'intonation emphatique (\esp{¡Lo que quiero ES estudiar derecho!}).
    \item \textbf{Culture / Guinée équatoriale :} Variante contextuelle : l'élève écrit à un correspondant de \emph{Malabo} ou \emph{Bata}. Lexique adapté : \esp{selectividad} $\rightarrow$ \esp{pruebas de acceso a la universidad}, \esp{bachillerato} $\rightarrow$ \esp{bachillerato unificado}, mention du \emph{français} comme langue officielle partagée.
\end{itemize}
\end{experience}

\begin{savaistu}[Le saviez-vous ? -- L'orientation post-bac en Espagne et au Mexique]
\begin{itemize}
    \item En \textbf{Espagne}, l'accès à l'université publique se fait via l'\textbf{EBAU} (\emph{Evaluación de Bachillerato para el Acceso a la Universidad}), anciennement \emph{Selectividad} ou \emph{PAU}. La note (sur 14) combine notes du bac (60\%) et épreuves spécifiques (40\%). Les filières les plus demandées (\emph{carreras con nota de corte alta}) sont Médecine, Doubles diplômes (Droit+ADE), Ingénierie aérospatiale.
    \item Au \textbf{Mexique}, l'\textbf{EXANI-II} (Examen Nacional de Ingreso) du CENEVAL est le test standard pour entrer dans les universités publiques (UNAM, IPN, UAM). La concurrence est extrême : à l'UNAM, seulement 10\% des candidats sont admis en Médecine ou Droit.
    \item Au \textbf{Cameroun}, l'orientation se fait via \textbf{MINESUP} (Ministère de l'Enseignement Supérieur) après le Bac. Les bourses \emph{État} (excellence, mérite, social) et les accords bilatéraux (Espagne, France, Maroc, Chine, Russie) sont les voies principales pour partir étudier à l'étranger. Le \emph{Guide de l'étudiant camerounais à l'étranger} (MINESUP) est une ressource officielle à consulter.
    \item \textbf{Guinée équatoriale} : Système calqué sur l'Espagne (Bachillerato $\rightarrow$ Selectividad/Acceso). L'\emph{Universidad Nacional de Guinea Ecuatorial (UNGE)} à Malabo et Bata est le pôle principal. Beaucoup d'étudiants poursuivent en Espagne grâce à des bourses du \emph{Gobierno de Guinea Ecuatorial} ou de l'\emph{AECID}.
\end{itemize}
\end{savaistu}

\newpage

\coursec{Méthodes et exercices résolus}

\courssub{Méthodes et exercices résolus}

\begin{popfigure}
\begin{center}\tcbox[colback=popink,colframe=popink,coltext=white,arc=9pt,boxsep=4pt,left=6pt,right=6pt]{\bfseries\small Leçon E21 Traduction Avancée}\end{center}
\vspace{-2pt}
\begin{tcbraster}[raster columns=2,raster equal height=rows,raster column skip=6pt,raster row skip=6pt,enhanced,blankest]
\begin{tcolorbox}[enhanced,colback=white,colframe=poporange,boxrule=0.9pt,arc=3pt,title={Cause},fonttitle=\bfseries\scriptsize,coltitle=white,colbacktitle=poporange,left=4pt,right=4pt,top=2pt,bottom=2pt,boxsep=1pt,toptitle=1.5pt,bottomtitle=1.5pt,halign=left]
\begin{itemize}[nosep,leftmargin=*,itemsep=1pt,font=\scriptsize]
\item porque (fin, neutre)
\item como, ya que, puesto que (début, connu)
\item dado que (formel)
\item debido a / a causa de + GN
\end{itemize}
\end{tcolorbox}
\begin{tcolorbox}[enhanced,colback=white,colframe=popgreen,boxrule=0.9pt,arc=3pt,title={Finalité},fonttitle=\bfseries\scriptsize,coltitle=white,colbacktitle=popgreen,left=4pt,right=4pt,top=2pt,bottom=2pt,boxsep=1pt,toptitle=1.5pt,bottomtitle=1.5pt,halign=left]
\begin{itemize}[nosep,leftmargin=*,itemsep=1pt,font=\scriptsize]
\item para + Inf. (même sujet)
\item para que + Subj (sujets diff.)
\item a fin de que + Subj (formel)
\end{itemize}
\end{tcolorbox}
\begin{tcolorbox}[enhanced,colback=white,colframe=poppurple,boxrule=0.9pt,arc=3pt,title={Futur Subordonné},fonttitle=\bfseries\scriptsize,coltitle=white,colbacktitle=poppurple,left=4pt,right=4pt,top=2pt,bottom=2pt,boxsep=1pt,toptitle=1.5pt,bottomtitle=1.5pt,halign=left]
\begin{itemize}[nosep,leftmargin=*,itemsep=1pt,font=\scriptsize]
\item Cuando / En cuanto / Hasta que + Subj Présent
\item Si + Indicatif Présent (condition)
\end{itemize}
\end{tcolorbox}
\begin{tcolorbox}[enhanced,colback=white,colframe=poppink,boxrule=0.9pt,arc=3pt,title={Emphase},fonttitle=\bfseries\scriptsize,coltitle=white,colbacktitle=poppink,left=4pt,right=4pt,top=2pt,bottom=2pt,boxsep=1pt,toptitle=1.5pt,bottomtitle=1.5pt,halign=left]
\begin{itemize}[nosep,leftmargin=*,itemsep=1pt,font=\scriptsize]
\item Es... quien/que (Cleft)
\item Lo que... es...
\item Topicalisation (Début phrase)
\end{itemize}
\end{tcolorbox}
\end{tcbraster}

\legende{Carte mentale récapitulative des 4 piliers grammaticaux de la leçon E21}
\end{popfigure}

\begin{methode}[Traduire la cause : connecteurs et registres]
\textbf{Objectif :} Choisir le bon connecteur causal selon le registre (oral/écrit), la place dans la phrase et la nature du complément (verbe ou nom).\\
\textbf{Étapes :}
\begin{enumerate}
    \item \textbf{Identifier le segment causal} dans la phrase française (introduit par \emph{parce que}, \emph{comme}, \emph{puisque}, \emph{vu que}, \emph{à cause de}, \emph{grâce à}, \emph{du fait de}, \emph{compte tenu de}).
    \item \textbf{Analyser la structure syntaxique :}
        \begin{itemize}
            \item Suivi d'un \textbf{verbe conjugué} (proposition subordonnée) $\rightarrow$ conjonction : \esp{porque}, \esp{como}, \esp{ya que}, \esp{puesto que}, \esp{dado que}.
            \item Suivi d'un \textbf{nom / GN / infinitif} (groupe prépositionnel) $\rightarrow$ locution prépositive : \esp{a causa de}, \esp{debido a}, \esp{por}, \esp{gracias a}, \esp{dado} (participe prépositionnel).
        \end{itemize}
    \item \textbf{Choisir le registre et la place :}
        \begin{itemize}
            \item \esp{Porque} : neutre, réponse à \esp{¿por qué?}, place naturelle en \textbf{fin de phrase} (information nouvelle).
            \item \esp{Como}, \esp{ya que}, \esp{puesto que} : cause connue, évidente, souvent en \textbf{début de phrase} (thème), suivis d'une virgule.
            \item \esp{Dado que} : formel, écrit (administratif, journalistique, académique), souvent en tête.
            \item \esp{A causa de} / \esp{Debido a} + GN : formel. \esp{A causa de} souvent péjoratif (cause négative) ; \esp{debido a} neutre.
            \item \esp{Gracias a} + GN : cause positive / bénéfique.
            \item \esp{Por} + GN / Infinitif : cause efficiente, motif (souvent oral ou locutions figées : \esp{por eso}, \esp{por qué}).
        \end{itemize}
    \item \textbf{Gérer l'accord de \esp{debido a} :} Traditionnellement accordé (\esp{debidos a}, \esp{debida a}) comme participe, mais l'usage moderne standard (et la recommandation RAE / Bac) le traite comme locution prépositive \textbf{invariable} : \esp{debido a la lluvia}, \esp{debido a los problemas}. \cle{Recommandation Bac : laisser \esp{debido a} invariable.}
    \item \textbf{Traduire} en respectant la flexibilité de l'ordre des mots en espagnol (cause en tête très fréquente à l'écrit).
\end{enumerate}
\end{methode}

\begin{exoresolu}[Cause : version et thème ciblés]
\textbf{Énoncé :} Traduire en espagnol les phrases suivantes. Justifier le choix du connecteur.
\begin{enumerate}
    \item Comme il pleut, le match de bendskin est annulé.
    \item Les prix du cacao ont baissé à cause de la sécheresse.
    \item Je ne suis pas venu parce que j'étais malade (et non pour une autre raison).
    \item Compte tenu de l'inflation, le gouvernement a augmenté les salaires. (Registre formel)
    \item Grâce à tes conseils, j'ai réussi l'examen.
\end{enumerate}
\tcblower
\textbf{Solution détaillée :}
\begin{enumerate}
    \item \esp{Como llueve, el partido de bendskin está cancelado.} \\
    \emph{Justification :} Cause connue/évidente (\emph{comme} en tête) $\rightarrow$ \esp{Como} + indicatif. Virgule obligatoire après la subordonnée antéposée.
    \item \esp{Los precios del cacao han bajado \textbf{debido a} la sequía.} (ou \esp{a causa de la sequía}). \\
    \emph{Justification :} Cause nominale (\emph{à cause de} + GN) $\rightarrow$ locution prépositive. \esp{Debido a} neutre/formel préféré. \esp{Debido a} laissé invariable.
    \item \esp{No vine \textbf{porque} estaba enfermo.} \\
    \emph{Justification :} Cause nouvelle, information principale (focus sur la cause), réponse à \emph{¿por qué?} $\rightarrow$ \esp{porque} en fin de phrase. Indicatif (fait réel passé).
    \item \esp{\textbf{Dado el aumento de la inflación}, el gobierno ha subido los salarios.} (ou \esp{Dado que la inflación ha aumentado...}). \\
    \emph{Justification :} Registre formel/administratif (\emph{compte tenu de}). \esp{Dado} (participe prépositionnel) + GN ou \esp{Dado que} + indicatif.
    \item \esp{\textbf{Gracias a} tus consejos, he aprobado el examen.} \\
    \emph{Justification :} Cause positive/bénéfique (\emph{grâce à}) $\rightarrow$ \esp{Gracias a} + GN.
\end{enumerate}
\textbf{Conclusion :} La distinction verbe/nom et le registre (neutre vs formel vs positif) dictent le choix. L'ordre cause-effet (thème-rhème) est la norme en espagnol écrit formel.
\end{exoresolu}

\begin{methode}[Exprimer la finalité : \esp{para} vs \esp{para que} + subjonctif]
\textbf{Objectif :} Maîtriser l'alternance infinitif / subjonctif selon l'identité ou la différence des sujets.\\
\textbf{Étapes :}
\begin{enumerate}
    \item \textbf{Repérer les deux actions :} Action principale (moyen) $\rightarrow$ Action visée (but).
    \item \textbf{Comparer les sujets des deux actions :}
        \begin{itemize}
            \item \textbf{Même sujet} $\rightarrow$ \esp{Para} + \textbf{infinitif}. (Ex: \esp{Estudio \textbf{para aprobar}.})
            \item \textbf{Sujets différents} $\rightarrow$ \esp{Para que} + \textbf{subjonctif} (présent si principal présent/futur/impératif ; imparfait si principal passé/conditionnel).
        \end{itemize}
    \item \textbf{Variantes formelles (sujets différents) :} \esp{A fin de que}, \esp{con el fin de que}, \esp{con el objeto de que} + \textbf{subjonctif}. Souvent en début de phrase pour l'emphase.
    \item \textbf{Négation :} \esp{Para que no} + subjonctif (but négatif / prévention). \esp{Para no} + infinitif (même sujet).
    \item \textbf{Piège majeur :} \emph{Pour que} + indicatif n'existe pas en espagnol standard. Toujours subjonctif si sujets différents.
\end{enumerate}
\end{methode}

\begin{exoresolu}[Finalité : transformation et thème]
\textbf{Énoncé :}
\begin{enumerate}
    \item Transformer en utilisant \esp{para que} + subjonctif (sujets différents) :
    \begin{itemize}
        \item \esp{Los padres trabajan mucho. Quieren que sus hijos estudien en la universidad.}
        \item \esp{Te doy mi número. Quiero que me llames mañana.}
    \end{itemize}
    \item Traduire (Thème) :
    \begin{itemize}
        \item Nous économisons de l'argent pour acheter une maison à Yaoundé.
        \item Le professeur explique la leçon pour que les élèves comprennent.
        \item J'ai apporté les documents afin que tu puisses les signer. (Passé + capacité)
    \end{itemize}
\end{enumerate}
\tcblower
\textbf{Solution détaillée :}
\begin{enumerate}
    \item \textbf{Transformations :}
        \begin{itemize}
            \item \esp{Los padres trabajan mucho \textbf{para que sus hijos estudien} en la universidad.} (Sujets : \emph{padres} vs \emph{hijos} $\rightarrow$ Subjonctif présent \esp{estudien}).
            \item \esp{Te doy mi número \textbf{para que me llames} mañana.} (Sujets : \emph{yo} vs \emph{tú} $\rightarrow$ Subjonctif présent \esp{llames}).
        \end{itemize}
    \item \textbf{Thème :}
        \begin{itemize}
            \item \esp{Ahorramos dinero \textbf{para comprar} una casa en \textbf{Yaundé}.} (Même sujet \emph{nosotros} $\rightarrow$ Infinitif. \cle{Graphie :} \esp{Yaundé}).
            \item \esp{El profesor explica la lección \textbf{para que los alumnos entiendan}.} (Sujets diff. \emph{profesor} vs \emph{alumnos} $\rightarrow$ \esp{para que} + Subj. présent \esp{entiendan}).
            \item \esp{Traje los documentos \textbf{para que pudieras} firmarlos.} (Principal \esp{traje} = passé $\rightarrow$ Subjonctif \textbf{imparfait} \esp{pudieras} / \esp{puedieras}. \cle{Note :} \esp{puedas} (présent) accepté seulement si le but est perçu comme encore actuel au moment de l'énoncé ; \esp{pudieras} (imparfait) est la norme de concordance des temps pour un passé narré). \textbf{Choix Bac :} \esp{para que pudieras firmarlos}.
        \end{itemize}
\end{enumerate}
\textbf{Conclusion :} Le test du sujet est l'unique critère fiable. La concordance des temps (principal passé $\rightarrow$ subj. imparfait) est exigible au Bac.
\end{exoresolu}

\begin{methode}[Le futur dans la subordonnée : \esp{cuando}, \esp{en cuanto}, \esp{tan pronto como}, \esp{si} + Subjonctif Présent]
\textbf{Objectif :} Appliquer la règle « Futur dans la principale $\rightarrow$ Subjonctif présent dans la subordonnée temporelle/conditionnelle » et distinguer \esp{si} conditionnel.\\
\textbf{Étapes :}
\begin{enumerate}
    \item \textbf{Repérer le mot-outil introducteur de subordonnée :}
        \begin{itemize}
            \item Temporels (antériorité/simultanéité future) : \esp{Cuando}, \esp{En cuanto}, \esp{Tan pronto como}, \esp{Hasta que}, \esp{Después de que}, \esp{Antes de que}, \esp{Mientras} (valeur future).
            \item Conditionnel : \esp{Si}.
        \end{itemize}
    \item \textbf{Analyser le temps de la principale :}
        \begin{itemize}
            \item Principale au \textbf{Futur simple} (ou Périphrastique \esp{ir a} + inf., ou Impératif, ou Présent de valeur future) $\rightarrow$ Subordonnée temporelle au \textbf{Subjonctif Présent}.
            \item Principale au \textbf{Passé} (ou Conditionnel) $\rightarrow$ Subordonnée au \textbf{Subjonctif Imparfait} (si antériorité/simultanéité) ou \textbf{Plus-que-parfait} (si antériorité marquée).
        \end{itemize}
    \item \textbf{Conjuguer le verbe de la subordonnée au Subjonctif Présent} (réguliers : \esp{-e, -es, -e, -emos, -éis, -en} / \esp{-a, -as, -a, -amos, -áis, -an} ; irréguliers fréquents : \esp{sea, vaya, tenga, haga, salga, ponga, venga, sepa, haya}).
    \item \textbf{Distinction cruciale \esp{Cuando} interrogatif vs temporel :}
        \begin{itemize}
            \item \esp{¿Cuándo vienes?} (Interrogatif direct/indirect) $\rightarrow$ \textbf{Indicatif}.
            \item \esp{Cuando vengas, avisamos.} (Temporel futur) $\rightarrow$ \textbf{Subjonctif}.
        \end{itemize}
    \item \textbf{Distinction cruciale \esp{Si} conditionnel vs interrogatif indirect :}
        \begin{itemize}
            \item \esp{Si llueve, no voy.} (Condition future réelle/possible) $\rightarrow$ \textbf{Indicatif Présent} (Règle classique : \esp{Si} + Présent Indicatif $\rightarrow$ Futur Indicatif).
            \item \esp{No sé si vendrá.} (Interrogatif indirect) $\rightarrow$ \textbf{Indicatif} (généralement) ou Subjonctif si doute fort.
            \item \cle{Règle d'or :} La règle « Subjonctif après \esp{cuando} futur » ne s'applique \textbf{JAMAIS} à \esp{si} conditionnel. \esp{*Si llueva} est une faute grave.
        \end{itemize}
\end{enumerate}
\end{methode}

\begin{exoresolu}[Futur dans la subordonnée : conjugaison et traduction]
\textbf{Énoncé :} Mettre le verbe entre parenthèses au temps convenable (Indicatif ou Subjonctif). Traduire.
\begin{enumerate}
    \item Te llamaré en cuanto \_\_\_\_\_\_\_\_ (llegar) a Douala.
    \item \_\_\_\_\_\_\_\_ (Cuando / ¿Cuándo?) \_\_\_\_\_\_\_\_ (venir) tú, iremos al mercado.
    \item Si \_\_\_\_\_\_\_\_ (haber) tiempo, \_\_\_\_\_\_\_\_ (visitar) el museo mañana.
    \item No me iré hasta que \_\_\_\_\_\_\_\_ (terminar) el trabajo.
    \item Dime \_\_\_\_\_\_\_\_ (cuándo / si) \_\_\_\_\_\_\_\_ (poder) venir.
\end{enumerate}
\tcblower
\textbf{Solution détaillée :}
\begin{enumerate}
    \item \esp{Te llamaré en cuanto \textbf{llegues} a Douala.} \\
    \emph{Traduction :} Je t'appellerai dès que tu arriveras à Douala. \\
    \emph{Justification :} Principale \esp{llamaré} (Futur) + \esp{en cuanto} (antériorité immédiate) $\rightarrow$ \textbf{Subjonctif Présent} (\esp{llegues}, 2e pers. sg).
    \item \esp{\textbf{Cuando} \textbf{vengas} tú, iremos al mercado.} \\
    \emph{Traduction :} Quand tu viendras, nous irons au marché. \\
    \emph{Justification :} Pas de signes \esp{¿ ?} $\rightarrow$ Temporel (pas interrogatif). Futur \esp{iremos} $\rightarrow$ \textbf{Subjonctif Présent} (\esp{vengas}, irrégulier \esp{venir}).
    \item \esp{Si \textbf{hay} tiempo, \textbf{visitaremos} el museo mañana.} \\
    \emph{Traduction :} S'il y a du temps, nous visiterons le musée demain. \\
    \emph{Justification :} \esp{Si} conditionnel + Futur dans principale $\rightarrow$ \textbf{Indicatif Présent} (\esp{hay}, impersonnel). \cle{Piège majeur :} On ne met JAMAIS de subjonctif après \esp{si} conditionnel.
    \item \esp{No me iré hasta que \textbf{termine} el trabajo.} \\
    \emph{Traduction :} Je ne partirai pas tant que le travail ne sera pas fini. \\
    \emph{Justification :} \esp{No me iré} (Futur) + \esp{hasta que} (jusqu'à ce que) $\rightarrow$ \textbf{Subjonctif Présent} (\esp{termine}).
    \item \esp{Dime \textbf{cuándo} \textbf{podrás} venir.} (ou \esp{puedas} si doute fort). \\
    \emph{Traduction :} Dis-moi quand tu pourras venir. \\
    \emph{Justification :} Interrogatif indirect (\esp{Dime...}). Verbe \esp{poder} au \textbf{Futur Indicatif} (\esp{podrás}) car c'est une question sur un fait futur certain/prévisible. Subjonctif \esp{puedas} possible seulement si incertitude subjective forte.
\end{enumerate}
\textbf{Conclusion :} La triade \esp{Cuando/En cuanto/Hasta que} + Futur Principal = Subjonctif Présent. \esp{Si} conditionnel = Indicatif Présent. Interrogatif indirect = Indicatif (généralement).
\end{exoresolu}

\begin{methode}[Révision de l'emphase : structures cleft et topicalisation]
\textbf{Objectif :} Mettre un élément en relief (sujet, COD, CC, verbe) pour corriger, insister ou opposer.\\
\textbf{Étapes :}
\begin{enumerate}
    \item \textbf{Identifier l'élément à souligner} (le \emph{focus} informationnel).
    \item \textbf{Choisir la structure adaptée :}
        \begin{itemize}
            \item \textbf{Cleft sujet (Sujet) :} \esp{Es / Era / Fue / Será + [Sujet mis en relief] + quien/que (relative)}. \cle{Accord :} Le verbe \esp{ser} s'accorde avec le sujet mis en relief (\esp{Son mis padres...}). \esp{Quién} pour personnes, \esp{que} pour choses/animaux.
            \item \textbf{Cleft objet / CC (COD, CCL, CCM, CCT...) :} \esp{Es / Era... + [Élément mis en relief] + que + [Phrase complète avec reprise pronominale si COD]}. \esp{Que} invariable.
            \item \textbf{Topicalisation (Fronting) :} Déplacer l'élément en \textbf{début de phrase} + \textbf{virgule} (écrit) ou intonation montante (oral). Ordre VSO ou VOS fréquent. Ex: \esp{El libro, lo leyó Juan.} / \esp{Mañana voy yo.}
            \item \textbf{Structure \esp{Lo que} + verbe + \esp{es/era} + [Focus] :} \esp{Lo que quiero es dormir.} (Focus sur l'infinitif/nom/action).
            \item \textbf{Emphase verbale (action) :} \esp{Lo que hice fue + infinitif} / \esp{Lo que hace es + infinitif}.
        \end{itemize}
    \item \textbf{Vérifier les accords et pronoms :}
        \begin{itemize}
            \item \esp{Ser} accordé au focus (sujet cleft).
            \item Pronoms toniques (\esp{mí, ti, él, ella, nosotros, vosotros, ellos/ellas}) après préposition ou en cleft sujet (\esp{Es para mí}, \esp{Soy yo}).
            \item Reprise du COD par \esp{lo/la/los/las} dans la relative si cleft COD (\esp{Es el libro que \textbf{lo} leí}).
        \end{itemize}
    \item \textbf{Traduire le \emph{C'est... qui/que} français :}
        \begin{itemize}
            \item Sujet personne : \esp{Es... quien...}
            \item Sujet chose / Autres fonctions : \esp{Es... que...}
            \item \emph{C'est moi / C'est toi} $\rightarrow$ \esp{Soy yo} / \esp{Eres tú} (pas \esp{*Es yo}).
        \end{itemize}
\end{enumerate}
\end{methode}

\begin{exoresolu}[Emphase : réécriture et traduction]
\textbf{Énoncé :}
\begin{enumerate}
    \item Mettre en relief l'élément souligné en utilisant une phrase clivée (\esp{Es... que/quien}) :
    \begin{itemize}
        \item \textbf{María} ha comprado el coche.
        \item He comprado \textbf{el coche} (y no la moto).
        \item Lo he comprado \textbf{ayer}.
    \end{itemize}
    \item Traduire en espagnol (emphase) :
    \begin{itemize}
        \item C'est \textbf{Paul} qui a mangé le ndolé.
        \item C'est \textbf{hier} que je l'ai vu.
        \item Ce que je veux, \textbf{c'est dormir}.
    \end{itemize}
\end{enumerate}
\tcblower
\textbf{Solution détaillée :}
\begin{enumerate}
    \item \textbf{Clivage (Cleft sentences) :}
        \begin{itemize}
            \item \esp{\textbf{Es María quien} ha comprado el coche.} (Sujet $\rightarrow$ \esp{quien}, \esp{ser} sg car \esp{María} sg).
            \item \esp{\textbf{Es el coche que} he comprado (y no la moto).} (COD $\rightarrow$ \esp{que} invariable. Reprise \esp{lo} possible : \esp{que lo he comprado}).
            \item \esp{\textbf{Es ayer cuando} lo he comprado.} (CC Temps $\rightarrow$ \esp{cuando} précis pour le temps, ou \esp{que}). \cle{Variante récit :} \esp{Fue ayer cuando lo compré} (Prétérit).
        \end{itemize}
    \item \textbf{Thème (Traduction emphatique) :}
        \begin{itemize}
            \item \esp{\textbf{Es Pablo quien} ha comido el ndolé.} (\emph{C'est Paul qui...} $\rightarrow$ Sujet cleft personne $\rightarrow$ \esp{quien}). \cle{Alternative orale :} \esp{Pablo ha comido el ndolé} (Topicalisation + intonation).
            \item \esp{\textbf{Es ayer cuando} lo vi.} (ou \esp{Fue ayer cuando lo vi} si récit passé achevé). \emph{C'est hier que...} $\rightarrow$ CC Temps cleft. \esp{Vi} (Prétérit indéfini) car action ponctuelle datée.
            \item \esp{\textbf{Lo que quiero es dormir.}} (Structure \esp{Lo que} + verbe + \esp{ser} + Focus Infinitif). \cle{Variante :} \esp{Lo que quiero hacer es dormir.}
        \end{itemize}
\end{enumerate}
\textbf{Conclusion :} L'emphase espagnole utilise massivement le clivage (\esp{Es... que/quien}) et la topicalisation. L'accord de \esp{ser} avec le focus (sujet) et l'usage de \esp{quien} (personne sujet) sont des marqueurs de niveau Bac.
\end{exoresolu}

\begin{methode}[Stratégie intégrée de traduction : Thème / Version « Monde du travail »]
\textbf{Objectif :} Mobiliser cause, finalité, futur subordonné et emphase dans un paragraphe cohérent.\\
\textbf{Étapes (Méthode BAC) :}
\begin{enumerate}
    \item \textbf{Analyse du texte source :} Identifier les segments \emph{Cause} (pourquoi), \emph{But} (pour quoi faire), \emph{Temporel/Conditionnel futur} (quand/si), \emph{Focus} (insistance/opposition).
    \item \textbf{Découpage en Unités de Traduction (UT) :} Une phrase française $\neq$ toujours une phrase espagnole. Regrouper ou scinder selon la syntaxe espagnole (ex: cause en tête, emphase par cleft).
    \item \textbf{Transcodage grammatical par UT :}
        \begin{itemize}
            \item Cause $\rightarrow$ Choix connecteur (\esp{como/ya que} en tête, \esp{porque} en queue, \esp{debido a} + GN).
            \item But $\rightarrow$ Test sujets (Infinitif vs \esp{para que} + Subj).
            \item Futur sub. $\rightarrow$ \esp{Cuando/En cuanto/Si} + Subj. Présent (si Futur princ.) / Indicatif (si \esp{Si} cond.).
            \item Emphase $\rightarrow$ Structure cleft (\esp{Es... que/quien/donde}) ou fronting si marqueur \emph{C'est... qui} ou italique/gras.
        \end{itemize}
    \item \textbf{Rédaction de la cible (Espagnol) :} Assembler les UT. Vérifier la cohésion (connecteurs logiques, pronoms relatifs, répétitions évitées).
    \item \textbf{Relecture « Check-list Bac » :} Accords (participes, \esp{debido a}), Accents (\esp{qué/cuándo/sí} interrogatifs vs relatifs), Subjonctifs (finalité, futur, négation), \esp{Ser/Estar}, \esp{Por/Para}, Faux-amis.
\end{enumerate}
\end{methode}

\begin{exoresolu}[Traduction intégrée : Paragraphe « Emploi jeunes »]
\textbf{Énoncé :} Traduire le paragraphe suivant en espagnol.
« \emph{Comme le chômage des jeunes augmente au Cameroun, le gouvernement a lancé un programme pour que les jeunes trouvent du travail. Quand ils auront une formation, ils pourront créer leur entreprise. C'est surtout dans l'agriculture que se trouvent les opportunités. Si tu es jeune, ne te plains pas : forme-toi ! }»
\tcblower
\textbf{Solution détaillée :}
\textbf{Analyse segment par segment :}
\begin{enumerate}
    \item \emph{Comme le chômage des jeunes augmente au Cameroun} $\rightarrow$ Cause connue, en tête. \esp{Como el desempleo juvenil aumenta en Camerún,} (ou \esp{Ya que...} / \esp{Dado que...} formel). \cle{Vocab :} \esp{desempleo juvenil} / \esp{paro juvenil}.
    \item \emph{le gouvernement a lancé un programme} $\rightarrow$ Passé (Prétérit action ponctuelle achevée). \esp{el gobierno lanzó un programa.}
    \item \emph{pour que les jeunes trouvent du travail} $\rightarrow$ Finalité, sujets diff (\emph{gouvernement} vs \emph{jeunes}). \esp{para que los jóvenes encuentren trabajo.} (Subj. présent \esp{encuentren} car programme encore actif/but actuel ; \esp{encontraran} imparfait possible par concordance stricte). \textbf{Choix :} \esp{encuentren}.
    \item \emph{Quand ils auront une formation} $\rightarrow$ \esp{Cuando} + Futur antérieur français = \textbf{Subjonctif Présent} espagnol. \esp{Cuando tengan una formación} (ou \esp{reciban una formación}). \esp{Tener} irrégulier $\rightarrow$ \esp{tengan}.
    \item \emph{ils pourront créer leur entreprise} $\rightarrow$ Principale Futur. \esp{podrán crear su empresa} (ou \esp{mountar su negocio}). \cle{Distributif :} \esp{su empresa} (sg) = une entreprise par jeune.
    \item \emph{C'est surtout dans l'agriculture que se trouvent les opportunités} $\rightarrow$ Emphase sur CC Lieu \emph{dans l'agriculture}. Structure Cleft : \esp{Es sobre todo en la agricultura donde se encuentran las oportunidades.} \cle{Note :} \esp{Donde} pour lieu. \esp{Se encuentran} (réfléchi passif/impersonnel) accordé avec \esp{las oportunidades} (pluriel).
    \item \emph{Si tu es jeune, ne te plains pas : forme-toi !} $\rightarrow$ \esp{Si} conditionnel + Présent Indicatif (\esp{eres}). Impératif négatif \esp{no te quejes} (Subj. présent). Impératif affirmatif \esp{fórmate} (accent sur \emph{á} pronom enclitique).
\end{enumerate}
\textbf{Version finale (Espagnol) :}
\esp{Como el desempleo juvenil aumenta en Camerún, el gobierno lanzó un programa para que los jóvenes encuentren trabajo. Cuando tengan una formación, podrán crear su empresa. Es sobre todo en la agricultura donde se encuentran las oportunidades. Si eres joven, no te quejes: ¡fórmate!}

\textbf{Justifications clés :} \esp{Como} (cause tête), \esp{para que} + Subj (finalité sujets diff), \esp{Cuando tengan} (futur sub), \esp{Es... donde} (emphase lieu), \esp{Si eres} (cond. indicatif), \esp{no te quejes / fórmate} (impératifs), \esp{¡...!} (ponctuation).
\end{exoresolu}

\begin{aretenir}[Boîte à outils : Connecteurs et structures clés pour le Bac]
\begin{center}
\begin{tabularx}{\linewidth}{|X|X|X|}
\hline
\rowcolor{popblueL} \textbf{Fonction} & \textbf{Structure (Espagnol)} & \textbf{Points de vigilance} \\
\hline
\multicolumn{3}{|c|}{\cellcolor{poporangeL}\textbf{CAUSE}} \\
\hline
Verbe (connu, tête) & \esp{Como / Ya que / Puesto que} + Indicatif & Virgule après la subordonnée. Pas \esp{porque}. \\
\hline
Verbe (nouveau, queue) & \esp{Porque} + Indicatif & Réponse à \esp{¿por qué?}. \\
\hline
Verbe (formel, tête) & \esp{Dado que} + Indicatif & Écrit, administratif. \\
\hline
Nom / GN (neutre/formel) & \esp{Debido a} + GN (\textbf{invariable}) & Préféré à \esp{a causa de}. \\
\hline
Nom / GN (négatif) & \esp{A causa de} + GN & Souvent péjoratif. \\
\hline
Nom / GN (positif) & \esp{Gracias a} + GN & Bénéfice. \\
\hline
\multicolumn{3}{|c|}{\cellcolor{popgreenL}\textbf{FINALITÉ}} \\
\hline
Même sujet & \esp{Para} + Infinitif & \esp{Para no} + Inf. (négatif). \\
\hline
Sujets différents & \esp{Para que} + Subjonctif & Présent (princ. P/F) / Imparfait (princ. Passé/Cond). \\
\hline
Formel (suj. diff.) & \esp{A fin de que / Con el fin de que} + Subj & Souvent en tête de phrase. \\
\hline
\multicolumn{3}{|c|}{\cellcolor{poppurpleL}\textbf{FUTUR DANS SUB. (Temporel)}} \\
\hline
\esp{Cuando, En cuanto, Tan pronto como, Hasta que, Antes de que, Después de que} & + \textbf{Subjonctif Présent} & Si Principale = Futur / Impératif / \esp{Ir a} + Inf. \\
\hline
\esp{Si} (Condition) & + \textbf{Indicatif Présent} & \textbf{JAMAIS} de subjonctif. \esp{Si llueve...} \\
\hline
\esp{Cuando} (Interrogatif) & + \textbf{Indicatif} & \esp{¿Cuándo vienes?} / \esp{No sé cuándo viene.} \\
\hline
\multicolumn{3}{|c|}{\cellcolor{poppinkL}\textbf{EMPHASE}} \\
\hline
Sujet (Personne) & \esp{Es... quien...} & \esp{Ser} accordé au focus. \\
\hline
Sujet (Chose) / COD / CC & \esp{Es... que...} & \esp{Que} invariable. Reprise COD (\esp{lo}). \\
\hline
CC Lieu / Temps & \esp{Es... donde / cuando...} & Précision spatiale/temporelle. \\
\hline
Action / Infinitif & \esp{Lo que... es...} + Infinitif & \esp{Lo que quiero es salir.} \\
\hline
\end{tabularx}
\end{center}
\end{aretenir}

\begin{methode}[Auto-correction et relecture : Check-list « Zéro Faute » Bac]
\textbf{Objectif :} Systématiser la relecture en 5 passes ciblées pour éliminer les erreurs coûteuses.\\
\textbf{Les 5 passes de relecture (ordre impératif) :}
\begin{enumerate}
    \item \textbf{Passe 1 : Accents \& Ponctuation (Visuel pur).} Repérer \textbf{tous} les mots interrogatifs/exclamatifs (\esp{¿qué? quién? cuándo? cómo? dónde? por qué?}) $\rightarrow$ Accent obligatoire. Repérer \textbf{tous} les monosyllabes à accent diacritique (\esp{sí, tú, él, mí, té, dé, sé, más}). Vérifier \esp{¿ ? ¡ !} en début/fin de phrase.
    \item \textbf{Passe 2 : Subjonctifs « Déclencheurs » (Grammaire pure).} Surligner tous les \esp{que}. Vérifier à gauche : \esp{Para que, para que no, a fin de que, cuando (futur), en cuanto, hasta que, antes de que, sin que, a menos que, es importante que, quiero que, dudo que, no creo que, ojalá}. $\rightarrow$ Verbe = Subjonctif ? Temps correct (Présent vs Imparfait) ?
    \item \textbf{Passe 3 : Ser / Estar \& Por / Para (Paires fatales).} \emph{Identité/Caractère/Heure/Origine/Passive/Definition} $\rightarrow$ \esp{Ser}. \emph{État/Lieu/Action en cours (gérondif)/Résultat} $\rightarrow$ \esp{Estar}. \emph{Cause/Moyen/Durée/Mouvement/Échange/Prix} $\rightarrow$ \esp{Por}. \emph{But/Destinataire/Échéance/Opinion/Proportion} $\rightarrow$ \esp{Para}.
    \item \textbf{Passe 4 : Accords \& Faux-amis (Lexique/Morpho).} Participes passés (\esp{-ado/-ido} accordés si \esp{ser/estar} + participe attributif / voix passive). \esp{Debido a} (invariable recommandé). Faux-amis fréquents : \esp{Actualmente} (actuellement $\neq$ actuellement), \esp{Carpeta} (dossier $\neq$ tapis), \esp{Éxito} (succès $\neq$ sortie), \esp{Largo} (long $\neq$ large), \esp{Realidad} (réalité $\neq$ realité), \esp{Sensible} (sensible $\neq$ sensible), \esp{Recordar} (se souvenir $\neq$ enregistrer), \esp{Asistir} (assister à $\neq$ assister), \esp{Éxamen} (examen $\neq$ examen), \esp{Librería} (librairie $\neq$ bibliothèque).
    \item \textbf{Passe 5 : Cohérence \& Registre (Style).} Connecteurs logiques (\esp{por tanto, sin embargo, además, en consecuencia}). Répétitions évitées (pronoms relatifs \esp{lo que, el que, cual, cuyo}). Ordre des mots naturel (SVO ou VSO si emphase). Registre soutenu si sujet formel (\esp{dado que, a fin de que, debido a, en consecuencia}).
\end{enumerate}
\end{methode}

\begin{exoresolu}[Détection et correction d'erreurs « Type Bac »]
\textbf{Énoncé :} Corriger les 10 erreurs contenues dans ce paragraphe d'élève. Expliquer chaque correction.
\esp{« Porque el desempleo aumenta, el gobierno ha lanzado un programa para que los jóvenes encuentran trabajo. Cuando ellos tendrán formación, podran crear sus empresas. Es en la agricultura donde están las oportunidades. Si tú eres joven, no te quejas: formate ! »}

\tcblower
\textbf{Solution détaillée : Corrections et Explications}
\begin{enumerate}
    \item \esp{Porque} $\rightarrow$ \textbf{\esp{Como / Ya que / Dado que}}. \emph{Erreur :} \esp{Porque} ne s'utilise pas en tête de phrase pour une cause connue (réponse à \emph{¿por qué?}). \esp{Como} place la cause en thème (structure thème-rhème).
    \item \esp{encuentran} (Indicatif) $\rightarrow$ \textbf{\esp{encuentren}} (Subjonctif Présent). \emph{Erreur :} \esp{Para que} + sujets différents (\emph{gobierno} vs \emph{jóvenes}) $\rightarrow$ Subjonctif obligatoire.
    \item \esp{Cuando ellos tendrán} (Futur Indicatif) $\rightarrow$ \textbf{\esp{Cuando tengan}} (Subjonctif Présent). \emph{Erreur :} Futur dans principale (\esp{podrán}) $\rightarrow$ Subjonctif après \esp{cuando} (temporel futur). \esp{Tener} irrégulier $\rightarrow$ \esp{tengan}. \esp{Ellos} redondant (sujet inclus).
    \item \esp{podran} $\rightarrow$ \textbf{\esp{podrán}}. \emph{Erreur :} Accent tonique manquant sur \emph{á} (futur 3e pl \esp{-án}).
    \item \esp{sus empresas} $\rightarrow$ \textbf{\esp{su empresa}} (ou \esp{sus empresas} si chacun plusieurs). \emph{Préférence :} \esp{su entreprise} (distributif singulier : une entreprise par jeune).
    \item \esp{Es en la agricultura donde están} $\rightarrow$ \textbf{\esp{Es en la agricultura donde se encuentran}}. \emph{Erreur :} \esp{Están} (verbe personnel « ils sont ») $\rightarrow$ Structure impersonnelle/passive réfléchie \esp{se encuentran} pour « s'y trouvent / il y a ». Accord \esp{se encuentran} (pluriel \esp{oportunidades}).
    \item \esp{Si tú eres} $\rightarrow$ \textbf{\esp{Si eres}}. \emph{Erreur :} Pronom \esp{tú} redondant (sujet inclus dans la terminaison). Garder \esp{tú} seulement si emphase/contraste (\emph{Si TÚ eres joven...}).
    \item \esp{no te quejas} (Indicatif) $\rightarrow$ \textbf{\esp{no te quejes}} (Subjonctif Présent / Impératif négatif). \emph{Erreur :} Impératif négatif = forme du subjonctif.
    \item \esp{formate} $\rightarrow$ \textbf{\esp{fórmate}}. \emph{Erreur :} Accent manquant sur \emph{á} (impératif affirmatif 2e sg \esp{fórma} + \esp{te} $\rightarrow$ \esp{fórmate}). Règle : accent sur la voyelle tonique du verbe original si pronom enclitique ajoute une syllabe.
    \item \esp{!} $\rightarrow$ \textbf{\esp{¡ ... !}}. \emph{Erreur :} Ponctuation espagnole : points d'exclamation/ouvrant \esp{¡} et fermant \esp{!} obligatoires.
\end{enumerate}
\textbf{Version corrigée :}
\esp{« \textbf{Como} el desempleo aumenta, el gobierno ha lanzado un programa \textbf{para que los jóvenes encuentren} trabajo. \textbf{Cuando tengan} formación, \textbf{podrán} crear \textbf{su empresa}. \textbf{Es en la agricultura donde se encuentran} las oportunidades. \textbf{Si eres} joven, \textbf{no te quejes: ¡fórmate!} »}
\end{exoresolu}

\begin{attention}[Les 5 erreurs qui coûtent des points au Bac]
\begin{enumerate}
    \item \textbf{Subjonctif oublié après \esp{Para que / Cuando (futur) / Antes de que / Sin que}.} \cle{Règle :} Sujets différents + Finalité = Subj. / Futur Principal + \esp{Cuando/Hasta que/En cuanto} = Subj. Présent. \emph{Ex:} \esp{*Para que ellos vienen} $\rightarrow$ \esp{Para que vengan}.
    \item \textbf{Confusion \esp{Si} conditionnel (Indicatif) vs \esp{Cuando} temporel (Subjonctif).} \cle{Piège :} \esp{Si llueve, iré} (Indicatif) MAIS \esp{Cuando llueva, iré} (Subjonctif). Ne jamais écrire \esp{*Si llueva} (sauf hypothétique très littéraire \esp{Si lloviera} = imparfait). Au Bac : \esp{Si} + Présent Indicatif.
    \item \textbf{Accents diacritiques et interrogatifs/exclamatifs :} \esp{*Si, tu, el, mas, se, mi, te} (sans accent) changent le sens ou sont fautifs. \esp{¿Cuando vienes?} (faute) vs \esp{¿Cuándo vienes?}. \esp{¡Que bien!} (faute) vs \esp{¡Qué bien!}. \esp{¿Por que?} (faute) vs \esp{¿Por qué?}.
    \item \textbf{\esp{Ser} vs \esp{Estar} sur adjectifs de nationalité/origine vs état/lieu.} \esp{Es camerunés} (identité/nationalité) vs \esp{Está en Camerún} (lieu). \esp{Está cansado} (état passager) vs \esp{Es cansado} (caractère permanent). \emph{Emphase :} \esp{Es mi padre quien está aquí} (pas \esp{*está mi padre}).
    \item \textbf{Calques français sur cause/finalité :} \esp{*Por qué} (interrogatif) utilisé comme relatif (\esp{*La razón por qué...} $\rightarrow$ \esp{La razón \textbf{por la que}...}). \esp{*Para qué} au lieu de \esp{Para que} + Subj. \esp{*A causa que} (inexistant) $\rightarrow$ \esp{A causa de} + GN / \esp{Porque} + Verbe. \esp{Debido a que} (anglicisme/calque) $\rightarrow$ \esp{Dado que} / \esp{Ya que} / \esp{Puesto que}.
\end{enumerate}
\end{attention}

\newpage

\coursec{Exercices}

\courssub{Je vérifie mes connaissances}

\begin{exercice}[QCM : Connecteurs de cause et de finalité]
Choisissez la réponse correcte pour compléter chaque phrase.
\begin{enumerate}
\item \esp{No pudo venir \_\_\_\_\_\_ estaba enfermo.} \quad A) \esp{para que} \quad B) \esp{porque} \quad C) \esp{a fin de que} \quad D) \esp{cuando}
\item \esp{Estudia mucho \_\_\_\_\_\_ sus padres estén orgullosos.} \quad A) \esp{porque} \quad B) \esp{como} \quad C) \esp{para que} \quad D) \esp{dado que}
\item \esp{\_\_\_\_\_\_ la lluvia, el partido fue cancelado.} \quad A) \esp{Debido a} \quad B) \esp{Para que} \quad C) \esp{Cuando} \quad D) \esp{Ya que}
\item \esp{Te llamaré \_\_\_\_\_\_ llegue a casa.} \quad A) \esp{porque} \quad B) \esp{cuando} \quad C) \esp{para} \quad D) \esp{ya que}
\item \esp{Lo hice \_\_\_\_\_\_ tú me lo pediste.} \quad A) \esp{con el fin de} \quad B) \esp{puesto que} \quad C) \esp{a causa de} \quad D) \esp{para}
\end{enumerate}
\end{exercice}

\begin{exercice}[Vrai ou Faux justifié]
Indiquez si les affirmations suivantes sont vraies (V) ou fausses (F) et justifiez votre réponse en citant la règle grammaticale correspondante.
\begin{enumerate}
\item Après \esp{para que}, le verbe se met toujours à l'indicatif.
\item \esp{Como} en tête de phrase introduit une cause connue de l'énonciateur.
\item Dans une subordonnée temporelle introduite par \esp{cuando} faisant référence à un événement futur, on utilise le subjonctif présent.
\item La structure emphatique \esp{Lo que + verbe + es + nom} permet de mettre l'accent sur le sujet.
\item \esp{Debido a} et \esp{a causa de} s'emploient de la même manière (même nature grammaticale).
\end{enumerate}
\end{exercice}

\begin{exercice}[Vocabulaire : Le monde du travail]
Complétez les phrases suivantes avec le mot approprié parmi la liste : \esp{contrato, entrevista, currículum, salario, prestaciones, jornada, despido, huelga, sindicato, prácticas}.
\begin{enumerate}
\item Antes de firmar el \_\_\_\_\_\_, lee bien todas las cláusulas.
\item Mañana tengo una \_\_\_\_\_\_ de trabajo en una empresa de Douala.
\item Los trabajadores exigen un aumento de \_\_\_\_\_\_ y mejores \_\_\_\_\_\_ sociales.
\item La \_\_\_\_\_\_ laboral en España es de 40 horas semanales como máximo.
\item Hice \_\_\_\_\_\_ profesionales en un hotel durante el verano.
\item El \_\_\_\_\_\_ convocó una \_\_\_\_\_\_ para protestar contra los despidos.
\end{enumerate}
\end{exercice}

\begin{exercice}[Conjugaison : Indicatif ou Subjonctif ?]
Conjuguez le verbe entre parenthèses au mode et au temps appropriés (Présent de l'indicatif ou Présent du subjonctif).
\begin{enumerate}
\item \esp{Cuando yo \_\_\_\_\_\_ (llegar) a Yaoundé, te avisaré inmediatamente.}
\item \esp{Como \_\_\_\_\_\_ (ser) temprano, no hay nadie en la oficina.}
\item \esp{Estudio español para que \_\_\_\_\_\_ (poder) trabajar en Guinea Ecuatorial.}
\item \esp{Cuando era niño, \_\_\_\_\_\_ (vivir) en Douala con mis abuelos.}
\item \esp{Te llamaré en cuanto \_\_\_\_\_\_ (saber) el resultado de la entrevista.}
\item \esp{No saldré hasta que \_\_\_\_\_\_ (terminar) el informe.}
\item \esp{Lo hago porque \_\_\_\_\_\_ (querer) ayudarte, no por obligación.}
\end{enumerate}
\end{exercice}

\courssub{Je m'entraîne}

\begin{exercice}[Traduction thématique : Cause, finalité et futur]
Traduisez les phrases suivantes en espagnol. Attention aux connecteurs, aux modes verbaux et à l'ordre des mots.
\begin{enumerate}
\item Comme il pleut, nous restons à la maison. (\esp{como} + indicatif)
\item Il travaille dur pour que sa famille vive mieux. (\esp{para que} + subjonctif)
\item Dès que j'aurai mon diplôme, je chercherai du travail. (\esp{en cuanto} / \esp{tan pronto como} + subjonctif)
\item C'est à cause de la crise économique qu'il a perdu son emploi. (\esp{debido a} / \esp{a causa de} + nom)
\item Nous étudions l'espagnol afin de pouvoir voyager en Amérique latine. (\esp{a fin de} / \esp{con el fin de} + infinitif)
\end{enumerate}
\end{exercice}

\begin{exercice}[Transformation de structures]
Réécrivez les phrases en appliquant la consigne indiquée entre parenthèses.
\begin{enumerate}
\item \esp{Estudio para aprobar el examen.} (Changer de sujet à la subordonnée de finalité : \emph{mi familia}).
\item \esp{No vine porque estaba enfermo.} (Remplacer \esp{porque} + indicatif par \esp{debido a} + nom).
\item \esp{Cuando llegues, llama a la puerta.} (Mettre la principale au futur et la subordonnée au subjonctif — déjà fait, vérifier l'accord ; ici transformer en : \esp{Cuando \_\_\_\_\_\_ (llegar), \_\_\_\_\_\_ (llamar)}).
\item \esp{Lo importante es la salud.} (Utiliser la structure \esp{Lo que} + verbe + \esp{es}...).
\item \esp{Como no tengo dinero, no puedo comprarlo.} (Inverser l'ordre : cause en second, utiliser \esp{ya que}).
\item \esp{El jefe quiere que los empleados trabajen más.} (Transformer en finalité avec sujet différent : \esp{El jefe les paga bien \_\_\_\_\_\_}).
\end{enumerate}
\end{exercice}

\begin{exercice}[Texte à trous : L'emploi des jeunes]
Complétez le texte suivant avec les mots ou formes verbales appropriés (connecteurs, verbes au mode convenable, prépositions).
\begin{textebox}[Extrait adapté — Empleo juvenil]
En Camerún, muchos jóvenes buscan trabajo \_\_\_\_\_\_ (1. finalidad) mantener a sus familias. \_\_\_\_\_\_ (2. Cause : Comme) la economía es informal, \_\_\_\_\_\_ (3. hay) pocos contratos fijos. El gobierno crea programas \_\_\_\_\_\_ (4. finalidad : a fin de que) los jóvenes \_\_\_\_\_\_ (5. poder) acceder a microcréditos. \_\_\_\_\_\_ (6. Tiempo : Cuando) estos proyectos \_\_\_\_\_\_ (7. funcionar) bien, la pobreza \_\_\_\_\_\_ (8. disminuir). \_\_\_\_\_\_ (9. Emphase : C'est l'effort) \_\_\_\_\_\_ (10. qui) \_\_\_\_\_\_ (11. apporte) des résultats, pas la chance.
\end{textebox}
\end{exercice}

\begin{exercice}[Appariement : Fonction et structure]
Reliez chaque énoncé de la colonne A à sa fonction/structure grammaticale dans la colonne B.
\begin{center}
\begin{tabularx}{\linewidth}{|X|X|}
\hline
\textbf{Colonne A (Énoncé / Extrait)} & \textbf{Colonne B (Fonction / Structure)} \\
\hline
\esp{Debido a la huelga, no hay autobuses.} & 1. Finalité (sujet différent) — Subjonctif \\
\hline
\esp{Estudio para que apruebes tú.} & 2. Cause (nom / SN) — Préposition \\
\hline
\esp{Te aviso cuando llegue.} & 3. Cause (proposition) — Conjonction (antéposée, connue) \\
\hline
\esp{Como no estudias, suspendes.} & 4. Futur dans subordonnée temporelle — Subjonctif \\
\hline
\esp{Lo que quiero es justicia.} & 5. Emphase (focus sur l'objet/attribut) — \esp{Lo que} \\
\hline
\esp{A fin de que no haya errores, revisa.} & 6. Finalité (forme soutenue) — Subjonctif \\
\hline
\end{tabularx}
\end{center}
\end{exercice}

\courssub{Je me prépare au Bac}

\begin{exercice}[Compréhension et analyse linguistique — Type Bac]
Lisez le texte suivant et répondez aux questions en français (sauf indication contraire).
\begin{textebox}[La realidad del joven trabajador en Yaoundé]
En Yaoundé, la vida es cara y los jóvenes lo saben bien. \textbf{Porque} los salarios son bajos, muchos chicos trabajan en el sector informal \textbf{para que} sus hermanos menores puedan ir a la escuela. \textbf{Cuando} termine \textbf{la jornada}, vuelven a casa agotados, pero orgullosos. \textbf{Lo que} les duele \textbf{es} la inseguridad del mañana: \textbf{cuando} \_\_\_\_\_\_ (tener) una emergencia médica, no hay ahorros. El gobierno promete crear empleo, \textbf{ya que} la estabilidad social \_\_\_\_\_\_ (depender) de ello. \textbf{A fin de que} estas promesas \_\_\_\_\_\_ (cumplir), los jóvenes exigen transparencia. \textbf{Debido a} su determinación, el cambio llegará.
\end{textebox}

\textbf{Questions de compréhension (en français) :}
\begin{enumerate}
\item Quel est le problème principal évoqué dans la première phrase ?
\item Pourquoi les jeunes acceptent-ils des jobs informels ? (Cause explicite dans le texte).
\item Quel sentiment exprime \esp{orgullosos} malgré la fatigue ?
\item Quelle est la revendication finale des jeunes ?
\end{enumerate}

\textbf{Questions de langue :}
\begin{enumerate}
\setcounter{enumi}{4}
\item Conjuguez les verbes entre parenthèses au mode et au temps convenables (lignes 5, 6, 7).
\item Justifiez l'emploi du subjonctif (ou de l'indicatif) pour \esp{puedan} (ligne 2) et pour \esp{tenga} (votre réponse à la q.5).
\item Identifiez et analysez la structure emphatique présente dans la phrase \esp{« Lo que les duele es la inseguridad... »}. Traduisez cette phrase en français.
\item Repérez deux connecteurs de cause différents dans le texte et précisez la nature grammaticale de ce qui les suit (proposition / nom).
\item Transformez la phrase \esp{« Debido a su determinación, el cambio llegará »} en utilisant \esp{como} + indicatif en début de phrase.
\end{enumerate}
\end{exercice}

\begin{exercice}[Thème et Version — Type Bac]

\textbf{Thème (Français $\rightarrow$ Espagnol) :} Traduisez les 5 phrases suivantes en espagnol. Vous mobiliserez : cause (\esp{porque, como, ya que, debido a}), finalité (\esp{para, para que, a fin de que}), futur dans la subordonnée (\esp{cuando, en cuanto} + subjonctif), emphase (\esp{lo que, lo de}).
\begin{enumerate}
\item Comme la vie est chère à Douala, il travaille le soir pour que ses enfants puissent étudier.
\item Je t'appellerai dès que j'aurai des nouvelles de l'entretien.
\item C'est l'effort qui réussit, pas la chance. (Emphase sur le sujet \emph{l'effort}).
\item À cause de la pluie, le match a été annulé. (Préposition + nom).
\item Nous manifestons afin que le gouvernement entende nos revendications. (Finalité soutenue, sujet différent).
\end{enumerate}

\textbf{Version (Espagnol $\rightarrow$ Français) :} Traduisez le paragraphe suivant en français soigné.
\begin{textebox}[Texte à traduire]
Muchos licenciados cameruneses no encuentran trabajo \textbf{porque} la economía no genera suficientes plazas formales. \textbf{Por eso}, emigran \textbf{para que} sus familias tengan un futuro mejor. \textbf{Cuando} \textbf{lleguen} a Europa, \textbf{se darán cuenta} de que \textbf{lo que} \textbf{brilla} no siempre \textbf{es} oro. \textbf{Debido a} la nostalgia y las barreras legales, \textbf{el camino} \textbf{es} duro. \textbf{Sin embargo}, \textbf{como} tienen fe, \textbf{siguen} adelante.
\end{textebox}
\end{exercice}

\begin{exercice}[Expression écrite — Sujet de rédaction Type Bac]
\textbf{Sujet :} « \emph{El desempleo juvenil es el mayor desafío para el desarrollo de África.} »
Rédigez un article d'opinion (entre 150 et 200 mots) pour le journal de votre lycée ou un blog jeunesse, dans lequel vous :
\begin{itemize}
\item Présentez la situation du chômage des jeunes dans votre pays (Cameroun) ou en Afrique.
\item Analysez les causes (éducation inadaptée, croissance démographique, secteur privé faible, corruption...).
\item Proposez des solutions concrètes (entrepreneuriat, formation professionnelle, politiques publiques, agriculture...).
\item Utilisez au moins \textbf{trois} connecteurs de cause (\esp{porque, como, ya que, puesto que, debido a, a causa de}), \textbf{deux} structures de finalité (\esp{para + inf., para que + subj., a fin de que, con el fin de}) et \textbf{une} structure emphatique (\esp{lo que, lo de, ser... quien/lo que}).
\end{itemize}
\textit{Conseil : Structurez votre article (Titre, Introduction, Développement en 2-3 paragraphes, Conclusion). Soignez la correction linguistique (accords, accents, subjonctifs).}
\end{exercice}

\newpage

\coursec{Corrigés des exercices}

\begin{corrige}[Exercice 1 — QCM : Connecteurs de cause et de finalité]
\begin{enumerate}
\item \textbf{B) \esp{porque}} : \esp{No pudo venir porque estaba enfermo.} « Il n'a pas pu venir parce qu'il était malade. » Cause réelle et affirmée $\Rightarrow$ indicatif ; \esp{para que} et \esp{a fin de que} expriment la finalité, \esp{cuando} le temps.
\item \textbf{C) \esp{para que}} : \esp{Estudia mucho para que sus padres estén orgullosos.} « Il étudie beaucoup pour que ses parents soient fiers. » Finalité avec changement de sujet $\Rightarrow$ \esp{para que} + subjonctif (\esp{estén}).
\item \textbf{A) \esp{Debido a}} : \esp{Debido a la lluvia, el partido fue cancelado.} « À cause de la pluie, le match a été annulé. » \esp{Debido a} + nom ; les autres options exigeraient une proposition avec verbe conjugué.
\item \textbf{B) \esp{cuando}} : \esp{Te llamaré cuando llegue a casa.} « Je t'appellerai quand j'arriverai à la maison. » Subordonnée temporelle renvoyant au futur $\Rightarrow$ subjonctif présent (\esp{llegue}).
\item \textbf{B) \esp{puesto que}} : \esp{Lo hice puesto que tú me lo pediste.} « Je l'ai fait puisque tu me l'as demandé. » Cause présentée comme connue, registre soutenu, + indicatif.
\end{enumerate}
\textbf{Points de vigilance :} \esp{para} + infinitif (même sujet) ; \esp{para que} + subjonctif (sujets différents) ; \esp{debido a} / \esp{a causa de} + nom (ou infinitif), jamais + verbe conjugué.
\end{corrige}

\begin{corrige}[Exercice 2 — Vrai ou Faux justifié]
\begin{enumerate}
\item \textbf{FAUX.} Après \esp{para que}, le verbe se met toujours au \textbf{subjonctif} : \esp{Trabajo para que vivas mejor.} « Je travaille pour que tu vives mieux. » La finalité avec sujet différent exige le subjonctif (action visée, non encore réalisée).
\item \textbf{VRAI.} \esp{Como} en tête de phrase présente une cause connue ou évidente pour l'interlocuteur : \esp{Como llueve, nos quedamos en casa.} « Comme il pleut, nous restons à la maison. »
\item \textbf{VRAI.} Après \esp{cuando} (ainsi que \esp{en cuanto, tan pronto como, hasta que}) renvoyant à un événement futur non encore réalisé, on emploie le subjonctif présent : \esp{Cuando llegue, te llamaré.} « Quand j'arriverai, je t'appellerai. »
\item \textbf{FAUX.} \esp{Lo que + verbe + es + nom} met l'accent sur le \textbf{complément ou l'attribut}, pas sur le sujet : \esp{Lo que quiero es justicia.} « Ce que je veux, c'est la justice. » Pour insister sur le sujet : \esp{Es Juan quien lo hizo.} « C'est Juan qui l'a fait. »
\item \textbf{VRAI.} \esp{Debido a} et \esp{a causa de} sont deux locutions prépositives suivies d'un nom (ou d'un infinitif) : \esp{debido a la crisis} / \esp{a causa de la crisis} « à cause de la crise ».
\end{enumerate}
\end{corrige}

\begin{corrige}[Exercice 3 — Vocabulaire : Le monde du travail]
\begin{enumerate}
\item \esp{Antes de firmar el \textbf{contrato}, lee bien todas las cláusulas.} « Avant de signer le contrat, lis bien toutes les clauses. »
\item \esp{Mañana tengo una \textbf{entrevista} de trabajo en una empresa de Duala.} « Demain j'ai un entretien d'embauche dans une entreprise de Douala. »
\item \esp{Los trabajadores exigen un aumento de \textbf{salario} y mejores \textbf{prestaciones} sociales.} « Les travailleurs exigent une augmentation de salaire et de meilleures prestations sociales. »
\item \esp{La \textbf{jornada} laboral en España es de 40 horas semanales como máximo.} « La durée hebdomadaire du travail en Espagne est de 40 heures au maximum. »
\item \esp{Hice unas \textbf{prácticas} profesionales en un hotel durante el verano.} « J'ai fait un stage professionnel dans un hôtel pendant l'été. »
\item \esp{El \textbf{sindicato} convocó una \textbf{huelga} para protestar contra los despidos.} « Le syndicat a appelé à une grève pour protester contre les licenciements. »
\end{enumerate}
\textbf{Attention :} \esp{prácticas} = stage (faux ami avec « pratiques ») ; \esp{despido} = licenciement (ne pas confondre avec \esp{despedida}, les adieux) ; \esp{currículum} (CV) n'apparaît pas ici mais est à connaître.
\end{corrige}

\begin{corrige}[Exercice 4 — Conjugaison : Indicatif ou Subjonctif ?]
\begin{enumerate}
\item \esp{Cuando yo \textbf{llegue} a Yaundé, te avisaré inmediatamente.} Subjonctif : futur dans la subordonnée temporelle. « Quand j'arriverai à Yaoundé, je te préviendrai immédiatement. »
\item \esp{Como \textbf{es} temprano, no hay nadie en la oficina.} Indicatif : cause réelle, constatée. « Comme il est tôt, il n'y a personne au bureau. »
\item \esp{Estudio español para que mi familia \textbf{pueda} vivir mejor.} Subjonctif après \esp{para que} (finalité, sujet différent). « J'étudie l'espagnol pour que ma famille puisse vivre mieux. »
\item \esp{Cuando era niño, \textbf{vivía} en Duala con mis abuelos.} Imparfait de l'indicatif : habitude passée, pas de futur. « Quand j'étais enfant, je vivais à Douala avec mes grands-parents. »
\item \esp{Te llamaré en cuanto \textbf{sepa} el resultado de la entrevista.} Subjonctif : futur après \esp{en cuanto}. « Je t'appellerai dès que je saurai le résultat de l'entretien. »
\item \esp{No saldré hasta que \textbf{termine} el informe.} Subjonctif : futur après \esp{hasta que}. « Je ne sortirai pas avant d'avoir terminé le rapport. »
\item \esp{Lo hago porque \textbf{quiero} ayudarte, no por obligación.} Indicatif : cause réelle et affirmée. « Je le fais parce que je veux t'aider, pas par obligation. »
\end{enumerate}
\textbf{Règle-clé :} cause réelle $\Rightarrow$ indicatif ; finalité ou futur non encore réalisé $\Rightarrow$ subjonctif.
\end{corrige}

\begin{corrige}[Exercice 5 — Traduction thématique : Cause, finalité et futur]
\begin{enumerate}
\item \esp{Como llueve, nos quedamos en casa.} (\esp{como} en tête + indicatif présent.)
\item \esp{Trabaja duro para que su familia viva mejor.} (Subjonctif \esp{viva} : changement de sujet.)
\item \esp{En cuanto tenga mi título, buscaré trabajo.} ou \esp{Tan pronto como obtenga mi diploma, buscaré empleo.} (Subjonctif présent \esp{tenga} pour le futur antérieur français « dès que j'aurai ».)
\item \esp{Ha perdido su empleo debido a la crisis económica.} ou, avec emphase : \esp{Es debido a la crisis económica que ha perdido su empleo.} (\esp{debido a} + nom.)
\item \esp{Estudiamos español a fin de poder viajar por América Latina.} ou \esp{...con el fin de poder viajar...} (Même sujet $\Rightarrow$ \esp{a fin de} + infinitif.)
\end{enumerate}
\textbf{Points de vigilance :} le futur antérieur français (« dès que j'aurai obtenu ») se rend par le subjonctif présent espagnol ; ne jamais dire \esp{*para que su familia vive} (indicatif impossible après \esp{para que}).
\end{corrige}

\begin{corrige}[Exercice 6 — Transformation de structures]
\begin{enumerate}
\item \esp{Estudio para aprobar el examen.} $\rightarrow$ \esp{Estudio para que mi familia esté orgullosa de mí.} « J'étudie pour que ma famille soit fière de moi. » (Changement de sujet $\Rightarrow$ \esp{para que} + subjonctif \esp{esté}.)
\item \esp{No vine porque estaba enfermo.} $\rightarrow$ \esp{No vine debido a una enfermedad.} / \esp{No vine a causa de mi enfermedad.} (Proposition causale $\rightarrow$ locution prépositive + nom.)
\item \esp{Cuando \textbf{llegues}, \textbf{llamarás} a la puerta.} « Quand tu arriveras, tu frapperas à la porte. » (Subjonctif dans la subordonnée, futur dans la principale.)
\item \esp{La salud es importante.} $\rightarrow$ \esp{Lo que es importante es la salud.} « Ce qui est important, c'est la santé. » (Emphase avec \esp{Lo que... es...}.)
\item \esp{Como no tengo dinero, no puedo comprarlo.} $\rightarrow$ \esp{No puedo comprarlo, ya que no tengo dinero.} (Cause en seconde position avec \esp{ya que} + indicatif.)
\item \esp{El jefe les paga bien. Así trabajan más.} $\rightarrow$ \esp{El jefe les paga bien para que trabajen más.} « Le patron les paie bien pour qu'ils travaillent davantage. » (Finalité avec sujet différent $\Rightarrow$ subjonctif \esp{trabajen}.)
\end{enumerate}
\end{corrige}

\begin{corrige}[Exercice 7 — Texte à trous : L'emploi des jeunes]
Texte complété :
\begin{enumerate}
\item \esp{para} (finalité, même sujet $\Rightarrow$ infinitif \esp{mantener})
\item \esp{Como} (cause connue, en tête de phrase)
\item \esp{hay} (indicatif : constat réel)
\item \esp{a fin de que} (finalité soutenue, sujet différent)
\item \esp{puedan} (subjonctif après \esp{a fin de que})
\item \esp{Cuando}
\item \esp{funcionen} (subjonctif : futur dans la subordonnée temporelle)
\item \esp{disminuirá} (futur dans la principale)
\item \esp{Es el esfuerzo} (début de la structure emphatique)
\item \esp{lo que}
\item \esp{aporta} (indicatif : \esp{Es el esfuerzo lo que aporta resultados, no la suerte.})
\end{enumerate}
Traduction du texte complété : « Au Cameroun, beaucoup de jeunes cherchent du travail pour subvenir aux besoins de leurs familles. Comme l'économie est informelle, il y a peu de contrats fixes. Le gouvernement crée des programmes afin que les jeunes puissent accéder aux microcrédits. Quand ces projets fonctionneront bien, la pauvreté diminuera. C'est l'effort qui apporte des résultats, pas la chance. »
\end{corrige}

\begin{corrige}[Exercice 8 — Appariement : Fonction et structure]
\begin{center}
\begin{tabularx}{\linewidth}{|X|X|}
\hline
\rowcolor{popblueL} \textbf{Énoncé} & \textbf{Fonction / Structure} \\
\hline
\esp{Debido a la huelga, no hay autobuses.} & 2. Cause (nom / SN) — locution prépositive \\
\hline
\esp{Estudio para que apruebes tú.} & 1. Finalité (sujet différent) — subjonctif \\
\hline
\esp{Te aviso cuando llegue.} & 4. Futur dans la subordonnée temporelle — subjonctif \\
\hline
\esp{Como no estudias, suspendes.} & 3. Cause (proposition) — conjonction antéposée, cause connue \\
\hline
\esp{Lo que quiero es justicia.} & 5. Emphase (focus sur l'objet) — \esp{Lo que... es...} \\
\hline
\esp{A fin de que no haya errores, revisa.} & 6. Finalité (registre soutenu) — subjonctif \\
\hline
\end{tabularx}
\end{center}
\textbf{Justifications :} \esp{debido a} est suivi du nom \esp{la huelga} ; \esp{para que} introduit un sujet différent (\esp{tú}), d'où le subjonctif \esp{apruebes} ; \esp{cuando llegue} renvoie à un événement futur ; \esp{como} antéposé signale une cause connue de l'interlocuteur ; \esp{lo que... es} focalise \esp{justicia} ; \esp{a fin de que} est la forme soutenue de \esp{para que} et exige le subjonctif (\esp{haya}).
\end{corrige}

\begin{corrige}[Exercice 9 — Compréhension et analyse linguistique — Type Bac]
\textbf{Barème indicatif : compréhension 4 pts (1 pt/question) ; langue 6 pts.}

\textbf{Compréhension :}
\begin{enumerate}
\item Le problème principal est le coût élevé de la vie à Yaoundé combiné à la faiblesse des salaires, ce qui pousse les jeunes vers le secteur informel.
\item Ils acceptent ces emplois \textbf{parce que les salaires (du secteur formel) sont bas} (\esp{porque los salarios son bajos}) et pour permettre à leurs frères et sœurs cadets d'aller à l'école.
\item \esp{Orgullosos} exprime la fierté : malgré l'épuisement, les jeunes sont fiers de subvenir aux besoins de leur famille.
\item Les jeunes exigent la \textbf{transparence} pour que les promesses de création d'emplois soient tenues.
\end{enumerate}

\textbf{Langue :}
\begin{enumerate}
\setcounter{enumi}{4}
\item \esp{cuando \textbf{tengan} una emergencia médica} (subjonctif : futur) ; \esp{la estabilidad social \textbf{depende} de ello} (indicatif : cause réelle après \esp{ya que}) ; \esp{a fin de que estas promesas \textbf{se cumplan}} (subjonctif : finalité).
\item \esp{Puedan} : subjonctif après \esp{para que} (finalité, sujet différent : \esp{sus hermanos}). \esp{Tengan} : subjonctif après \esp{cuando} renvoyant à un événement futur incertain.
\item Structure emphatique : \esp{Lo que + verbe + es + nom}. Elle met en relief le complément \esp{la inseguridad del mañana}. Traduction : « Ce qui leur fait mal, c'est l'insécurité du lendemain. »
\item Deux connecteurs de cause : \esp{porque} + proposition verbale (\esp{los salarios son bajos}) ; \esp{debido a} + nom (\esp{su determinación}). On peut ajouter \esp{ya que} + proposition.
\item \esp{Como están decididos, el cambio llegará.} ou \esp{Como su determinación es grande, el cambio llegará.} « Comme ils sont déterminés, le changement viendra. »
\end{enumerate}
\textbf{Points de vigilance :} bien distinguer \esp{porque} / \esp{ya que} + verbe conjugué et \esp{debido a} + nom ; ne jamais mettre l'indicatif après \esp{a fin de que}.
\end{corrige}

\begin{corrige}[Exercice 10 — Thème et Version — Type Bac]
\textbf{Barème indicatif : thème 5 pts (1 pt/phrase) ; version 5 pts.}

\textbf{Thème — propositions de traduction :}
\begin{enumerate}
\item \esp{Como la vida es cara en Duala, trabaja por la noche para que sus hijos puedan estudiar.}
\item \esp{Te llamaré en cuanto tenga noticias de la entrevista.}
\item \esp{Es el esfuerzo lo que triunfa, no la suerte.}
\item \esp{Debido a la lluvia, el partido fue cancelado.} / \esp{A causa de la lluvia, se anuló el partido.}
\item \esp{Nos manifestamos a fin de que el gobierno escuche nuestras reivindicaciones.}
\end{enumerate}
\textbf{Points de vigilance :} subjonctifs \esp{puedan}, \esp{tenga}, \esp{escuche} ; \esp{debido a} + nom, sans verbe conjugué ; \esp{Es... lo que...} pour l'emphase sur le sujet ; \esp{manifestarse} est pronominal (\esp{nos manifestamos}).

\textbf{Version — traduction modèle :}
« Beaucoup de diplômés camerounais ne trouvent pas de travail parce que l'économie ne génère pas suffisamment d'emplois formels. C'est pourquoi ils émigrent pour que leurs familles aient un avenir meilleur. Quand ils arriveront en Europe, ils se rendront compte que tout ce qui brille n'est pas or. À cause de la nostalgie et des barrières juridiques, le chemin est dur. Cependant, comme ils ont la foi, ils continuent d'avancer. »
\textbf{Points de vigilance :} \esp{licenciados} = diplômés (faux ami avec « licenciés ») ; \esp{plazas} = postes, emplois ; \esp{se darán cuenta de que} = ils se rendront compte que ; rendre \esp{no todo lo que brilla es oro} par le proverbe français correspondant.
\end{corrige}

\begin{corrige}[Exercice 11 — Expression écrite — Sujet de rédaction Type Bac]
\textbf{Barème indicatif (sur 20) :} respect du sujet et de la structure 4 pts ; richesse du lexique du travail 4 pts ; correction grammaticale (accords, accents) 4 pts ; emploi des connecteurs exigés (3 causes, 2 finalités, 1 emphase) 6 pts ; cohérence et articulation 2 pts.

\textbf{Proposition de rédaction modèle :}

\esp{\textbf{El desempleo juvenil: la batalla del mañana}}

\esp{En Camerún, miles de jóvenes terminan sus estudios cada año, pero no encuentran empleo. \textbf{Como} la economía formal es débil, muchos licenciados venden en los mercados o conducen mototaxis \textbf{para} sobrevivir. Esta situación es dramática \textbf{porque} desperdicia el talento de toda una generación.}

\esp{Las causas son múltiples: la educación es a veces inadaptada a las necesidades de las empresas, \textbf{ya que} privilegia la teoría sobre la práctica. Además, \textbf{debido a} la corrupción y a la debilidad del sector privado, los puestos de trabajo son escasos. \textbf{Lo que} más duele \textbf{es} ver a jóvenes diplomados sin esperanza.}

\esp{Sin embargo, hay soluciones. El Estado debe crear programas \textbf{a fin de que} los jóvenes accedan a la formación profesional y a los microcréditos. También hay que apoyar el emprendimiento \textbf{para que} los jóvenes creen sus propias empresas. \textbf{Puesto que} la juventud es el futuro de África, invertir en ella es invertir en el desarrollo. \textbf{Cuando} los gobiernos \textbf{entiendan} esto, el continente avanzará.}

Traduction (sens général) : « Au Cameroun, des milliers de jeunes finissent leurs études chaque année sans trouver d'emploi. Comme l'économie formelle est faible, beaucoup de diplômés vendent dans les marchés ou conduisent des motos-taxis pour survivre. Cette situation est dramatique parce qu'elle gaspille le talent de toute une génération. Les causes sont multiples : l'éducation est parfois inadaptée aux besoins des entreprises, puisqu'elle privilégie la théorie sur la pratique. De plus, à cause de la corruption et de la faiblesse du secteur privé, les postes de travail sont rares. Ce qui fait le plus mal, c'est de voir des jeunes diplômés sans espoir. Cependant, il y a des solutions. L'État doit créer des programmes afin que les jeunes accèdent à la formation professionnelle et aux microcrédits. Il faut aussi soutenir l'entrepreneuriat pour que les jeunes créent leurs propres entreprises. Puisque la jeunesse est l'avenir de l'Afrique, investir en elle, c'est investir dans le développement. Quand les gouvernements comprendront cela, le continent avancera. »

\textbf{Points de vigilance :}
\begin{itemize}
\item Vérifier chaque subjonctif après \esp{para que}, \esp{a fin de que}, \esp{cuando} + futur : \esp{accedan, creen, entiendan}.
\item \esp{Debido a} / \esp{a causa de} + nom, jamais + verbe conjugué.
\item Accents : \esp{economía, dramática, teoría, práctica, además, corrupción} ; pas d'accent sur \esp{esperanza}.
\item Faux amis : \esp{licenciado} = diplômé (et non « licencié ») ; « un emploi » = \esp{un empleo / un puesto de trabajo} ; \esp{embarazada} $\neq$ « embarrassée ».
\end{itemize}
\end{corrige}

\newpage

\coursec{Fiche bilan}

\coursec{Leçon E21 : Traduction : phrase causale, finalité, futur dans la subordonnée, emphase (révision)}

\begin{popfigure}
\begin{center}\tcbox[colback=poporange,colframe=poporange,coltext=white,arc=9pt,boxsep=4pt,left=6pt,right=6pt]{\bfseries\small Leçon E21 Traduction Causalité, Finalité, Futur sub., Emphase}\end{center}
\vspace{-2pt}
\begin{tcbraster}[raster columns=3,raster equal height=rows,raster column skip=6pt,raster row skip=6pt,enhanced,blankest]
\begin{tcolorbox}[enhanced,colback=white,colframe=poporange,boxrule=0.9pt,arc=3pt,title={Phrase causale},fonttitle=\bfseries\scriptsize,coltitle=white,colbacktitle=poporange,left=4pt,right=4pt,top=2pt,bottom=2pt,boxsep=1pt,toptitle=1.5pt,bottomtitle=1.5pt,halign=left]
\begin{itemize}[nosep,leftmargin=*,itemsep=1pt,font=\scriptsize]
\item \textit{porque, como}
\item \textit{puesto que}
\item \textit{ya que}
\item \textit{a causa de / debido a + nom}
\end{itemize}
\end{tcolorbox}
\begin{tcolorbox}[enhanced,colback=white,colframe=popgreen,boxrule=0.9pt,arc=3pt,title={Finalité},fonttitle=\bfseries\scriptsize,coltitle=white,colbacktitle=popgreen,left=4pt,right=4pt,top=2pt,bottom=2pt,boxsep=1pt,toptitle=1.5pt,bottomtitle=1.5pt,halign=left]
\begin{itemize}[nosep,leftmargin=*,itemsep=1pt,font=\scriptsize]
\item \textit{para + inf.}
\item \textit{para que + subj.}
\item \textit{a fin de que + subj.}
\item \textit{con el fin de + inf.}
\end{itemize}
\end{tcolorbox}
\begin{tcolorbox}[enhanced,colback=white,colframe=poppurple,boxrule=0.9pt,arc=3pt,title={Futur dans la subordonnée},fonttitle=\bfseries\scriptsize,coltitle=white,colbacktitle=poppurple,left=4pt,right=4pt,top=2pt,bottom=2pt,boxsep=1pt,toptitle=1.5pt,bottomtitle=1.5pt,halign=left]
\begin{itemize}[nosep,leftmargin=*,itemsep=1pt,font=\scriptsize]
\item \textit{cuando, en cuanto}
\item \textit{hasta que}
\item \textit{tan pronto como}
\item Subjonctif présent
\end{itemize}
\end{tcolorbox}
\begin{tcolorbox}[enhanced,colback=white,colframe=poppink,boxrule=0.9pt,arc=3pt,title={Emphase (Révision)},fonttitle=\bfseries\scriptsize,coltitle=white,colbacktitle=poppink,left=4pt,right=4pt,top=2pt,bottom=2pt,boxsep=1pt,toptitle=1.5pt,bottomtitle=1.5pt,halign=left]
\begin{itemize}[nosep,leftmargin=*,itemsep=1pt,font=\scriptsize]
\item Cleft: \textit{Es... quien/lo que}
\item Fronting (topicalisation)
\item \textit{Lo que... es...}
\item Pronoms toniques
\end{itemize}
\end{tcolorbox}
\begin{tcolorbox}[enhanced,colback=white,colframe=popgold,boxrule=0.9pt,arc=3pt,title={Connecteurs clés},fonttitle=\bfseries\scriptsize,coltitle=white,colbacktitle=popgold,left=4pt,right=4pt,top=2pt,bottom=2pt,boxsep=1pt,toptitle=1.5pt,bottomtitle=1.5pt,halign=left]
\begin{itemize}[nosep,leftmargin=*,itemsep=1pt,font=\scriptsize]
\item \textit{por / a causa de}
\item \textit{para / con objeto de}
\item \textit{si (condition)}
\item \textit{como (cause)}
\end{itemize}
\end{tcolorbox}
\begin{tcolorbox}[enhanced,colback=white,colframe=popdark,boxrule=0.9pt,arc=3pt,title={Stratégies traduction},fonttitle=\bfseries\scriptsize,coltitle=white,colbacktitle=popdark,left=4pt,right=4pt,top=2pt,bottom=2pt,boxsep=1pt,toptitle=1.5pt,bottomtitle=1.5pt,halign=left]
\begin{itemize}[nosep,leftmargin=*,itemsep=1pt,font=\scriptsize]
\item Identifier la relation
\item Choisir le connecteur
\item Vérifier le mode (ind./subj.)
\item Adapter l'emphase
\end{itemize}
\end{tcolorbox}
\end{tcbraster}

\legende{Carte mentale récapitulative : 6 pôles pour maîtriser la traduction des relations logiques et de l'emphase.}
\end{popfigure}

\begin{textebox}[Observation guidée : Témoignage d'un jeune entrepreneur à Yaoundé]
\esp{Como no encuentro trabajo en la administración pública, he decidido crear mi propia empresa de transformación de cacao. Lo hago para que mi familia viva mejor. Cuando tenga el capital necesario, compraré las máquinas. Lo que más me motiva es la independencia.}
\end{textebox}

\begin{propriete}[Lexique clé : Connecteurs et marqueurs]
\begin{center}
\begin{tabularx}{\linewidth}{|X|X|X|}
\hline
\rowcolor{popblueL}
\textbf{Français} & \textbf{Español} & \textbf{Niveau / Registre} \\
\hline
Parce que / Comme & \esp{porque}, \esp{como} (début phrase) & Courant / Neutre \\
\hline
Puisque / Vu que & \esp{puesto que}, \esp{ya que}, \esp{dado que} & Soutenu / Écrit \\
\hline
À cause de / En raison de & \esp{a causa de}, \esp{debido a} (+ nom / infinitif) & Formel / Administratif \\
\hline
Pour / Afin de & \esp{para}, \esp{a fin de}, \esp{con el fin de} (+ inf.) & Standard \\
\hline
Pour que / Afin que & \esp{para que}, \esp{a fin de que}, \esp{con el fin de que} (+ subj.) & Standard / Formel \\
\hline
Quand / Dès que & \esp{cuando}, \esp{en cuanto}, \esp{tan pronto como} & Temps / Subjonctif si futur \\
\hline
Jusqu'à ce que & \esp{hasta que} & Subjonctif si futur \\
\hline
\end{tabularx}
\end{center}
\end{propriete}

\begin{propriete}[Règles de grammaire essentielles]
\begin{center}
\begin{tabularx}{\linewidth}{|p{3.2cm}|X|X|}
\hline
\rowcolor{popblueL}
\textbf{Notion} & \textbf{Règle / Structure} & \textbf{Pièges / Exceptions} \\
\hline
\textbf{Cause : Conjonctions} & \esp{porque} (réponse à \esp{¿por qué?}), \esp{como} (cause connue, en tête), \esp{puesto que / ya que / dado que} (cause évidente, formel). & \esp{Por que} (2 mots) = \esp{por} + \esp{que} (relatif) ou \esp{porqué} (nom). \textbf{Accents diacritiques !} \\
\hline
\textbf{Cause : Locutions} & \esp{a causa de} / \esp{debido a} + \textbf{Nom / Infinitif}. Accord : \esp{debido a} (adj.) $\rightarrow$ \esp{debida a}, \esp{debidos a}, \esp{debidas a} (accord avec le nom noyau du sujet ou \textit{motivo/causa}). & \textit{Debido a que} + indicatif (cause réelle). Éviter *\esp{debido a que} + subj. \esp{A causa de} invariable. \\
\hline
\textbf{Finalité : Même sujet} & \esp{Para / A fin de / Con el fin de} + \textbf{Infinitif}. & \esp{Para} = but immédiat/concret. \esp{A fin de / Con el fin de} = but formel, planifié, souvent écrit. \\
\hline
\textbf{Finalité : Sujets différents} & \esp{Para que / A fin de que / Con el fin de que} + \textbf{Subjonctif} (présent ou imparfait). & \textbf{JAMAIS d'indicatif.} Principal Présent/Futur $\rightarrow$ Subj. Présent. Principal Passé/Conditionnel $\rightarrow$ Subj. Imparfait (\esp{-ra} ou \esp{-se}). \\
\hline
\textbf{Futur dans la sub. (Temps)} & Conjonctions : \esp{cuando, en cuanto, tan pronto como, hasta que, después (de) que, antes (de) que} + \textbf{Subjonctif présent} (valeur futur/antériorité). & \textit{Después de que} / \textit{antes de que} $\rightarrow$ \textbf{TOUJOURS} subjonctif. \textit{Cuando} + Indicatif = habitude ou passé accompli (\esp{Cuando llego, como} / \esp{Cuando llegué, comí}). \\
\hline
\textbf{Emphase : Cleft} & \esp{Es / Era / Fue / Será} + \textbf{élément focalisé} + \esp{quien / que / lo que / lo cual}. & \esp{Quien} pour personnes (antécédent humain). \esp{Que / lo que} pour choses/phrases. \esp{Lo que} = « ce qui / ce que » (antécédent nul). \\
\hline
\textbf{Emphase : Fronting} & \textbf{Objet / Complément} + \textbf{Verbe} + \textbf{Sujet}. Ex: \esp{El libro lo leyó Juan.} & Marque l'info nouvelle (thème/rhème). Fréquent à l'oral et littéraire. Pronom redondant (\esp{lo/la/le/les}) si objet défini/animé. \\
\hline
\textbf{Emphase : \textit{Lo que}} & \esp{Lo que} + \textbf{verbe} + \esp{es/era} + \textbf{infinitif / nom / adj.} Ex: \esp{Lo que quiero es dormir.} & Structure nominale figée. Équivaut à « Ce que je veux, c'est... ». \esp{Lo que} = pronom neutre « ce qui/ce que ». \\
\hline
\end{tabularx}
\end{center}
\end{propriete}

\begin{methode}[Traduire une phrase causale (Version $\rightarrow$ Thème)]
\begin{enumerate}
\item Analyser la nature du mot suivant le connecteur : \textbf{Verbe conjugué} $\rightarrow$ conjonction (\esp{porque, como, ya que...}) ; \textbf{Nom / Infinitif} $\rightarrow$ locution (\esp{a causa de, debido a}).
\item Vérifier le registre : familier (\esp{porque}), neutre (\esp{ya que}), formel (\esp{puesto que, debido a}).
\item Attention à \esp{debido a} (accord adjectif : \esp{La falta fue debida a...}) vs \esp{a causa de} (invariable).
\item Si cause + conséquence dans même phrase : ordre \esp{Como llueve, no salgo} (cause en tête = \esp{como}).
\end{enumerate}
\end{methode}

\begin{methode}[Traduire une finalité (Thème $\rightarrow$ Version)]
\begin{enumerate}
\item Identifier les \textbf{sujets} des deux propositions.
\item \textbf{Même sujet} $\rightarrow$ \esp{Para / A fin de / Con el fin de} + \textbf{Infinitif}.
\item \textbf{Sujets différents} $\rightarrow$ \esp{Para que / A fin de que / Con el fin de que} + \textbf{Subjonctif}.
\item Choisir le temps du subjonctif : Principal au Présent/Futur $\rightarrow$ Subj. Présent ; Principal au Passé/Conditionnel $\rightarrow$ Subj. Imparfait (\esp{hablara/hablase}).
\item Négatif : \esp{Para que no} + subj. (ne pas oublier \esp{no} avant le verbe).
\end{enumerate}
\end{methode}

\begin{methode}[Gérer le futur dans la subordonnée temporelle]
\begin{enumerate}
\item Repérer le connecteur temporel : \esp{cuando, en cuanto, hasta que, tan pronto como, después (de) que, antes (de) que}.
\item Regarder le verbe de la principale : s'il exprime \textbf{Futur, Ordre, Infinitif, Probabilité} $\rightarrow$ \textbf{Subjonctif Présent} dans la subordonnée.
\item \esp{Después de que} / \esp{Antes de que} : \textbf{Toujours} subjonctif (même au passé : \esp{Después de que llegó} $\times$ $\rightarrow$ \esp{Después de que llegara}).
\item \esp{Cuando} + Indicatif = action habituelle ou passée accomplie (\esp{Cuando llego, como} / \esp{Cuando llegué, comí}).
\end{enumerate}
\end{methode}

\begin{methode}[Reconnaître et restituer l'emphase]
\begin{enumerate}
\item \textbf{Cleft (\textit{Es... que})} : Isoler l'élément focalisé. \esp{Juan compró el coche} $\rightarrow$ \esp{Fue Juan quien compró el coche} (sujet) / \esp{Fue el coche lo que compró Juan} (objet).
\item \textbf{Fronting} : Déplacer le complément en début de phrase. Ajouter pronom redondant si nécessaire (\esp{El coche lo compró Juan}).
\item \textbf{\textit{Lo que}} : Pour focaliser une action ou une idée entière. \esp{Quiero dormir} $\rightarrow$ \esp{Lo que quiero es dormir}.
\item \textbf{Traduction FR $\rightarrow$ ES} : « C'est... qui... » $\rightarrow$ \esp{Es... quien/que...} ; « Ce que..., c'est... » $\rightarrow$ \esp{Lo que..., es...}.
\end{enumerate}
\end{methode}

\begin{exercice}[Entraînement à la traduction : Thème et Version]
\textbf{A. Thème (Français $\rightarrow$ Espagnol)} \\
1. Puisque tu es ici, nous pouvons commencer la réunion. (formel) \\
2. Il a raté l'examen à cause de son manque d'étude. \\
3. Nous travaillons dur pour que nos enfants aient un avenir meilleur. \\
4. Dès que j'aurai le diplôme, je chercherai du travail à Douala. \\
5. C'est le directeur qui a signé le contrat hier. \\
\\
\textbf{B. Version (Espagnol $\rightarrow$ Français)} \\
1. \esp{Dado que la economía crece, hay más empleo.} \\
2. \esp{El proyecto se retrasó debido a las lluvias torrenciales.} \\
3. \esp{Te explico esto para que lo entiendas bien.} \\
4. \esp{Cuando termines el informe, tráigamelo inmediatamente.} \\
5. \esp{Lo que necesitamos es más inversión en agricultura.}
\end{exercice}

\begin{corrige}[Exercice Leçon E21]
\textbf{A. Thème} \\
1. \esp{Puesto que / Dado que estás aquí, podemos empezar la reunión.} (\textit{Registre formel : puesto que/dado que + indicatif}). \\
2. \esp{Suspendió el examen a causa de / debido a su falta de estudio.} (\textit{Nom après la locution ; debido à invariabilité car pas d'accord apparent, ou accord si "la causa fue debida..."}). \\
3. \esp{Trabajamos duro para que nuestros hijos tengan un futuro mejor.} (\textit{Sujets différents : nosotros / nuestros hijos $\rightarrow$ para que + subj. présent \esp{tengan}}). \\
4. \esp{En cuanto / Tan pronto como tenga el título, buscaré trabajo en Douala.} (\textit{Futur dans la principale \esp{buscaré} $\rightarrow$ subj. présent \esp{tenga} après \esp{en cuanto/tan pronto como}}). \\
5. \esp{Fue el director quien firmó el contrato ayer.} (\textit{Cleft : \esp{Fue} + focalisé \esp{el director} + \esp{quien} (personne) + verbe}). \\
\\
\textbf{B. Version} \\
1. \textit{Vu que / Puisque l'économie croît, il y a plus d'emploi.} (\esp{Dado que} + indicatif, cause évidente, registre soutenu). \\
2. \textit{Le projet a été retardé en raison des pluies torrentielles.} (\esp{Debido a} + nom, accord implicite avec \esp{el retraso} $\rightarrow$ \esp{debido}, ou \esp{a causa de} invariable). \\
3. \textit{Je t'explique cela pour que tu le comprennes bien.} (\esp{Para que} + subj. présent \esp{entiendas}, sujets différents \esp{yo/tú}). \\
4. \textit{Quand tu auras fini le rapport, apporte-le-moi immédiatement.} (\esp{Cuando} + subj. présent \esp{termines} car ordre/imminence dans la principale \esp{tráigamelo}). \\
5. \textit{Ce dont nous avons besoin, c'est plus d'investissement dans l'agriculture.} (\esp{Lo que} = « Ce que / Ce dont » ; \esp{necesitamos} + \esp{es} + nom).
\end{corrige}

\begin{aretenir}[Checklist avant le Bac : Leçon E21]
$\square$ Je distingue \esp{porque} (cause nouvelle), \esp{como} (cause connue en tête), \esp{puesto que/ya que} (cause évidente, formel).\\
$\square$ J'accorde \esp{debido a} (\esp{debida, debidos, debidas}) et j'utilise \esp{a causa de} (invariable) correctement.\\
$\square$ Je choisis \esp{para + inf.} (même sujet) vs \esp{para que + subj.} (sujets différents) pour la finalité.\\
$\square$ Je conjugue le subjonctif présent (\esp{hable, coma, viva}) et imparfait (\esp{hablara/comiera, hablase/comiese}) après \esp{para que} selon le temps de la principale.\\
$\square$ J'applique le \textbf{subjonctif présent} après \esp{cuando, en cuanto, hasta que, tan pronto como} si la principale exprime le futur/ordre.\\
$\square$ Je sais que \esp{después de que} et \esp{antes de que}$\rightarrow$ \textbf{toujours} subjonctif.\\
$\square$ Je construis une phrase cleft : \esp{Es/Fue + [Focalisé] + quien/que/lo que + [Verbe + Reste]}.\\
$\square$ J'utilise \esp{quien} pour les personnes, \esp{que/lo que} pour les choses, \esp{lo que} = « ce qui/ce que ».\\
$\square$ Je pratique le \textbf{Fronting} (OSV) avec pronom redondant (\esp{La carta la escribió María}).\\
$\square$ Je traduis « C'est... qui... » par \esp{Es... quien/que...} et « Ce que je veux, c'est... » par \esp{Lo que quiero es...}.\\
$\square$ Je vérifie les accents diacritiques : \esp{porqué} (nom), \esp{por qué} (interrog.), \esp{por que} (relatif), \esp{porque} (conj.).\\
$\square$ En version, je rends la nuance de registre (ex: \esp{debido a} $\rightarrow$ « en raison de » / « du fait de »).
\end{aretenir}

\findecours

\end{document}
